
% 2026 revision, stage 0 (clean source tree, no mathematical changes).
% Comments marked "% Claude (2026):" are additions by Claude; see ../revision_log.md.
\documentclass[12pt]{article}

\usepackage{amsmath}
\usepackage{amssymb}
\usepackage{amsthm}
\usepackage{amsfonts}
\usepackage[hypertexnames=false,colorlinks=true,linkcolor=blue,citecolor=blue,urlcolor=blue]{hyperref}
\numberwithin{equation}{section}

\newtheorem{theorem}{Theorem}[section]
\newtheorem{corollary}[theorem]{Corollary}
\newtheorem{proposition}[theorem]{Proposition}
\newtheorem{example}[theorem]{Example}
\newtheorem{remark}[theorem]{Remark} % my addition
\newtheorem{CA}[theorem]{Computer Algebra Commands} % my addition
\newtheorem{definition}[theorem]{Definition} % my addition
\newtheorem{lemma}[theorem]{Lemma}
\newtheorem{fact}[theorem]{Fact}
\newtheorem{conjecture}[theorem]{Conjecture}
\newtheorem{algorithm}[theorem]{Algorithm}
\newtheorem{axiom}[theorem]{Axiom}
\newtheorem{ponderable}[theorem]{Exercise}
\newtheorem{exercise}[theorem]{Exercise}
\newtheorem{solution}[theorem]{solution}
\def\unif{\varpi_F}
\def\off{{\mathcal O}_F}
\def\hhh{{\mathcal H}}
\def\ogf{\overline{G(F)}}
\def\oaf{\overline{A(F)}}
\def\oaNf{\overline{A(F)^N}}
\def\oapnf{\overline{A'(F)^n}}
\def\oapf{\overline{A'(F)}}
\def\ogpf{\overline{G'(F)}}
\def\oslf{\overline{SL_2(F)}}
\def\aaa{\mathbb{A}}
\def\fff{\mathbb{F}}
\def\qqq{\mathbb{Q}}
\def\rrr{\mathbb{R}}
\def\ccc{\mathbb{C}}
\def\zzz{\mathbb{Z}}
\def\nnn{\mathbb{N}}
\def\pf{{\bf proof}:\ }
\addtolength{\oddsidemargin}{-.5in}
\addtolength{\evensidemargin}{-.5in}


\begin{document}

\author{David Joyner}
\title{A hitch-hiker's guide to the
metaplectic trace formula for SL(2): basic theory
\thanks{Slightly revised 6-2002}}
\date{1-27-95}
\maketitle

\begin{abstract}
This is an account of the Arthur trace formula for the
n-fold metaplectic cover of $SL(2)$, following closely
Arthur's derivation of the invariant trace formula in this case.
We also discuss the stabilized trace formula of
the covers of $SL(2)$ by defining
stable conjugacy on a cover of $SL(2)$ with conjugacy
in a corresponding cover of $GL(2)$,
then sketch a $\overline{GL(2,\aaa _S)}$-invariant
trace formula as in \cite{LL}.
\end{abstract}

\tableofcontents

\section{Introduction}\label{sec:intro}

We will write down (with some proofs) the invariant trace
formula following \cite{A2}--\cite{A15} in the case of the
n-fold metaplectic cover of $G=SL(2)$.
This is essentially an explicit formula for the trace of
the right regular representation of
$\overline{G(\aaa )}$ on the space of square-integrable
genuine cusp forms.
Then we shall sketch a partial stabilization following \cite{LL}.
Our formula resembles that of \cite{LL} except that there
is no term in our formula corresponding to the
contribution to the residual spectrum
\cite[(5.5), p. 754]{LL}. Moreover, we do not stabilize the spectral
side of the trace formula.

We have left out almost entirely the details of the theory of
Eisenstein
series (there are several excellent presentations available -
\cite{L1,L2}, \cite{K2}). We have also left out most of the
local representation theory needed, for which we refer to the
papers \cite{J1,J3}.

Any comments on this paper may be sent to
(wdj@sma.usna.navy.mil or Math Dept, USNA, Annapolis MD 21402).

\section{Background on metaplectic extensions}\label{sec:background}

Notation: Let $k$ denote a number field containing the group of $n^{th}$
roots of unity $\mu _n\subset \ccc ^{\times }$. In particular, if $n>2$
then $k$ is totally imaginary. Let $\Sigma _k$ denote the set of places of $k
$, $\Sigma _k^\infty $ denote the set of archimedian places of $k$ (which we
may identify with the set of inequivalent complex embeddings $k\hookrightarrow 
\ccc $), and let $\Sigma _k^f$ denote the set of non-archimedian places of 
$k$ (which we may identify with the set of non-zero prime ideals of $%
{\mathcal{O}}_k$). Let $\aaa $ $\;$denote the adele ring of $\qqq$ and
let $\aaa _k$ denote the adele ring of $k$ and, for $v\in \Sigma _k$, let $%
k_v$ denote the completion of $k$ with respect to $v$. For $F=k_v$ p-adic,
let ${\mathcal{O}}={\mathcal{O}}_F={\mathcal{O}}_v$ denote the ring of integers of 
$F$, $\pi =\pi _F=\pi _v$ a local uniformizer, $q$ the cardinality of the
residue field, and $\nu =\nu _F=\nu _v:F^{\times }\rightarrow {\mathbb{Z}}$ \ the
normalized valuation, so that $\nu (\pi )=q^{-1}$.

Let $G$ denote the reductive group over $\qqq$ defined by $SL(2)$\ via
restriction of scalars $k/\qqq$, so for example $G(\qqq)=SL(2,k)$.

\subsection{Local aspects}\label{ssec:local}

Let $F=k_v$, where $v\in \Sigma _k$. For $g=\left( 
\begin{array}{cc}
a & b \\ 
c & d
\end{array}
\right) \in SL_2(F)$, let 
\[
x(g)=\left\{ 
\begin{array}{c}
c,\;\;\;c\neq 0, \\ 
d,\;\;\;c=0,
\end{array}
\right. 
\]
and let $\beta =\beta _v:SL_2(F)\times SL_2(F)\rightarrow \mu _n$ be defined
by 
\begin{eqnarray*}
\beta (g_1,g_2) &=&(x(g_1),x(g_2))_n\cdot (-x(g_1)^{-1}x(g_2),x(g_1g_2))_n \\
&=&(\frac{x(g_1g_2)}{x(g_1)},\frac{x(g_1g_2)}{x(g_2)})_n,
\end{eqnarray*}
where $(...,...)_n=(...,...)_{n,v}:F^{\times }\times F^{\times }\rightarrow $
$\mu _n$ denotes the Hilbert symbol \cite{W1}. Note that this is trivial if $F=%
\ccc $.

\begin{lemma}\label{lem:2.1}
(Kubota \cite{K1}, Moore \cite{M}) The map $\beta $ defines a continuous Borel
2-cocycle (giving $\mu _n$ the discrete topology of course). Its cohomology
class in the continuous cohomology group $H^2(SL_2(F),\mu _n)$ corresponds
to an n-fold topological central extension $\overline{SL_2(F)}$ which fits
into the short exact sequence 
\[
1\rightarrow \mu _n\rightarrow \overline{SL_2(F)}\stackrel{p}{\rightarrow }%
SL_2(F)\rightarrow 1.
\]
\end{lemma}

If $H<SL_2(F)$ is a subgroup then we define $\overline{H}=p^{-1}(H)\subset 
\overline{SL_2(F)}$.

Elements of $\overline{SL_2(F)}$ will be denoted by $(g,\varsigma )$, 
where $g\in SL_2(F)$, $\varsigma \in \mu _n$. Multiplication on 
$\overline{SL_2(F)}$ is given by 
\begin{equation}\label{eq:2.1}
(g_1,\varsigma _1)\cdot (g_2,\varsigma _2)=(g_1g_2,\beta (g_1,g_2)\varsigma
_1\varsigma _2).
\end{equation}
so inverses are given by 
\begin{equation}\label{eq:2.2}
(g,\varsigma )^{-1}=(g^{-1},\beta (g,g^{-1})^{-1}\varsigma ^{-1})
\end{equation}
and conjugation is given by 
\begin{equation}\label{eq:2.3}
(g,\varsigma )\cdot (h,\varsigma ^{\prime })\cdot (g,\varsigma
)^{-1}=(ghg^{-1},\beta (g,h)\beta(gh,g^{-1})
\beta (g,g^{-1})^{-1}\varsigma ^{\prime }).
\end{equation}

We temporarily let $B=\{\left( 
\begin{array}{cc}
\ast & * \\ 
0 & *
\end{array}
\right) \}\subset SL_2$ denote the standard Borel, $N$ its unipotent
radical, and $A$ its Levi component. For
$\theta\in F^\times -F^{\times 2}$, let
\[
T=T_\theta=
\{
\left(
\begin{array}{cc}
a &b\theta \\
b & a
\end{array}
\right)
\ |\ a^2+b^2\theta =1,\ a,b\in F\}.
\]


\begin{lemma}\label{lem:2.2}
(a) $\overline{N(F)}=N(F)\times \mu _n$ (direct product),

(b) if $p$ is relatively prime to $n$ then 
$\overline{SL_2({\mathcal{O}})}=SL_2({\mathcal{O}})\times \mu _n$ 
(direct product),

(c) if $K_0=SL_2({\mathcal{O}})$ and 
\[
K_m=\{\left( 
\begin{array}{cc}
a & b \\ 
c & d
\end{array}
\right) \in SL_2({\mathcal{O}})\;|\;a-1,d-1,b,c\in \pi ^m{\mathcal{O}}\},\;m>0,
\]
then $\overline{K_m}=K_m\times \mu _n$ (direct product) for all $m\geq n$,

(d) if $F=k_v=\mathbb{R}$ then 
$\overline{A(k_v)^0}=A(k_v)^0\times \mu _n$,

(e) The center of $\overline{SL_2(F)}$ is 
\[
Z(\overline{SL_2(F)})=
\{\pm 1_2\}\times \mu _n,
\]

(f) For $a_i=\left( 
\begin{array}{cc}
t_i & 0 \\ 
0 & t_i^{-1}
\end{array}
\right) \in A(F)$, $i=1,2$, and $\varsigma _i\in \mu _n$,
we have $(a_1,\varsigma _1)^{-1}\cdot
(a_2,\varsigma _2)\cdot (a_1,\varsigma
_1)=(a_2,(t_1,t_2)_n^2\varsigma _2)$.
(In particular, $\overline{A(F)}$ is abelian if $n=1,2$.)

(g) $ \overline{T}$ is abelian.

(h) If $g\in G(F)$ and
$n=
\left(
\begin{array}{cc}
1 & 1\\
0 & 1
\end{array}
\right)
$
then $(g,\zeta)(n,\zeta ')(g,\zeta)^{-1}
=(gng^{-1},\zeta')$.
\end{lemma}

proof: Except for (g), (h)
these are all immediate conseqences of the well-known properties of
the Hilbert symbol \cite{W1}. Part (g) is equivalent to the statement
$\beta(\gamma_1, \gamma_2)=\beta(\gamma_2, \gamma_1)$,
for all $\gamma_1, \gamma_2\in T$. This 
is a special case of the calculation in \cite[p. 128]{Fl}.
Part (h) is a direct calculation which we omit. 
$\Box $

The covers of the maximal comapct subgroups do not always split
at the archimedian places either.
We recall (see for example \cite{Ge2}) that the $2$-fold cover of 
$SO(2,{\mathbb R})$ does not split, though for $n>2$ the $n$-fold
covers of $SU(2,{\mathbb C})$ do split.

Note that $G(\qqq_p)=\prod_{v|p}SL_2(k_v)$. Define 
$\beta _p:G(\qqq_p)\times G(\qqq_p)\rightarrow \mu _n$ by 
\[
\beta _p(g_1,g_2)=\prod_{v|p}\beta _v(g_{1,v},g_{2,v}).  
\]
This is a 2-cocycle by Lemma~\ref{lem:2.1} above. Define
$\overline{G(\qqq_p)}$ to be
the n-fold covering of $G(\qqq_p)$ associated to this 2-cocycle. We may
identify $\overline{G(\qqq_p)}$ with a subgroup of 
$\prod_{v|p}\overline{SL_2(k_v)}$.

Let $B=\{\left( 
\begin{array}{cc}
\ast & * \\ 
0 & *
\end{array}
\right) \}\subset G$ denote the standard minimal parabolic
defined over $\qqq$, $N$ its unipotent radical, and
$A$ its Levi component. Let
$T\subset G(\qqq_p)$ denote a Cartan subgroup of the form
$T=\prod_{v|p}T_v$, where $T_v$ is a non-split
Cartan of $SL_2(k_v)$.

\begin{lemma}\label{lem:2.3}
(a) $\overline{N(\qqq_p)}=N(\qqq_p)\times \mu _n$ (direct product),

(b) if $p$ is relatively prime to $n$
then $\overline{G({\mathbb{Z}}_p)}=G({\mathbb{Z}}_p)\times \mu _n$ (direct product),

(c) if $K_0=G({\mathbb{Z}}_p)$ and 
\[
K_m=\prod_{v|p}\{\left( 
\begin{array}{cc}
a & b \\ 
c & d
\end{array}
\right) \in SL_2(k_v)\;|\;a-1,d-1,b,c\in \pi _v^m{\mathcal{O}}_v\},\;m>0,
\]

then $\overline{K_m}=K_m\times \mu _n$ (direct product) for all $m\geq n$,

(d) if $n=2,$ $k=\qqq$, $p=\infty $ then 
$\overline{A({\mathbb R})^0}=A({\mathbb R})^0 \times \mu _n$,

(e) The center of $\overline{G(\qqq_p)}$ is 
\[
Z(\overline{G(\qqq_p)})=
\{\pm 1_2\}\times \mu _n,
\]

(f) For $a_i=\left( 
\begin{array}{cc}
t_i & 0 \\ 
0 & t_i^{-1}
\end{array}
\right) \in A(\qqq_p)$, $i=1,2$, and $\varsigma _i\in \mu _n$, we have $%
(a_1,\varsigma _1)^{-1}\cdot (a_2,\varsigma _2)\cdot (a_1,\varsigma
_1)=(a_2,\prod_{v|p}(t_1,t_2)_{n,v}^2\varsigma _2)$.
(In particular, $\overline{A(\qqq_p)}$ is abelian if $n=1,2$.)

(g) $\overline{T}$ is abelian.
\end{lemma}

proof: These are all immediate conseqences of the previous lemma. $\Box $

\subsubsection{$NAK$ decompositions}

By the above lemma, we may identify $N(\qqq_p)$ with a subgroup of $%
\overline{G(\qqq_p)}$ via $n\longmapsto (n,1)$. We may identify $%
K_0=K_{0,p}$ with a subset of $\overline{G(\qqq_p)}$ via $k\longmapsto
(k,1)$. If $p|n$ then this image is a subgroup.

\begin{lemma}\label{lem:2.4}
We have 
\[
\overline{G(\qqq_p)}=N(\qqq_p)\cdot \overline{A(\qqq_p)}\cdot K_p.
\]

Furthermore, if $(g,\varsigma )=(n,1)\cdot (a,\varsigma ^{\prime })\cdot
(k,1)$ then the map 
\[
(g,\varsigma )\longmapsto a\;({\rm mod}\;Z(G(\qqq_p))
\]
is well-defined.
\end{lemma}

\subsubsection{Heights}

For $x\in \prod_{v|p}x_v$, let $|x|_p={\max }_{v | p}\{|x_v|_v\}$.
Let $g=\left( 
\begin{array}{cc}
a & b \\ 
c & d
\end{array}
\right) \in G(\qqq_p)$ and define 
\[
||g||_p=\left\{ 
\begin{array}{c}
\max (|a|_p,|b|_p,|c|_p,|d|_p),\;\;p<\infty , \\ 
\sqrt{|a|_\infty ^2+|b|_\infty ^2+|c|_\infty ^2+|d|_\infty ^2},\;p=\infty .
\end{array}
\right. 
\]
This is called the height or norm of $g$. We extend this to
$\overline{G({\qqq}_p)}$ by $||(g,\zeta)||_p=||g||_p$.

\subsubsection{Characteristic functions of Weyl chambers}

Let $P_0=B$ denote the minimal parabolic defined above Lemma~\ref{lem:2.3}.
This will be our fixed
minimal parabolic in the set-up of \cite{A2,A3}. The set of all standard
parabolics is simply the set $\{B,G\}$. Let $P=MN$ denote a standard
parabolic and let $A_P$ denote the center of $M$. Thus $A_B=A$ and
$A_G=\{\pm 1\}$.

Define $H=H_B:\overline{G(\qqq_p)}\rightarrow \mathbb{R}$
by $H(g)=\log |t|_{\qqq_p}$, where $g=nak$ and $a=\left( 
\begin{array}{cc}
t & 0 \\ 
0 & t^{-1}
\end{array}
\right) \in A(\qqq_p)$, $t\in \qqq_p$.

\begin{lemma}\label{lem:2.5}
(a) $H(aa^{\prime })=H(a)+H(a^{\prime }),
\;\forall a,a^{\prime }\in \overline{A(\aaa )}$.

(b) $H(kxk^{\prime })=H(x),\;\forall k,k^{\prime }\in K$.

(c) $|H(x)|<<1+\log ||x||,\;\forall x\in \overline{G(\aaa )}$.

(d) For all $x\in \overline{G(\aaa )}$,
$\gamma \in B(\qqq)$, $H(\gamma x)=H(x)$.
\end{lemma}

Let $X(A)=Hom_{\qqq}(A,{\mathbb{G}}_m)\cong {\mathbb{Z}}$ \
(send the character $\chi _n$ defined by $\chi _n(a)=t^n$
to $n\in {\mathbb{Z}}$) and identify ${{\mathfrak{a}}}=Hom(X(A),\mathbb{R})$
with $\mathbb{R}$ (send $X\in {\mathfrak{a}}$ to $X(\chi _1)$).
Let $X_\alpha\in {\mathfrak{a}}$ denote the map such that
$X_\alpha (\chi_1)=1$.
Let ${\mathfrak{a}}^{*}=Hom({\mathfrak{a}},\mathbb{R})$ and let $\alpha \in {\mathfrak{a}}^{*}$
be defined by $\alpha (X)=2X(\chi _1)$.
We may identify ${\mathfrak{a}}$ and ${\mathfrak{a}}^*$ by means of the
linear map which sends
$X_\alpha\in {\mathfrak{a}}$ to $\alpha\in {\mathfrak{a}}^*$.
Note $\alpha(X_\alpha )=2$.
We also let $\alpha \in X(A)$ be the
element such that $\alpha (a)=t^2$, where $a=\left( 
\begin{array}{cc}
t & 0 \\ 
0 & t^{-1}
\end{array}
\right) \in A$, so $\alpha $ the simple positive root of $(G,A)$. By
duality, there is an $\varpi \in ({\mathfrak{a}}^{*})^{*}\cong {\mathfrak{a}}$
such that 
$\varpi (\alpha )=1$, so $\varpi $ is the dual root of $\alpha $. We
may identify $\varpi$ with ${1\over 2}X_alpha$.

Let ${\mathfrak{a}}^{+}=\{X\in {\mathfrak{a}}\;|\;\alpha (X)>0\}$
denote the positive
Weyl chamber and let
$({\mathfrak{a}}^{*})^{+}=\{X\in {\mathfrak{a}}^{*}\;|\;\varpi (X)>0\}$
denote the dual positive Weyl chamber. The elements of
${\mathfrak{a}}^{+} $ are called {\bf regular}. Let $\tau _B$ denote
the characteristic
function of ${\mathfrak{a}}^{+}$ and let $\tau _B^{\wedge }$ denote the
characteristic function of $({\mathfrak{a}}^{*})^{+}$.
Define $\tau _G=\tau _G^{\wedge }=1$.

More generally, if $P_1\subset P_2$ denote two standard
parabolics. let $\Delta _1^2\subset \Delta _1$, $\tau _1^2$,
and $\tau _1^{\wedge ,2}=(\tau^{\wedge })_1^2$ be as in \cite{A2}.
If we identify ${\mathfrak{a}}={\mathfrak{a}}_B$ with $R$
as above then we have 
\begin{eqnarray}
\tau _0 &=&\tau _B=\tau _0^G=\tau _B^G=\mathrm{char.\;fcn.\;of}
\;\{u\in {\mathbb{R}}\;|\;u>0\} \label{eq:2.4}\\
\tau _0^{\wedge }
&=&\tau _B^{\wedge }=\tau _0^{\wedge ,G}
=\tau _B^{\wedge,G}=\mathrm{char.\;fcn.\;of}
\;\{u\in {\mathbb{R}}\;|\;u>0\}  \nonumber \\
\tau _1^2 &=&\left\{ 
\begin{array}{c}
\tau _1,\;\;\mathrm{if}\;P_2=G, \\ 
1,\;\;\;\mathrm{if}\;P_2\neq G,
\end{array}
\right.  \nonumber
\end{eqnarray}
Finally, if $\sigma _1^2$ is as in \cite{A2}, then 
\begin{equation}\label{eq:2.5}
\sigma _1^2=\left\{ 
\begin{array}{c}
\tau _1,\;\;\mathrm{if}\;P_1=B,\;P_2=G, \\ 
0,\;\;\;\;\;\mathrm{if}\;P_2=P_1=B, \\ 
1,\;\;\;\;\mathrm{if}\;P_2=P_1=G.
\end{array}
\right.
\end{equation}

\subsubsection{Algebras of functions}\label{sssec:localalg}

If $H$ is the group of $F$-rational points of a reductive algebraic group
over a real or p-adic field $F$, let $C_c^\infty (H)$ denote the usual
algebra of compactly supported smooth functions under convolution and $%
C^\infty (H//K)$ the space of smooth functions which are bi-$K$-ivariant. A
function $f:\overline{G(\qqq_p)}\rightarrow \ccc $ such that 
\begin{equation}\label{eq:2.6}
f(g,\varsigma )=\varsigma \cdot f(g,1)
\end{equation}
will be called {\bf genuine}. Let 
\begin{eqnarray*}
C_c^\infty (\overline{G(\qqq_p)})_{gen}
&=&\{f:\overline{G(\qqq_p)}%
\rightarrow \ccc \
|\;(i)\;f|_{G(\qqq_p)}\in C_c^\infty (G(\qqq_p)),
\\
&&(ii)\;f\;\mathrm{is\;genuine}\}
\end{eqnarray*}
denote the Schwartz-Bruhat algebra of smooth genuine functions.
If $(p,2n)=1$ then let 
\begin{eqnarray*}
{\mathcal{H}}(\overline{G(\qqq_p)},K_p)
&=&\{f:\overline{G(\qqq_p)}\rightarrow \ccc 
\ |\;(i)\;f|_{G(\qqq_p)}\in
C^\infty (G(\qqq_p)//K_p), \\
&&\ (ii)\;f\;\mathrm{is\;genuine}\}
\end{eqnarray*}
denote the Hecke algebra of genuine spherical functions.

\subsubsection{Measures}

Let $dk$ denote the Haar measure on $K_p$ normalized so that $K_p$ has
volume 1. Fix an additive non-trivial character $\psi _v$ of $k_v$, for each$%
\;v\in \Sigma _k$, and let $dn$ denote a Haar measure on $N(\qqq_p)$
which is induced by a self-dual measure on the additive group 
\[
\bigoplus_{v\in \Sigma _k, v|p}k_v
\]
(with respect to the $\psi _v,$ $v|p$).
Let $da$ be a Haar measure on $A(\qqq_p)$. If $T$ is a Cartan subgroup of 
$G$ defined over $\qqq_p$ and $\psi :T\rightarrow A$ is a $\overline{\mathbb{%
Q}_p}$-isomorphism given by $h\longmapsto t^h$, some $h\in G(\qqq_p)$,
then we define $dt$ to be the pull-back of $da$ from $A$ to $T$ as in \cite{R}.
If $T^{\prime }$ is another Cartan subgroup of $G$ defined over $\qqq_p$
which is stably conjugate to $T$ then we assume that its Haar measure
$dt^{\prime }$ is compatible with $dt$ as in \cite{R}. Let $dg$ be a Haar measure
on $G(\qqq_p)$ such that $dg=\delta _B(a)^{1/2}dkdadn$, $g=kan$, where 
\begin{equation}\label{eq:2.7}
\delta _B(a)=|t|^2,\;\;a=\left( 
\begin{array}{cc}
t & 0 \\ 
0 & t^{-1}
\end{array}
\right) .
\end{equation}

We extend these measures to the n-fold cover by imposing the discrete
measure $d\varsigma $ on $\mu _n$ which gives $\mu _n$ total volume 1. For
example, the measure on $\overline{G(\qqq_p)}$ will be denoted $d%
\overline{g}=dg\cdot d\varsigma $.

\subsection{Global aspects}\label{ssec:global}

Define a map $\beta =\beta _{\aaa }:G(\aaa )\times G(\aaa )\rightarrow
\mu _n$ by 
\[
\beta (g_1,g_2)=\prod_{p\leq \infty }\beta _p(g_{1,v},g_{2,v}).
\]
This product converges by Lemma~\ref{lem:2.3}(b) and the definition of the restricted
direct product. As a product of local 2-cocycles it inherits the cocycle
property and therefore defines an element of $H^2(G(\aaa ),\mu _n)$
associated to an n-fold central topological covering group $\overline{G(\mathbb{%
A})}$ fitting into the following exact sequence 
\begin{equation}\label{eq:2.8}
1\rightarrow \mu _n\rightarrow \overline{G(\aaa )}\stackrel{p}{\rightarrow 
}G(\aaa )\rightarrow 1.
\end{equation}
The rules of multiplication, inverses, and conjugation are still as in
equations \eqref{eq:2.1}--\eqref{eq:2.3} above.

Note that we may identify $\overline{G(\aaa )}$ with a subgroup of the
restricted direct product
$\prod_{v\in \Sigma _k}\overline{SL_2(k_v)}$. Again,
if $H<G(\aaa )$ is a subgroup then we define
$\overline{H}=p^{-1}(H)<\overline{G(\aaa )}$.

\begin{lemma}\label{lem:2.6}
(a)\ $\overline{G(\qqq)}\cong G(\qqq)\times \mu _n$
(direct product), where the isomorphism in $\cong$ will be
described below,

(b)\ $\overline{N(\qqq)}=N(\qqq)\times \mu _n$ (direct product),

(c) if $K(n)^f=\prod_{p|n}K_{n,p}$, $K(n)_f=\prod_{(p,n)=1}K_{0,p}$, and
$K(n)=K(n)^f\times K(n)_f$ then 
\[
\overline{K(n)}=K(n)\times \mu _n,
\ \ \ {\rm (direct product)}.
\]

(d) The center of $\overline{G(\aaa )}$ is
$Z(\overline{G(\aaa )})=Z(G({\aaa }))\times \mu _n$,

(e) For $a_i=\left( 
\begin{array}{cc}
t_i & 0 \\ 
0 & t_i^{-1}
\end{array}
\right) \in A(\aaa )$, $i=1,2$, and $\varsigma _i\in \mu _n$, we have $%
(a_1,\varsigma _1)^{-1}\cdot (a_2,\varsigma _2)\cdot (a_1,\varsigma
_1)=(a_2,(t_1,t_2)_n^2\varsigma _2)$. Here $(...,...)_n$ denotes the global
Hilbert symbol.
\end{lemma}

proof: This is an immediate consequnce of the above lemmas or well-known
properties of the global Hilbert symbol \cite{W1}. $\Box $

From \cite{K2}, we see that
the splitting of the above exact sequence indicated in part (a) of the above
lemma is not given by $\gamma \longmapsto (\gamma ,1)$ but by $\gamma
\longmapsto (\gamma ,s_{\aaa }(\gamma )^{-1})$, where 
\begin{equation}\label{eq:2.9}
\begin{array}{c}
s_{\aaa }(\gamma )=
\prod_{v\in \Sigma _k}
s_v(\gamma),\\
s_v(\gamma)=\left\{ 
\begin{array}{c}
(c,d)_{v,n}^{-1},
\;\;\;cd\neq 0{\rm \ and\ }
n\;\mathrm{does\;not\;divide}\;v(c)
{\rm \ and\ }v<\infty, \\ 
1,\;\;\;\;\;\;\;\;\;\;\;\;\;\;\;\;
\mathrm{\;otherwise}.\;\;\;\;\;\;\;\;
\end{array}
\right.
\end{array}
\end{equation}
Here $\gamma =\left[ 
\begin{array}{cc}
a & b \\ 
c & d
\end{array}
\right] \in SL_2(k)$.

{\em From now on, we identify $G(\qqq)$ with its image
in $\overline{G(\aaa )}$ under this splitting (and the diagonal
embedding).} We identify $N(\aaa )$ with the
subgroup
$N(\aaa )\times 1\subset \overline{G(\aaa )}$.

Let $K_0=\Pi_{p\leq infty}K_{0,p}$
denote the ``usual'' maximal compact subgroup of $G(\aaa )$.
Here, if $p<\infty $ then $K_{0,p}=G({\mathbb Z}_p)$ and
otherwise
\[
K_{0,\infty}=\prod_{ v|\infty, v\in\Sigma_k}K_v, 
\]
where $K_v=SO(2,{\mathbb{R}})$ if $v$ is real and
$K_v=SU(2,\ccc )$ if $v$ is complex.

The following result is essentially due to Kubota \cite[Proposition 5]{K2} (see also \cite[Lemma 3.2]{Ge2})

\begin{lemma}\label{lem:2.7}
For each compact open subgroup $K=\Pi _{p\leq \infty }K_p\subset K_0\subset
G(\aaa ),$ there are discrete subgroups of cofinite volume $\Gamma
_i\subset \overline{G({\mathbb{R}})}$ such that 
\[
G({\mathbb Q})\backslash \overline{G(\aaa )}/\overline{K}
\cong \coprod_i
\Gamma _i\backslash \overline{G({\mathbb{R}})}/\overline{K_\infty }.
\]
\end{lemma}

\subsubsection{$NAK$ decompositions}

By the above lemma, we may identify $N(\aaa )$ with a subgroup of $%
\overline{G(\aaa )}$ via $n\longmapsto (n,1)$. We may identify $%
K_0=\prod_{p\leq \infty }K_{0,p}$ with a {\em subset} of
$\overline{G(\aaa )}$ via $k\longmapsto (k,1)$.

\begin{lemma}\label{lem:2.8}
We have 
\[
\overline{G(\aaa )}
=N(\aaa )\cdot \overline{A(\aaa )}\cdot K_0.
\]

Furthermore, if $(g,\varsigma )=(n,1)\cdot (a,\varsigma ^{\prime })\cdot
(k,1)$ then the map
$(g,\varsigma )\longmapsto a\;({\rm mod}\;Z(G(\aaa )))$
is well-defined.
\end{lemma}

\subsubsection{Heights}

For $g\in \prod_{p\leq \infty }g_p\in G(\aaa )$,
$g_p=\left( 
\begin{array}{cc}
a_p & b_p \\ 
c_p & d_p
\end{array}
\right) \in G(\qqq_p)$, define 
\begin{equation}\label{eq:2.10}
||g||=\prod_{p\leq \infty }||g_p||_p
\end{equation}
This is called the height or norm of $g$. Extend this to
$\overline{G(\aaa )}$ by $||(g,\zeta)||=||g||$.

\begin{lemma}\label{lem:2.9}
(a) $||x||\geq 1,\;\forall x\in \overline{G(\aaa )}$.

(b) $||kxk^{\prime }||=||x||,\;\forall k,k^{\prime }\in K$.

(c) $||xy||\leq ||x||\cdot ||y||,\;\forall x,y\in \overline{G(\aaa )}$.

(d) There is a constant $N>0$ such that for all
$x\in \overline{G(\aaa )}$,
$||x^{-1}||\leq ||x||^N$.
\end{lemma}

\subsubsection{Algebras of functions}

A function $f:\overline{G(\aaa )}\rightarrow \ccc $
such that \eqref{eq:2.6}
holds will be called {\bf genuine}. Let 
${\mathcal{H}}(\overline{G(\aaa )})$ denote the algebra of 
functions generated by the elements of the
restricted tensor product 
$\bigotimes_{p\leq \infty }
C_c^\infty (\overline{G(\qqq_p)})_{gen}$, 
where almost all $f_p$ in the product 
$f=\otimes _{p\leq \infty}f_p\in 
\bigotimes_{p\leq \infty }C_c^\infty (\overline{G(\qqq_p)})_{gen}$ are
spherical in the sense of \S\ref{sssec:localalg}.

\subsubsection{Measures}

Let $dk$ denote the Haar measure on $K=\prod_{p\leq \infty }K_p$
defined by $dk=\prod_{p\leq \infty }dk_p$.
Let $dg$ be the Tamagawa measure on
$G(\aaa )$ \cite{W2}. For any Cartan subgroup
$T$ of $G$ defined over $\qqq$, let $dt$
denote the global Tamagawa measure on $T(\aaa )$. If $T=A$ then this will
be denoted by $da$. Let $dn$ denote the product measure on $N(\aaa )$.

We extend these measures to the n-fold cover by imposing the discrete
measure $d\varsigma $ on $\mu _n$ which gives $\mu _n$ total volume 1. For
example, the measure on $\overline{G(\aaa )}$ will be denoted
$d\overline{g}=dg\cdot d\varsigma $.

\section{The rank 0 case}\label{sec:rank0}

It might be appropriate here to discuss the rank 0 case. The metaplectic
covers of $A(\mathbb{A})$ are not in general abelian but nonetheless, the
trace formula takes an especially simple form in this case. In this section,
we shall prove the trace formula for metaplectic covers of $A({\mathbb{A}})$ and
apply it to prove a ``matching'' between automorphic forms on $A({\mathbb{A}})$
and on metaplectic covers of $A({\mathbb{A}})$. At the same time, we construct
invariant distributions $I_{\mathcal{O}}^A$ and $I_\chi ^A$ which will be
needed in an inductive definition of certain invariant distributions
on $\overline{G({\mathbb{A}})}$ in Section~\ref{sec:invariant}. 

\subsection{Metaplectic lifting over the split torus}\label{ssec:rank0lift}

Suppose that $k/{\mathbb{Q}}$ is a number field containing the $n-th$ roots of
unity. (If $n>2$ then this implies $k$ has no real places, i.e, $k_v=\mathbb{C}$
for all $v\in \Sigma _k^\infty $.)
Let $A$ and $A({\mathbb{A}})^1$ be as in the previous section,
so $A({\mathbb{Q}})=k^\times$.
\emph{We shall often identify }
$A({\mathbb{A}})$
\emph{\ with }
${\mathbb{A}}_k^{\times }$
\emph{\ in this section}. Let
$(..,...)_n:A({\mathbb{A}})\times A({\mathbb{A}})
\cong {\mathbb{A}}_k^{\times }\times {\mathbb{A}}_k^{\times}\rightarrow \mu _n$
denote the global Hilbert symbol defined by
\cite[section XIII.5, proposition 7]{W1} as a product of local
Hilbert symbols $(...,...)_{n,v}$, where
$\mu _n$ is the group of $n$-th roots of unity. Let $\mu _n$ act on 
$G({\mathbb{A}})={\mathbb{A}}_k^{\times }$ trivially.
A 2-cocycle on $A({\mathbb{A}})$ with
values in $\mu _n$ is a continuous mapping
$\beta :A({\mathbb{A}})\times A({\mathbb{A}})\rightarrow \mu _n$ satisfying 
\begin{equation}\label{eq:3.1}
\beta (g_2,g_3)\beta (g_1,g_2g_3)=\beta (g_1,g_2)\beta (g_1g_2,g_3),\;
\end{equation}
$g_1,g_2,g_3\in {\mathbb{A}}_k^{\times }$. We also denote by $\beta $ the
restriction
$\beta :A({\mathbb{A}})^1\times A({\mathbb{A}})^1\rightarrow \mu _n$. This
is a product of local 2-cocycles, denoted $\beta _v$. Let
$\overline{A(\mathbb{A})}^1$ denote the $n$-fold
metaplectic cover of $A({\mathbb{A}})^1$
defined by the 2-cocycle $\beta $ given by 
\begin{equation}\label{eq:3.2}
\beta (g_1,g_2)=(g_1^{-1},g_2^{-1})_n\cdot (-g_1g_2^{-1},g_1^{-1}g_2^{-1})_n.
\end{equation}
The fact that this satisfies the condition above may be verified by
multiplying each term out and using the properties of the Hilbert symbol
listed in \cite[Proposition 60, p. 176]{Sha} (this is long but
straightforward and in any case well-known).
By definition of group extension, multiplication
on $\overline{A({\mathbb{A}})^1}=A({\mathbb{A}})^1\times \mu _n$
(as \emph{sets}) is given by 
\begin{equation}\label{eq:3.3}
(g,\varsigma )\cdot (g^{\prime },\varsigma ^{\prime })=(gg^{\prime },\beta
(g_1,g_2)\varsigma \varsigma ^{\prime }),
\;\;\;\;g,g^{\prime }\in A({\mathbb{A}}),
\;\;\varsigma ,\varsigma ^{\prime }\in \mu _n.
\end{equation}
This is a central extension of $A({\mathbb{A}})^1$ (i.e.,
$\mu _n$ belongs to the
center of the metaplectic cover) which fits into the following exact
sequence 
\[
1\rightarrow \mu _n\rightarrow \overline{A({\mathbb{A}})^1}
\rightarrow A({\mathbb{A}})^1\rightarrow 1,
\]
where $\mu _n\rightarrow \overline{A({\mathbb{A}})^1}$ is the
inclusion of $\mu _n$ into the second component of
$\overline{A({\mathbb{A}})}^1$, and the map
$\overline{A({\mathbb{A}})^1}\rightarrow A({\mathbb{A}})^1$ is the
projection map.

Note that
$\overline{A({\mathbb{A}})}^1=\overline{A({\mathbb{A}})^1}$.

Using the local 2-cocycles $\beta _v$, we may in turn define local n-fold
covers $\overline{k_v^{\times }}$ of $k_v^{\times },\;v\in \Sigma _k$. The
rule for multiplication on the local covers
\[
\overline{A({\mathbb{Q}}_p)}=
\overline{\prod_{v|p}k_v^{\times }}
\subset \prod_{v|p}
\overline{k_v^{\times }}
\]
of $\prod_{v|p}k_v^{\times }$ is the same, except one must
replace $\beta $ by $\prod_{v|p}\beta _v$. The topology on
$\overline{k_v^{\times }}$
is defined by stipulating 
\[
\{U_v\times 1\;|\;\beta _v(U_v,U_v)=1\;\mathrm{and}\;1\in U_v\subset
k_v^{\times }\;\mathrm{is\;open}\}
\]
to be a basis of open neighborhoods of the identity.
If $(p,n)=1$ and if $v|p$ then the n-fold covers
$\overline{A(\mathbb{Z}_p)}$ and
$\overline{{\mathcal{O}}_v^{\times }}$ split by
the properties of the Hilbert symbol listed in \cite[section XIII.5]{W1}.

The global cover $\overline{A({\mathbb{A}})}$ is not the restricted direct
product of the local covers $\overline{k_v^{\times }}$. However, there is an
embedding 
\[
\overline{A({\mathbb{A}})}\hookrightarrow
\prod_{v\in \Sigma _k}^{\prime }\overline{k_v^{\times }}.
\]

We call a representation $\pi $ of $\overline{A({\mathbb{A}})^1}$
\textbf{genuine} if it satisfies

\begin{equation}\label{eq:3.4}
\pi (g,\varsigma )=\varsigma^{-1} \cdot \pi (g,1).
\end{equation}
A {\em function} satisfying this property will sometimes be
called {\bf anti-genuine}.
Similarly, we define a representation $\pi $ of
$\overline{A({\mathbb{A}})}=\overline{{\mathbb{A}}_k^{\times }}$ or
of $\overline{k_v^{\times }}$ to be 
\textbf{genuine} if it satisfies \eqref{eq:3.4}.

If $H<\overline{A({\mathbb{A}})}$ a locally compact totally disconnected group
of type 1, let $\Pi (H)$ denote
the set of equivalence classes of admissible genuine
representations of $H$ and
let $\Pi_u(H)$ denote the subset of unitary ones.
Let $\Pi_f (\overline{A({\mathbb{A}})})$
denote the subset (equivalance classes of) representations in
$\Pi(\overline{A({\mathbb{A}})})$ which can be written as the
restriction to $\overline{A({\mathbb{A}})}$
of a {\em factorizable} representation $\overline{\pi }$
of the restricted direct product 
\[
\Pi (\prod_v\overline{k_v^{\times }})=\prod_v
\Pi (\overline{k_v^{\times }}).
\]
Let $\Pi_f (\overline{A({\mathbb{A}})}^1)$ denote the
set of equivalence classes of
representations of $\overline{A({\mathbb{A}})^1}$ which can be written as the
restriction to $\overline{A({\mathbb{A}})^1}$ of an element of
$\Pi_f (\overline{A({\mathbb{A}})})$.
Let $\Pi_{u,f} (\overline{A({\mathbb{A}})})$,
$\Pi_{u,f} (\overline{A({\mathbb{A}})^1})$
denote the corresponding subset of factorizable unitary representations.

\begin{lemma}\label{lem:3.1} (\cite[\S 3]{Ge2})
\[
\begin{array}{c}    
\Pi_f (\overline{A({\mathbb{A}})})=\Pi (\overline{A({\mathbb{A}})}),\\
\Pi_{u,f} (\overline{A({\mathbb{A}})})=\Pi_u (\overline{A({\mathbb{A}})}),\\
\Pi_f (\overline{A({\mathbb{A}})^1})=\Pi (\overline{A({\mathbb{A}})^1}),\\
\Pi_{u,f} (\overline{A({\mathbb{A}})^1})=\Pi_u (\overline{A({\mathbb{A}})^1})
\end{array}
\]
\end{lemma}

See also \cite{Ge1}, \cite{F} for the non-metaplectic version.

We shall sketch a proof of the following

\begin{theorem}\label{thm:3.2}
There is a natural 1-1 correspondence between the subset of all genuine
irreducible representations
$\overline{\chi }\in \Pi (\overline{A({\mathbb{A}})}^1)$
of $\overline{A({\mathbb{A}})^1}$ and characters of the form
$\chi ^{n/2}$ (n even), $\chi ^n$ (n odd),
for $\chi \in $\ $\Pi (A({\mathbb{A}})^1)$.
\end{theorem}

First we prove a trace formula for $\overline{A({\mathbb{A}})}^1$.

\subsection{Irreducible representations of $\overline{A({\mathbb{A}})}^1$}

For any group $H$, let
Let 
\begin{equation}\label{eq:3.5}
H^m
=\{x\in H\;|\;x=y^m,\;\;\mathrm{some}
\;\;y\in H\}.
\end{equation}
Let 
\begin{equation}\label{eq:3.6}
N=
\left\{
\begin{array}{c}
n,\ \ n \mathrm{\ odd},\\
n/2,\ \ n \mathrm{\ even}.
\end{array}
\right.
\end{equation}

We have, by the properties of the Hilbert
symbol,
$\overline{A({{\mathbb{A}}})^{1,n}}=A({{\mathbb{A}}})^{1,n}\times \mu _n$,
that is to say, the metaplectic cover
splits over $A({{\mathbb{A}}})^{1,n}$.
By Lemma~\ref{lem:2.6}(e) above, the cover
$\overline{A({{\mathbb{A}}})^{1,N}}$ is abelian (though it
does not split).

For $\chi \in \Pi (\overline{A({\mathbb{A}})^1})$,
$h\in \overline{A({\mathbb{A}})}$, define
$\chi ^h(g)=\chi (h^{-1}gh)$, for $g\in 
\overline{A({\mathbb{A}})}^1$. We call such a
$\chi $ \textbf{self-conjugate} if
there is an $h$ as above such that $\chi =\chi ^h$
in $\Pi (\overline{A({\mathbb{A}})^1})$.
If $h=(x,\varsigma ),\;\varsigma \in \mu _n,$ then 
\[
\chi ^h(g)=\chi (h^{-1}gh)
=\chi (x^{\prime },(x,x^{\prime })_n^2\varsigma^{\prime }),
\;\;\;\;\;\;\;\;\;\;g=(x^{\prime },
\varsigma ^{\prime })\in \overline{A({\mathbb{A}})^1},
\]

If $p<\infty$, let $C_p\subset A({\mathbb{Q}}_p)$ denote a subgroup
of finite index such that
$\overline{C_p}\subset \overline{A({\mathbb{Q}}_p)}$ is
an abelian subgroup. (In particular, we may take
$C_p=A({\mathbb{R}})^N$.) If $p=\infty$ let $C_p=A({\mathbb{R}})$. Let
\begin{equation}\label{eq:3.7}
C_{{\mathbb{A}}}=\prod_{p\leq \infty }C_p
\end{equation}
denote the restricted direct product.

We have the following result (due to Clifford, Mackey,
...)

\begin{lemma}\label{lem:3.3}
If $\chi \in \Pi (\overline{C_{{\mathbb{A}}}})$ is not self-conjugate then 
\[
Ind_{\overline{C_{{\mathbb{A}}}}}^{\overline{A({\mathbb{A}})}}(\chi )
\cong Ind_{\overline{C_{{\mathbb{A}}}}}^{\overline{A({\mathbb{A}})}}(\chi ^h),
\;\;\;h\in  \overline{A({\mathbb{A}})},
\]
is irreducible and belongs to $\Pi (\overline{A({\mathbb{A}})})$.
\end{lemma}

\begin{remark}\label{rem:3.4} By analogy with Mackey's irreduciblity
criterion in \cite[\S 7.4]{Se}, the converse to the avove result also
holds. The detailed proof of the local result can be found in
\cite{J3}.
\end{remark}

\textsf{proof}: We must first reduce this down to a statement about local
representations and then prove the local statement.

Pick a $p\leq \infty $. By Clifford-Mackey theory (\cite{J3}, \cite[\S 7.4]{Se}),
if $\chi \in \Pi (\overline{A({\mathbb{Q}}_p)})$ is not self-conjugate then 
\begin{equation}\label{eq:3.8}
Ind_{\overline{C_p}}^{\overline{A({\mathbb{Q}}_p)}}(\chi )
\cong Ind_{\overline{C_p}}^{\overline{A({\mathbb{Q}}_p)}}(\chi ^h),
\;\;\;h\in \overline{A({\mathbb{Q}}_p)},
\end{equation}
is irreducible and belongs to $\Pi (\overline{A({\mathbb{Q}}_p)})$. We may take
the restricted direct product of these local representations to obtain an
irreducible (by Schur's lemma) factorizable representation on
the restricted direct product of of these local covers. $\Box $

In our case, it is easy to see that for any
$\chi \in \Pi (\overline{C_p})$,
the restriction of
$Ind_{\overline{C_p}}^{\overline{A({\mathbb{Q}}_p)}}(\chi )$
to the abelian subgroup
$\overline{C_p}$ must decompose into a direct sum of
characters: 
\begin{equation}\label{eq:3.9}
Ind_{\overline{C_p}}^{\overline{A({\mathbb{Q}}_p)}}(\chi )|_{\overline{A_p}}
\cong \bigoplus_{i=1}^dm_i\chi _i,
\;\;\;\chi _i\in \Pi (\overline{C_p}),
\end{equation}
where $m_i,d$ is are some integers and
one of these $\chi _i$ must be $\chi $ itself. By the isomorphism
(for any $h\in \overline{A({\mathbb{Q}}_p)}$)
\begin{equation}\label{eq:3.10}
Ind_{\overline{C_p}}^{\overline{A({\mathbb{Q}}_p)}}(\chi )
\cong Ind_{\overline{C_p}}^{\overline{A({\mathbb{Q}}_p)}}(\chi ^h),
\;\;\;h\in \overline{A({\mathbb{Q}}_p)},
\end{equation}
one of these $\chi _i$ must be equal to $\chi ^h$.
In fact, we have

\begin{lemma}\label{lem:3.5} (\cite[Lemma 2.1]{GeKn})
If $\chi \in \Pi (\overline{C_p})$ then all the $m_i's$ in \eqref{eq:3.9} are
equal (to $m$, say) and
\[
Ind_{\overline{C_p}}^{\overline{A({\mathbb{Q}}_p)}}(\chi )|_{\overline{C_p}}
\cong \bigoplus_{h\in \overline{A({\mathbb{Q}}_p)}
/\overline{A({\mathbb{Q}}_p)}_{\chi}}
m\chi ^h,
\]
where
\[
\overline{A({\mathbb{Q}}_p)}_{\chi}
=\{a\in\overline{A({\mathbb{Q}}_p)} \ |\ \chi^a=\chi\}.
\]
\end{lemma}

\begin{remark}\label{rem:3.6} By analogy with \cite[ch 7]{Se}, one might expect there to
be a more general form of this result, closely related to
Frobenius reciprocity.
\end{remark}
\begin{remark}\label{rem:3.7}
An aside: Even using the explicit local recipocity law of
(Sen, Wiles, ...) it does not appear to be
easy to explicitly determine $C_p$
such that $A({\mathbb{Q}}_p)/C_p$ is minimal.
\end{remark}


\subsection{The trace formula for \emph{$\overline{A({\mathbb{A}})^1}$}}

There is natural projection
$pr:\overline{A({\mathbb{A}})}\;\rightarrow \;A({\mathbb{A}}),
\;(x,\varsigma )\longmapsto x$. Let
${\mathcal{S}}(\overline{A({\mathbb{A}})}^1)$
denote the collection of all genuine functions $f$ on
$\overline{A({\mathbb{A}})^1}$ such
that $f\circ pr\in {\mathcal{S}}(A({\mathbb{A}})^1)$. Let 
\[
H=
L^2(\overline{A({\mathbb{Q}})}\backslash \overline{A({\mathbb{A}})^1)})_{gen},
\]
where \emph{$\overline{A({\mathbb{Q}})}$} is regarded as a subgroup
of $\overline{A({\mathbb{A}})^1}$ via the
diagonal embedding. Let $R$ denote
the right regular representation: 
\[
R:\overline{A({\mathbb{A}})^1}\rightarrow \;End(H),
\]
defined by $(R(y)\varphi )(x)
=\varphi (xy),\;x,y\in \overline{A({\mathbb{A}})}^1,
\;\varphi \in H$.
For $\overline{f}\in {\mathcal{S}}(\overline{A({\mathbb{A}})}^1)$, let 
\[
R(\overline{f})=\int_{\overline{A({\mathbb{A}})}^1}\overline{f}(x)R(x)dx.
\]
Since $R(x):H\rightarrow H$ is a unitary operator (in particular bounded),
this integral converges. This
operator $R(\overline{f}):H\rightarrow H$ may be written 
\begin{equation}\label{eq:3.11}
\begin{array}{c}
(R(\overline{f})\varphi )(x)
=\int_{\overline{A({\mathbb{A}})}^1}\overline{f}(x)(R(y)\varphi )(x)dy,
\;\;\;\varphi \in H \\ 
=\int_{\overline{A({\mathbb{A}})}^1}\overline{f}(x^{-1}y)\varphi (y)dy \\ 
=\int_{A({\mathbb{Q}})\backslash \overline{A({\mathbb{A}})}^1}
\sum_{\delta \in A({\mathbb{Q}})}
\overline{f}(x^{-1}\delta y)\varphi (\delta y)dy \\ 
=\int_{A({\mathbb{Q}})\backslash \overline{A({\mathbb{A}})^1}}
\varphi (y)
\sum_{\delta \in A({\mathbb{Q}})}\overline{f}(x^{-1}\delta y)dy.
\end{array}
\end{equation}
Let 
\[
K(x,y)=K_{\overline{f}}(x,y)
=\sum_{\delta \in A({\mathbb{Q}})}
\overline{f}(x^{-1}\delta y).
\]

For each $\overline{\chi }\in \Pi (\overline{A({\mathbb{A}})}^1)$, let 
\begin{equation}\label{eq:3.12}
\begin{array}{c}
\overline{\chi }(\overline{f})
=\int_{\overline{A({\mathbb{A}})}^1}
\overline{f}(y)\overline{\chi }(y)dy, \\ 
tr\;\overline{\chi }(\overline{f})
=\int_{\overline{A({\mathbb{A}})^1}}\overline{f}(y)
tr\;(\overline{\chi }(y))dy
=\overline{f}^{\wedge }(\overline{\chi }).
\end{array}
\end{equation}
We call the latter integral the \textbf{Fourier transform}
of $\overline{f}$ at $\overline{\chi }$, and the
former integral the \textbf{operator-valued}
\textbf{Fourier transform} of $\overline{f}$ at
$\overline{\chi }$. Note that these are
$0$ if $\overline{f}$ is not genuine.

Finally, for each $h\in $ $\overline{A({\mathbb{A}})^1}$, let
$\overline{A({\mathbb{A}})}_h^1$ denote the \textbf{centralizer} of $h$
in $\overline{A({\mathbb{A}})^1}$: 
\begin{equation}\label{eq:3.13}
\overline{A({\mathbb{A}})}_h^1
=\{A\in \overline{A({\mathbb{A}})}^1\;|\;ghg^{-1}=h\}.
\end{equation}

\begin{theorem}\label{thm:3.8}
(a) $tr\;R(\overline{f})
=\int_{A({\mathbb{Q}})\backslash \overline{A({\mathbb{A}})}^1}
K_{\overline{f}}(x,x)dx.$

(b) For each $\overline{\chi }\in \Pi (\overline{A({\mathbb{A}})}^1)$, there is
an integer $m(\overline{\chi })\geq 0$\ such that 
\[
R(\overline{f})
=\bigoplus_{\overline{\chi }\in \Pi_u (\overline{A({\mathbb{A}})}^1)}
m(\overline{\chi })\overline{\chi }(\overline{f}).
\]
\end{theorem}

Thus we have 
\begin{equation}\label{eq:3.14}
\begin{array}{c}
tr\;R(\overline{f})
=\sum_{\overline{\chi }\in \Pi (\overline{A({\mathbb{A}})}^1)}
m(\overline{\chi })tr(\overline{\chi }(\overline{f})) \\ 
=\int_{A({\mathbb{Q}})\backslash \overline{A({\mathbb{A}})}^1}
\sum_{\delta\in A({\mathbb{Q}})}
\overline{f}(x^{-1}\delta x)dx \\ 
=\int_{A({\mathbb{Q}})\backslash \overline{A({\mathbb{A}})}^1}
\sum_{\{\delta \}}
\sum_{\gamma \in A({\mathbb{Q}})_\delta\backslash A({\mathbb{Q}})}
\overline{f}(x^{-1}\gamma ^{-1}\delta \gamma x)dx \\ 
=\sum_{\{\delta \}}
\int_{A({\mathbb{Q}})\backslash \overline{A({\mathbb{A}})}^1}
\sum_{\gamma \in A({\mathbb{Q}})_\delta \backslash A({\mathbb{Q}})}
\overline{f}(x^{-1}\gamma ^{-1}\delta \gamma x)dx \\ 
=\sum_{\{\delta \}}
\int_{A({\mathbb{Q}})_\delta \backslash \overline{A({\mathbb{A}})}^1}
\overline{f}(x^{-1}\delta x)dx \\ 
=\sum_{\{\delta \}}
vol(A({\mathbb{Q}})_\delta \backslash \overline{A({\mathbb{A}})}_\delta ^1)
\sum_{x\in \overline{A({\mathbb{A}})}_\delta ^1
\backslash \overline{A({\mathbb{A}})}^1}
\overline{f}(x^{-1}\delta x),
\end{array}
\end{equation}
where $\{\delta \}$ runs over a complete set of representatives of
$\overline{A({\mathbb{A}})}^1$-conjugacy classes of $A({\mathbb{Q}})$.

This proves the

\begin{theorem}\label{thm:3.9}
(trace formula on $\overline{A({\mathbb{A}})}^1$):
If $\overline{f}\in {\mathcal{S}}(\overline{A({\mathbb{A}})}^1)$ then 
\[
\sum_{\overline{\chi }\in \Pi_u (\overline{A({\mathbb{A}})}^1)}
m(\overline{\chi })
tr(\overline{\chi }(\overline{f}))
=\sum_{\{\delta \}}
vol(A({\mathbb{Q}})_\delta \backslash \overline{A({\mathbb{A}})}_\delta ^1)
\sum_{x\in \overline{A({\mathbb{A}})}_\delta ^1
\backslash \overline{A({\mathbb{A}})^1}}
\overline{f}(x^{-1}\delta x),
\]
where $\{\delta \}$ runs over a complete set of representatives of
$\overline{A({\mathbb{A}})}^1$-conjugacy classes of
$A({\mathbb{Q}})$.
\end{theorem}

\subsection{proof of the correspondence}\label{ssec:rank0corr}

Let $\delta\in A({\mathbb{Q}})$.

We now prove the correspondence theorem mentioned at the beginning of this
subsection. To do this we need to relate the above twisted trace
formula on $\overline{A({\mathbb{A}})}^1$ to the ordinary trace
formula (Poisson summation) on $A({\mathbb{A}})^N$. This boils down to
relating the metaplectic orbital integrals 
\[
\int_{\overline{A({\mathbb{A}})}_\delta ^1
\backslash \overline{A({\mathbb{A}})^1}}
\overline{f}(x^{-1}\delta x)dx
\]
on $\overline{A({\mathbb{A}})}^1$ to ordinary orbital integrals
on $A({\mathbb{A}})^N$ (which are ``trivial'' since
$\overline{A({\mathbb{A}})}$ is abelian).

Let $\overline{f}$ be factorizable: 
\[
\overline{f}(x)=\prod_{p\leq \infty}\overline{f}_p(x_p),
\;\;\;\;\;\overline{f}_p\in {\mathcal{S}}(\overline{A({\mathbb{Q}}_p)}),
\]
so 
\begin{equation}\label{eq:3.15}
\int_{\overline{A({\mathbb{A}})}_\delta ^1
\backslash \overline{A({\mathbb{A}})^1}}
\overline{f}(x^{-1}\delta x)dx
=\prod_{p\leq \infty}
\sum_{x_p\in {\mathbb{Q}}_{p,\delta}^\times
\backslash {\mathbb{Q}}_p^\times}
\overline{f}_p(\delta ,(\delta,x_p)_n^2s_{\mathbb{A}}(\delta)^{-1} ),
\end{equation}
where
\[
\begin{array}{c}
{\mathbb{Q}}_{p,\delta}^\times
=\{x_p\in {\mathbb{Q}}_p^\times
\ |\ (x_p,\delta )_n=1\},\\
(x_p,\delta)_n=\prod_{v|p}(x_v,\delta_v)_{v,n},
\end{array}
\]
for $x_p=\prod_{v|p}x_v,\ \delta=\prod_{v|p}\delta_v$.
 
Let 
\begin{equation}\label{eq:3.16}
f_p(y)=
\sum_{x\in {\mathbb{Q}}_{p,y}^\times
\backslash {\mathbb{Q}}_p^\times}
\overline{f}_p(y ,(y,x)_n^2s_{\mathbb{A}}(y)^{-1} ),
\end{equation}
for $y\in \prod_{v|p}k_v^\times\cong A({\mathbb{Q}}_p)$.

\begin{remark}\label{rem:3.10}
Let $(h(t),1)\in \overline{A({\mathbb{Q}}_p}$ be the element
such that $h(t)$ is diagonal with diagonal entries $t,t^{-1}$.
For $y\in {\mathbb{Q}}_p^\times$, we have
$(h(y),1)^{(h(x),1)}=
(h(y) ,(y,x)_n^2) $.
Therefore, if $\overline{f}_p$ is genuine then
the sum in \eqref{eq:3.16} is zero unless
$y\in (\prod_{v|p}k_v^\times)^N\cong A({\mathbb{Q}}_p)^N$.
\end{remark}

\begin{lemma}\label{lem:3.11}
One may choose Haar measures so that

(a)\ If $p<\infty $\ then the map $\overline{f}_p\longmapsto f_p$
defines a map ${{\mathcal{S}}}(\overline{A({{\mathbb{Q}}}_p)})\rightarrow
{{\mathcal{S}}}(A({{\mathbb{Q}}}_p))$.

(b)\ If $p=\infty$ then the map
$\overline{f}_p\longmapsto f_p$\ defines a map
${{\mathcal{S}}}(\overline{A(\mathbb{R})})
\rightarrow {{\mathcal{S}}}(A({\mathbb{R}}))$.

(c) If $p<\infty $ and if $\overline{f}_p$
is a unit in the convolution
algebra ${\mathcal{S}}(\overline{A({\mathbb{Q}}_p)})$\ and
if $\overline{f}_p\longmapsto f_p$ then $f_p$
is the unit element of
the convolution algebra ${\mathcal{S}}(A({\mathbb{Q}}_p)^N)$.
\end{lemma}

For $\overline{f}\in {\mathcal{S}}(\overline{A({\mathbb{A}})^1})$
factorizable, let $\overline{f}\longmapsto f$ denote the map
obtained by taking the
product of the maps $\overline{f}_p\longmapsto f_p$.
By the above lemma, this defines a map 
\begin{equation}\label{eq:3.17}
{\mathcal{S}}(\overline{A({\mathbb{A}})}^1)\rightarrow {\mathcal{S}}(C_{{\mathbb{A}}}).
\end{equation}

As a corollary to the above lemma, we obtain

\begin{lemma}\label{lem:3.12}
Let\ $\overline{f}\in {\mathcal{S}}(\overline{A({\mathbb{A}})}^1)$ be a
factorizable function and define
$f\in {\mathcal{S}}(C_{\mathbb{A}})$ as above via
$\overline{f}\longmapsto f$. Then 
\[
\int_{\overline{A({\mathbb{A}})}_\delta ^1
\backslash \overline{A({\mathbb{A}})^1}}\overline{f}(x^{-1}\delta x)dx
=f(\delta ),
\]
if $\delta\in A({\mathbb{A}})^N$ and
\[
\int_{\overline{A({\mathbb{A}})}_\delta ^1
\backslash \overline{A({\mathbb{A}})^1}}\overline{f}(x^{-1}\delta x)dx
=0,
\]
if $\delta\in A({\mathbb{A}})-A({\mathbb{A}})^N$.
\end{lemma}

Note that if $\delta\in A({\mathbb{Q}})^N$ then
$\overline{A({\mathbb{A}})^1}_\delta=\overline{A({\mathbb{A}})^1}$.
We choose global measures so that
\[
vol(A({\mathbb{Q}}) \backslash \overline{A({\mathbb{A}})^1})
=vol(A({\mathbb{Q}})^N \backslash A({\mathbb{A}})^{1,N}).
\]
This and the above lemma imply that 
\begin{equation}\label{eq:3.18}
\begin{array}{c}
\sum_{\{\delta \}}
vol(A({\mathbb{Q}})_\delta \backslash \overline{A({\mathbb{A}})^1}_{\delta })
\int_{\overline{A({\mathbb{A}})^1}_{\delta }\backslash 
\overline{A({\mathbb{A}})^1}}
\overline{f}(x^{-1}\delta x)dx \\ 
=vol(A({\mathbb{Q}})^N \backslash A({\mathbb{A}})^{1,N})
\sum_{\gamma \in A({\mathbb{Q}})^N}
f(\gamma ),
\end{array}
\end{equation}
by equation \eqref{eq:3.16} and the remark above.

In this case, we have, by the Poisson summation formula for
$A({\mathbb{Q}})^N\backslash A({\mathbb{A}})^{1,N}$, we have 
\begin{equation}\label{eq:3.19}
\sum_{\overline{\chi }\in
\Pi_u (\overline{A({\mathbb{A}})^1})}m(\overline{\chi })
tr(\overline{\chi }(\overline{f}))
=\sum_{\chi \in \Pi (A({\mathbb{A}})^{1,N})}
f^{\wedge }(\chi ),
\end{equation}
for each factorizable
$\overline{f}\in {\mathcal{S}}(\overline{A({\mathbb{A}})}^1)$
which is genuine and $f\in {\mathcal{S}}(A({\mathbb{A}})^{1,N})$
is defined as above
via $\overline{f}\longmapsto f$.

Our correspondence theorem follows immediately from the two claims below:

\textsc{Claim 1}: For each fixed $\overline{\chi }$
$\in \Pi (\overline{A({\mathbb{A}})}^1)$ there is a sequence of
functions
$\overline{f}_m
=\overline{f}_{\overline{\chi },m}$ such that 
\[
\lim_{n\rightarrow \infty }
tr(\overline{\chi ^{\prime }}(\overline{f}_m))
=\left\{ 
\begin{array}{c}
1,\;\;\overline{\chi ^{\prime }}=\overline{\chi }, \\ 
0,\;\mathrm{otherwise}.
\end{array}
\right. 
\]

\textsc{Claim 2}: For each fixed
$\overline{\chi }\in \Pi (\overline{A({\mathbb{A}})}^1)$
there is a unique $\chi \in \Pi (A({\mathbb{A}})^{1,N})$ such that
$\overline{f}_m\longmapsto f_n=f_{\chi ,m}$ and 
\[
\lim_{m\rightarrow \infty }f_m^{\wedge }(\chi ^{\prime })=\left\{ 
\begin{array}{c}
1,\;\;\;\chi ^{\prime }=\chi , \\ 
0,\;\mathrm{otherwise}.
\end{array}
\right. 
\]

Claim 1 is an exercise in harmonic analysis on $A({\mathbb{A}})^1$ since
$tr(\overline{\chi ^{\prime }}(\overline{f}_m))
=\overline{f}_m^{\wedge }(\overline{\chi ^{\prime }})$.

To prove claim 2 we fix a genuine factorizable function $f$ on $A({\mathbb{A}})$
and work locally. If 
\[
\overline{\chi }=
Ind_{\overline{A({\mathbb{Q}}_p)^{1,N}}}^{\overline{A({\mathbb{Q}}_p)^1}}
(\chi )\in \Pi_u (\overline{A({\mathbb{A}})}^1)
\]
then we have 
\begin{equation}\label{eq:3.20}
\begin{array}{c}
\overline{f}^{\wedge }(\overline{\chi })
=\int_{\overline{A({\mathbb{Q}}_p)}}
\overline{f}(y)tr(\overline{\chi }(y))dy \\ 
=\sum_{h\in \overline{A({\mathbb{Q}}_p)^N}
\backslash \overline{A({\mathbb{Q}}_p)}}
\int_{\overline{A({\mathbb{Q}}_p)^N}}
\overline{f}(a)\chi ^h(a)da \\ 
=\sum_{h\in \overline{A({\mathbb{Q}}_p)^N}
\backslash \overline{A({\mathbb{Q}}_p)}}
\int_{\overline{A({\mathbb{Q}}_p)^N}}
\overline{f}(h^{-1}ah)\chi (a)da \\ 
=\int_{A({\mathbb{Q}}_p)^N\times 1}
f(a)\chi (a)da \\ 
=f^{\wedge }(\chi ).
\end{array}
\end{equation}

Claim 2 follows.

The 1-1 correspondence between $\Pi (\overline{A({\mathbb{A}})^1})$ and
$\Pi (A({\mathbb{A}})^{1,N})=\Pi (\overline{A({\mathbb{A}})^{1,N}})_{gen}$
is given as follows. The map 
\[
\Pi (\overline{A({\mathbb{A}})^{1,N}})
\rightarrow \Pi (\overline{A({\mathbb{A}})}^1)
\]
is the induction map. In fact, the map
$\Pi (\overline{A({\mathbb{A}})^1})\rightarrow
\Pi (A({\mathbb{A}})^{1,N})$
is given locally by the map 
\[
\overline{\chi }
=Ind_{\overline{A({\mathbb{Q}}_p)^N}}^{\overline{A({\mathbb{Q}}_p)}}
(\chi )\longmapsto \chi 
\]
This proves the desired correspondence. $\Box $

Let ${\mathcal{O}}$ denote a (singleton) conjugacy
class in $A({\mathbb{Q}})$ and let
$\chi\in \Pi_u(\overline{A({\mathbb{A}})})$ denote a genuine irreducible
representation. Define $I^A_{\mathcal{O}}$ by
\begin{equation}\label{eq:3.21}
I^A_{\mathcal{O}}(f)=
\sum_{\delta\in {\mathcal{O}}}
f(\delta)
\end{equation}
and $I^A_{\chi}$ by
\begin{equation}\label{eq:3.22}
I^A_{\chi}(f)=
{\rm tr}(\chi(f)) 
\end{equation}
for $f\in {\mathcal{H}}(\overline{A({\mathbb{A}})})$, where
${\mathcal{H}}(\overline{A({\mathbb{A}})})$ is as in \S\ref{sssec:defs} below.



\section{The geometric side of the metaplectic trace formula}\label{sec:geometric}

We follow \cite{A2,A3}.

\subsection{The kernel function}

Let ${\mathcal{C}}$ denote the set of equivalence classes of
$\gamma \in G({\mathbb{Q}})$ as in \cite[section 2]{A2}.
Recall that ${\mathcal{C}}$ is in bijection
with the set of equivalence classes of pairs $(M,{\mathcal{O}}_M)$,
where $P=MN $ is a standard parabolic and
${\mathcal{O}}_M\subset M({\mathbb{Q}})$ is a
conjugacy class of elliptic elements.

For ${\mathcal{O}}\in {\mathcal{C}}$, let 
\[
K_{{\mathcal{O}}}(x,y)=\sum_{\gamma \in {\mathcal{O}}}f(x\gamma y^{-1}),
\]
and more generally, let 
\begin{equation}\label{eq:4.1}
K_{P,{\mathcal{O}}}(x,y)=\sum_{\gamma \in M({\mathbb{Q}})\cap {\mathcal{O}}}
\int_{N({\mathbb{A}})}f(x\gamma y^{-1})dn,
\end{equation}
where $P$ is a parabolic as above.

\begin{lemma}\label{lem:4.1}
For $\overline{K_0}$-finite
$f\in {\mathcal{H}}(\overline{G({\mathbb{A}})})$, the sum
$K_P(x,y)=\sum_{{\mathcal{O}}\in {\mathcal{C}}}K_{{\mathcal{O}}}(x,y)$
converges to the kernel
function of $R_P$, where $R_P$ denotes the right
regular representation of
${\mathcal{H}}(\overline{G({\mathbb{A}})})_{gen}$ on
$L^2(N({\mathbb{A}})M({\mathbb{Q}})\backslash \overline{G({\mathbb{A}})})_{gen}$.
\end{lemma}

proof: For
$\phi \in C_c^\infty (N({\mathbb{A}})M({\mathbb{Q}})
\backslash \overline{G({\mathbb{A}})})_{gen}$, we have 
\begin{eqnarray*}
(R_P(f)\phi )(x) &=&\int_{\overline{G({\mathbb{A}})}}\phi (g)f(xg^{-1})dg \\
&=&\int_{\overline{G({\mathbb{A}})}}\phi (n^{-1}\gamma ^{-1}g)f(xg^{-1})dg \\
&=&\int_{\overline{G({\mathbb{A}})}}\phi (g)f(xn\gamma g^{-1})dg \\
&=&\int_{M({\mathbb{Q}})\backslash \overline{G({\mathbb{A}})}}\phi (g)\sum_{\gamma \in
M({\mathbb{Q}})}f(xn\gamma g^{-1})dg \\
&=&\int_{N({\mathbb{A}})M({\mathbb{Q}})\backslash \overline{G({\mathbb{A}})}}\phi
(g)\sum_{\gamma \in M({\mathbb{Q}})}\int_{N({\mathbb{A}})}f(xn\gamma g^{-1})dndg.
\end{eqnarray*}

The convergence of the inner sum is established, for example, by passing to
the place at infinity and carrying out the estimates in the usual manner.
For the description of
$G({\mathbb{Q}})\cap (\overline{G({\mathbb{R}})\times K^f})$,
see \cite{K2} for example.$\Box $

\subsubsection{Some analytic reduction theory}

Let 
\[
N({\mathbb{A}})_u=\{\left( 
\begin{array}{cc}
1 & n \\ 
0 & 1
\end{array}
\right) \;|\;|n|_{{\mathbb{A}}}\leq u\}. 
\]
For $\omega \subset N({\mathbb{A}})_u\cdot A({\mathbb{A}})^1$
relatively compact, $T_0\in {\mathfrak{a}}^{+}$, let 
\[
S(T_0,\omega )=
\{nak\in \overline{\omega \cdot A({\mathbb{R}})^0\cdot K_0}
\;|\;\alpha (H(a)-T_0)>0\}. 
\]
This is a Siegel domain of $\overline{G({\mathbb{A}})}$. It has finite volume
with respect to the measure $\delta (a)dndadk$. For $T_0\in {\mathfrak{a}}^{+}$
and $\omega \subset N({\mathbb{A}})_u\cdot A({\mathbb{A}})^1$ relatively
compact both
``sufficiently large'' ($u\geq 1/2,\;T_0>\sqrt{3}/2$), we have 
\[
\overline{G({\mathbb{A}})}=G({\mathbb{Q}})\cdot S(T_0,\omega ). 
\]
In particular, $G({\mathbb{Q}})\backslash \overline{G({\mathbb{A}})}$ is
finite volume.

\begin{lemma}\label{lem:4.2}
There exists an $N>0$ such that for all $g\in \overline{G({\mathbb{A}})}$,
$T_0\in {\mathfrak{a}}^{+}$, we have 
\[
\sum_{\delta \in B({\mathbb{Q}})\backslash G({\mathbb{Q}})}
\tau _B^{\wedge }(H(\delta x)-T_0)\leq 1.
\]
\end{lemma}

This is an immediate consequence of the corresponding result in \cite{A2}.

For $T_0>T_1$ in ${\mathfrak{a}}^{+}$, let 
\[
S(T_0,T_1,\omega )=\{nak\in S(T_0,\omega )\;|\;\omega (H(a)-T_0)\leq 0\}. 
\]
Let $F^G(x,T_1,T_0)$ denote the characteristic function of the image
$\overline{S(T_0,T_1,\omega )}\subset
G({\mathbb{Q}})\backslash \overline{G({\mathbb{A}})}$ of
$S(T_0,T_1,\omega )$ under the quotient map
$\overline{G({\mathbb{A}})}\rightarrow
G({\mathbb{Q}})\backslash \overline{G({\mathbb{A}})}$. Let
$F^B(x,T_1,T_0)$
denote the characteristic function of the set of all
$xB({\mathbb{Q}})\in B({\mathbb{Q}})\backslash \overline{G({\mathbb{A}})}$ such
that $\delta x\in \overline{\omega \cdot A({\mathbb{R}})^o\cdot K_0}$
for some $\delta \in B({\mathbb{Q}})$.

\begin{lemma}\label{lem:4.3}
For $\omega \subset N({\mathbb{A}})_u\cdot A({\mathbb{A}})^1$ relatively
compact and
$T_0\in {\mathfrak{a}}^{+}$ sufficiently large,
$x\in \overline{G({\mathbb{A}})}$, we
have 
\[
\sum_{P\in \{B,G\}}
\sum_{\delta \in P({\mathbb{Q}})\backslash G({\mathbb{Q}})}
F^P(\delta x,T_1,T_0)\tau _P(H(\delta x)-T_1)=1.
\]
\end{lemma}

This is an immediate consequence of the corresponding result in \cite{A2}.

\subsubsection{The truncated kernel}

The auxiliary kernel function is by definition 
\[
k^T(x,f)=\sum_{P\in \{B,G\}}(-1)^{\dim (A_P)}
\sum_{\delta \in P({\mathbb{Q}})\backslash G({\mathbb{Q}})}
K_P(\delta x,\delta x)\tau _P^{\wedge }(H(\delta x)-T). 
\]
This sum is finite by Lemma~\ref{lem:4.2} above. Using the previous Lemma~\ref{lem:4.3},
this may be rewritten as 
\begin{eqnarray*}
k^T(x,f) &=&\sum_{P,Q\in \{B,G\}}(-1)^{\dim (A_P)}\times \\
&&\;\;\;\;\;\;\;\;\;\;\;\times
\sum_{\delta \in Q({\mathbb{Q}})\backslash G({\mathbb{Q}})}
K_Q(\delta x,\delta x)\tau _Q^{\wedge }(H(\delta x)-T)
F^P(\delta x,T,T_0)\tau _P^Q(H(\delta x)-T),
\end{eqnarray*}
for $T>T_0$ regular. This may be rewritten as
\begin{eqnarray}
k^T(x,f) &=&K_G(x,x)F^G(x,T,T_0)+ \label{eq:4.2}\\
&&\sum_{\delta \in B({\mathbb{Q}})\backslash G({\mathbb{Q}})}
\tau _B(H(\delta x)-T)
F^B(\delta x,T,T_0)(K_G(\delta x,\delta x)-K_B(\delta x,\delta x)). 
\nonumber
\end{eqnarray}

We want to prove

\begin{lemma}\label{lem:4.4}
$\int_{G({\mathbb{Q}})\backslash \overline{G({\mathbb{A}})}}|k^T(x,f)|dx<\infty .$
\end{lemma}

We majorize this by the sum of the integrals of the absolute values of the
two terms on the right hand side of \eqref{eq:4.2}. The integral over the first term
on the right hand side of \eqref{eq:4.2}, 
\[
\int_{G({\mathbb{Q}})\backslash \overline{G({\mathbb{A}})}}|K_G(x,x)F^G(x,T,T_0)|dx 
\]
clearly converges since $F^G(x,T,T_0)$ is compactly supported. To estimate
the second term, we first write that 
\begin{eqnarray}
&&\int_{G({\mathbb{Q}})\backslash \overline{G({\mathbb{A}})}}|
\sum_{\delta \in B({\mathbb{Q}})\backslash G({\mathbb{Q}})}
\tau _B(H(\delta x)-T)\times \label{eq:4.3}\\
&&\;\;\;\;\;\;\;\;\;\;\;\;\;\;\;\;\;\;\;\;\times F^B(\delta
x,T,T_0)(K_G(\delta x,\delta x)-K_B(\delta x,\delta x))|dx \label{eq:4.4}\\
&<&<\int_{G({\mathbb{Q}})\backslash \overline{G({\mathbb{A}})}}
\tau _B(H(\delta x)-T)F^B(\delta x,T,T_0)\times  \nonumber \\
&&\;\;\;\;\;\;\;\;\;\;\;\;\;\;\;\;\;\;\;\;\;\;\;\times
|\sum_{\delta \in G({\mathbb{Q}})}f(x^{-1}\delta x)
-\int_{N({\mathbb{A}})}f(x^{-1}\mu n^{\prime }x)dn^{\prime }|dx.
\label{eq:4.5}\end{eqnarray}
Next, we again pass to the archimedian place. Let 
\begin{eqnarray}
\Gamma _0 &=&N({\mathbb{Q}})\cap (\overline{G({\mathbb{R}})\times K^f}),
\;\Gamma_\infty
=N({\mathbb{R}})\cap (\overline{G({\mathbb{R}})\times K_0^f}),\; \label{eq:4.6}\\
A({\mathbb{R}})_T &=&\{\prod_{v|\infty}
\left( 
\begin{array}{cc}
t_v & 0 \\ 
0 & t_v^{-1}
\end{array}
\right) \in A({\mathbb{R}})\;|\;|t_v|>T^{1/2}\},
\;G({\mathbb{R}})_T=N({\mathbb{R}})\cdot A({\mathbb{R}})_T\cdot K_\infty ,  
\nonumber
\end{eqnarray}
so $\Gamma _0\subset \Gamma _\infty $ is a co-compact discrete abelian
group. The above upper bound \eqref{eq:4.5} is finite if 
\begin{equation}\label{eq:4.7}
\int_{\Gamma _0\backslash \overline{G({\mathbb{R}})_T}}|\sum_{\gamma \in \Gamma
_0}f(x^{-1}\gamma x)-\int_{\Gamma _\infty }f(x^{-1}nx)dn|dx<\infty .
\end{equation}
We may apply the theory of Fourier transforms, as in \cite{W1}, to the 
quotient $\Gamma _0\backslash \Gamma _\infty $. 
Choosing Haar measures appropriately,
there is a non-trivial character 
$\psi $ of $\Gamma _\infty $ such that the Fourier transform 
\[
\varphi ^{\wedge }(\gamma )
=\int_{\Gamma _\infty }\varphi (\gamma ^{\prime})
\psi (<\gamma ,\gamma ^{\prime }>)d\gamma ^{\prime },
\;\;\;\varphi \in L^2(\Gamma _0\backslash \Gamma _\infty ), 
\]
satisfies 
\[
\varphi ^{\wedge }(1)=\int_{\Gamma _\infty }\varphi (\gamma ^{\prime
})d\gamma ^{\prime }. 
\]

Take $\varphi (\gamma )=\varphi _x(\gamma )=f(x^{-1}\gamma x)$, so by the
Poisson summation formula \cite{W1}, 
\[
\sum_{\gamma \in \Gamma _0}f(x^{-1}\gamma x)-\int_{\Gamma _\infty
}f(x^{-1}nx)dn=\sum_{\gamma ^{*}\in \Gamma _0^{*}}\varphi _x^{\wedge
}(\gamma ^{*}), 
\]
where $\Gamma _0^{*}$ denotes the dual lattice of $\Gamma _0$. Since $f$ is
smooth, $\varphi _x^{\wedge }$ is rapidly decreasing. For sufficiently
regular $T$, standard arguments now suffice to show that, for any constant $%
C>0$, we have 
\[
\int_{\Gamma _0\backslash \overline{G({\mathbb{R}})_T}}|\sum_{\gamma ^{*}\in
\Gamma _0^{*}}\varphi _x^{\wedge }(\gamma ^{*})|dx<<_{C,f,n}T^{-C}. 
\]

This verifies \eqref{eq:4.7}, which proves the lemma. $\Box $

The above proof actually establishes more, namely,

\begin{corollary}\label{cor:4.5}
(of the above proof) For each constant $C>0$, we have 
\[
\int_{G({\mathbb{Q}})\backslash \overline{G({\mathbb{A}})}}k^T(x,f)dx-\int_{G({\mathbb{Q}}%
)\backslash \overline{G({\mathbb{A}})_T}}K^G(x,x)dx<<_{C,f,n}T^{-C}. 
\]
Here 
\[
G({\mathbb{A}})_T=N({\mathbb{A}})_u\cdot A({\mathbb{A}})^T\cdot K 
\]
where $A({\mathbb{A}})^T=\{a\in A({\mathbb{A}})\;|\;|\delta (a)|_{\mathbb A}\leq T\}$.
\end{corollary}

In fact, the proof yields more than this. If we let 
\[
k_{{\mathcal{O}}}^T(x,f)=\sum_{P\in \{B,G\}}(-1)^{\dim (A_P)}\sum_{\delta \in
P({\mathbb{Q}})\backslash G({\mathbb{Q}})}K_{P,{\mathcal{O}}}(\delta x,\delta x)\tau
_P^{\wedge }(H(\delta x)-T), 
\]
for ${\mathcal{O}}\in {\mathcal{C}}$, then the above proof establishes

\begin{corollary}\label{cor:4.6}
(of the above proof) For each constant $C>0$, we have 
\[
\sum_{{\mathcal{O}}\in {\mathcal{C}}}\int_{G({\mathbb{Q}})\backslash \overline{G(\mathbb{%
A})}}|k_{{\mathcal{O}}}^T(x,f)dx-F^G(x,T,T_0)K_{G,{\mathcal{O}}%
}(x,x)|dx<<_{C,f,n}T^{-C}. 
\]
\end{corollary}

\subsection{Computing the integral of $k^T(x,f)$}

For a class ${\mathcal{O}}\in {\mathcal{C}}$ and a standard parabolic $P=MN$,
let 
\[
J_{P,{\mathcal{O}}}(x,y)=\sum_{\mu \in M({\mathbb{Q}})\cap {\mathcal{O}}%
}\sum_{\varsigma \in N_{\mu _s}({\mathbb{Q}})\backslash N({\mathbb{Q}})}\int_{N_{\mu
_s}({\mathbb{A}})}f(x^{-1}\varsigma ^{-1}\mu n\varsigma y)dn, 
\]
where $\mu =\mu _s\mu _u$ denotes the Jordan decomposition (over $\overline{%
{\mathbb{Q}}}$), and let 
\[
j_{{\mathcal{O}}}^T(x,f)=\sum_{P\in \{B,G\}}(-1)^{\dim (A_P)}\sum_{\delta \in
P({\mathbb{Q}})\backslash G({\mathbb{Q}})}J_{P,{\mathcal{O}}}(\delta x,\delta x)\tau
_P^{\wedge }(H(\delta x)-T). 
\]
the above proof establishes

\begin{corollary}\label{cor:4.7}
(of the above proof) For each constant $C>0$, we have 
\[
\sum_{{\mathcal{O}}\in {\mathcal{C}}}\int_{G({\mathbb{Q}})\backslash \overline{G(\mathbb{%
A})}}|j_{{\mathcal{O}}}^T(x,f)dx-F^G(x,T,T_0)J_{G,{\mathcal{O}}%
}(x,x)|dx<<_{C,f,n}T^{-C}. 
\]
\end{corollary}

Let 
\begin{eqnarray*}
J_{{\mathcal{O}}}^T(f) &=&\int_{G({\mathbb{Q}})\backslash \overline{G({\mathbb{A}})}}j_{%
{\mathcal{O}}}^T(x,f)dx, \\
K_{{\mathcal{O}}}^T(f) &=&\int_{G({\mathbb{Q}})\backslash \overline{G({\mathbb{A}})}}k_{%
{\mathcal{O}}}^T(x,f)dx.
\end{eqnarray*}

\begin{proposition}\label{prop:4.8}
Both $
J_{{\mathcal{O}}}^T(f)
$ and $K_{{\mathcal{O}}}^T(f)$ are polynomials in
$T$.
\end{proposition}

Define
\begin{equation}\label{eq:4.8}
J_{{\mathcal{O}}}(f)
=J_{{\mathcal{O}}}^T(f)|_{T=0}.
\end{equation}

We will prove this below. First, we record a

\begin{corollary}\label{cor:4.9}
$\sum_{{\mathcal{O}}\in {\mathcal{C}}}J_{{\mathcal{O}}}^T(f)=\sum_{{\mathcal{O}}\in 
{\mathcal{C}}}K_{{\mathcal{O}}}^T(f)$.
\end{corollary}

proof: The difference between them is a polynomial, $p$ say, in $T$ by the
proposition. Since $J_{G,{\mathcal{O}}}(x,y)=K_{G,{\mathcal{O}}}(x,y)$, this
polynomial $p$ must be the zero polynomial by Corollaries~\ref{cor:4.6}
and~\ref{cor:4.7}. $\Box $

Let ${\mathcal{O}}\in {\mathcal{C}}$ denote an unramified semi-simple class \cite[p. 949]{A2}. If $\mu \in {\mathcal{O}}$ then its centralizer $G_\mu $ is either $G$
or a Cartan subgroup of $G$ defined over ${\mathbb{Q}}$.

As in \cite{A2}, we have 
\begin{equation}\label{eq:4.9}
j_{{\mathcal{O}}}^T(x,f)=
\sum_{\eta \in G_\mu ({\mathbb{Q}})\backslash G({\mathbb{Q}})}
f(x^{-1}\eta ^{-1}\mu \eta x)\sum_{P\in \{B,G\}}(-1)^{\dim (A_P)}
\tau_P^{\wedge }(H(\eta x)-T),
\end{equation}
so 
\begin{eqnarray}
J_{{\mathcal{O}}}^T(f)
&=&\int_{G_\mu ({\mathbb{Q}})\backslash \overline{G({\mathbb{A}})}}
f(x^{-1}\mu x)\sum_{P\in \{B,G\}}(-1)^{\dim (A_P)}
\tau _P^{\wedge}(H(x)-T)dx \label{eq:4.10}\\
&=&vol(A_{{\mathcal{O}}}({\mathbb{R}})^0G_\mu ({\mathbb{Q}})
\backslash \overline{G_\mu ({\mathbb{A}})}
\int_{\overline{G_\mu ({\mathbb{A}})}\backslash \overline{G({\mathbb{A}})}}
f(x^{-1}\mu x)v_{{\mathcal{O}}}(x,T)dx,  \nonumber
\end{eqnarray}
where $A_{{\mathcal{O}}}$ denotes the center of $M_\mu $ and 
\[
v_{{\mathcal{O}}}(x,T)
=\int_{A_{{\mathcal{O}}}({\mathbb{R}})^0}\sum_{P\in
\{B,G\}}(-1)^{\dim (A_P)}\tau _P^{\wedge }(H(ax)-T)da. 
\]
Note that the global cocycle $\beta $ splits over
$A_{{\mathcal{O}}}({\mathbb{R}})^o\subset A({\mathbb{R}})^o$.
In the case where $A_{\mathcal{O}}=A$, this may be
computed to be 
\[
v_{{\mathcal{O}}}(x,T)=2(T-H(x)). 
\]

If $\mu \in {\mathcal{O}}$ is semi-simple but not regular then 
\[
j_{{\mathcal{O}}}^T(x,f)
=f(\mu )\sum_{P\in \{B,G\}}(-1)^{\dim (A_P)}
\tau_P^{\wedge }(H(x)-T) 
\]
so 
\begin{equation}
J_{{\mathcal{O}}}^T(f)=f(\mu )
\int_{G({\mathbb{Q}})\backslash \overline{G({\mathbb{A}})}}
\sum_{P\in \{B,G\}}(-1)^{\dim (A_P)}
\tau _P^{\wedge }(H(x)-T)dx.  \nonumber
\end{equation}



\subsection{Basic properties}

\subsubsection{Definitions}\label{sssec:defs}

Let $S$ be a subset of places of ${\mathbb{Q}}$ and let
$||...||:G({\mathbb{A}}_S)\rightarrow {\mathbb{R}}_+$ be a ``height'' function
as in Section~\ref{sec:background}. Let $N>0$ and let
$\Gamma\subset \Pi(\overline{K})$
be a finite set of inequivalent genuine irreducible representations
of the n-fold cover of $K=K_{0,S}=\prod_{p\in S}K_{0,p}$.

We call a genuine factorizable function
$f=\prod_{p\in S}f_p:\overline{G({\mathbb{A}}_S)}\rightarrow {\mathbb{C}}$
{\bf smooth} if

(a) $f_\infty\in C^\infty(\overline{G({\mathbb{R}})})_{gen}$
if $p=\infty \in S$, and

(b) $f_p\in C_c^\infty(\overline{G({\mathbb{Q}}_p)})_{gen}$
if $p\in S-\{\infty\}$.

\noindent
In general, a genuine
function $f:\overline{G({\mathbb{A}}_S)}\rightarrow {\mathbb{C}}$
is {\bf smooth} if it is a finite linear combination of factorizable smooth
functions. Define
${\mathcal{H}}_N(\overline{G({\mathbb{A}}_S)})_\Gamma$
to be the vector space of
all genuine smooth complex valued functions
$f:\overline{G({\mathbb{A}}_S)}\rightarrow {\mathbb{C}}$
such that
\begin{equation}\label{eq:4.11}
\begin{array}{c}
{\rm supp}(f)\subset \overline{G({\mathbb{A}}_S,N)}
=\{ g\in \overline{G({\mathbb{A}}_S)}\ |\ ||g||\leq e^N\}\\
{\rm The \ spans\ of\ }
k\longmapsto f(kg){\rm \ generates\ a \ repn\ in\ }\Gamma,\\
{\rm The \ spans\ of\ }
k\longmapsto f(gk){\rm \ generates\ a \ repn\ in\ }\Gamma .
\end{array}
\end{equation}
The semi-norms
\[
||f||_D=\sup_{x\in \overline{G({\mathbb{A}}_S)}} |Df(x)|,
\]
$D$ a differential operator on $\overline{G({\mathbb{R}})}$,
induces a topology on
${\mathcal{H}}_N(\overline{G({\mathbb{A}}_S)})_\Gamma$
making it into a topological
vector space. By \cite{Ru}, ch 1,
we may complete this space using the above
family of seminorms and obtain a Frech\'et space.

Define
\[
{\mathcal{H}}(\overline{G({\mathbb{A}}_S)})=\lim_{\Gamma}
{\mathcal{H}}(\overline{G({\mathbb{A}}_S)})_\Gamma
\]
where
\[
{\mathcal{H}}(\overline{G({\mathbb{A}}_S)})_\Gamma =\lim_{N}
{\mathcal{H}}_N(\overline{G({\mathbb{A}}_S)})_\Gamma .
\]
Therefore, each function in
${\mathcal{H}}(\overline{G({\mathbb{A}}_S)})$
is both right and left $\overline{K}$-finite.

Similarly, if $A$ denotes the usual
diagonal torus of $G=SL(2)$, then we define
${\mathcal{H}}(\overline{A({\mathbb{A}}_S)})$.

From \cite[section 11]{A11}, another space of functions of
``almost compact'' support,
${\mathcal{H}}_{ac}(\overline{G({\mathbb{A}}_S)})$, is introduced,
as well as its subspace of ``moderate'' functions,
which we denote by
${\mathcal{H}}_{mod}(\overline{G({\mathbb{A}}_S)})$.
However, in the case $G=SL(2)$ it turns out that
\[
\begin{array}{c}
{\mathcal{H}}_{mod}(\overline{G({\mathbb{A}}_S)})=
{\mathcal{H}}_{ac}(\overline{G({\mathbb{A}}_S)})=
{\mathcal{H}}(\overline{G({\mathbb{A}}_S)}),\\  
{\mathcal{H}}_{mod}(\overline{A({\mathbb{A}}_S)})=
{\mathcal{H}}_{ac}(\overline{A({\mathbb{A}}_S)})=
{\mathcal{H}}(\overline{A({\mathbb{A}}_S)}).
\end{array}
\]

We shall only sketch the definition of the Schwartz
space ${\mathcal{C}}(\overline{G({\mathbb{A}}_S)})$: this is the space spanned by
factorizable functions $f$ each of whose local components
belong to the local Schwartz space, defined for example
in \cite{J1}.
% Claude (2026): the 1997 book cites its Chapter 6 here.
This space will play a bigger role in Section~\ref{sec:invariant} below.

\subsection{An equivalence relation on $G({\mathbb{Q}})$}

Let $\gamma ,\gamma ^{\prime }\in G({\mathbb{Q}})$ and let
$\gamma =\gamma_u\gamma _s$ denote its Jordan decomposition.
We say that $\gamma ,\gamma^{\prime }$ are \textbf{equivalent},
written $\gamma \sim \gamma ^{\prime }$,
if $\gamma _s$ is $G({\mathbb{Q}})$-conjugate to
$\gamma _s^{\prime }$. Clearly
if $\gamma ,\gamma ^{\prime }$ are conjugate and
semi-simple then they are
equivalent. In fact, by the uniqueness of the Jordan decomposition, it can
be shown that if $\gamma ,\gamma ^{\prime }$ are conjugate then they are
equivalent. Define $\mathcal{C}$ to be the set of all equivalence classes.
This definition holds in general.

At least for $G=SL(n)$, the classical analog of this definition is the
following: Let $\gamma ,\gamma ^{\prime }\in G({\mathbb{Z}})$. We say that
$\gamma ,\gamma ^{\prime }$ are \textbf{equivalent}, written $\gamma \sim
\gamma ^{\prime }$, if $\gamma _s$ is $G(\mathbb{C})$-conjugate to $\gamma
_s^{\prime } $. In other words, $\gamma ,\gamma ^{\prime }$ are equivalent
if they have the same set of eigenvalues.

We call $(\gamma ,\zeta)\in \overline{G({\mathbb{A}})}$
semi-simple if $\gamma\in G({\mathbb{A}})$ is semi-simple.
Let ${\mathcal{O}}\in \mathcal{C}$ denote a
class consisting entirely of semisimple
elements (we call such a class \textbf{semi-simple}).
If $\gamma \in {\mathcal{O}}$
then the centralizer $G_\gamma $ is reductive and anisotropic modulo its
center:

\begin{equation}\label{eq:4.12}
\begin{array}{c}
G_\gamma =\left\{ 
\begin{array}{c}
G,\;\;\;\;\;\;\;\;\;\;\;\;\;\;\;\;\;\;\;\;\;\gamma =\pm 1, \\ 
A^x,\;\;\;\gamma \neq \pm 1\;and\;\gamma \sim a\in A, \\ 
T,\;\;\;\;\;\;\;\;\;\;\;\;\;\;\;\;\;\;\mathrm{otherwise},
\end{array}
\right.
\end{array}
\end{equation}
where $A^x=x^{-1}Ax$ and $x^{-1}\gamma x\in A$.
In particular, all such classes are unramified in the sense of
\cite{A2}, p. 949.

Let ${\mathcal{O}}\in \mathcal{C}$ denote an unramified
semi-simple class, let $P=MN\in\{B,G\}$, and let
$\mu \in {\mathcal{O}}\cap M({\mathbb{Q}})$.
Then

\begin{equation}\label{eq:4.13}
\begin{array}{c}
J_{B,{\mathcal{O}}}(\delta x,\delta x)
=\sum_{\eta \in M_\mu ( {\mathbb{Q}})\backslash M(\mathbb{Q})}
\sum_{\varsigma \in N({\mathbb{Q}})}
f(x^{-1}\delta ^{-1}\varsigma ^{-1}\eta^{-1}\mu \eta \varsigma \delta x) \\ 
=\sum_{\eta \in M_\mu ({\mathbb{Q}})\backslash P({\mathbb{Q}})}
f(x^{-1}\delta^{-1}\eta ^{-1}\mu \eta \delta x).
\end{array}
\end{equation}

If $P=B$ then $M=A$ and $M_\mu =A$ for
$\mu \in {\mathcal{O}}$ as above. If $P=G$ then
$M=G$ and $M_\mu =G_\mu $ is one of the above
centralizers for $\mu \in {\mathcal{O}}$ as
above. As we saw above, if $P=B$ then

\begin{equation}\label{eq:4.14}
\begin{array}{c}
J_{B,{\mathcal{O}}}(\delta x,\delta x)
=\sum_{\varsigma \in N( {\mathbb{Q}})}
f(x^{-1}\delta^{-1}\varsigma ^{-1}\mu \varsigma \delta x) \\ 
=\sum_{\eta \in A({\mathbb{Q}})\backslash B({\mathbb{Q}})}
f(x^{-1}\delta ^{-1}\eta^{-1}\mu \eta \delta x).
\end{array}
\end{equation}
We have then by an equation above

\begin{equation}\label{eq:4.15}
\begin{array}{c}
j_{\mathcal{O}}^T(x,f)
=\sum_{P\in \{B,G\}}(-1)^{\dim A_P}
\sum_{\delta \in P( {\mathbb{Q}})\backslash G({\mathbb{Q}})}
J_{P,{\mathcal{O}}}(\delta x,\delta x) \tau _P^{\wedge}(H_B(\delta x)-T) \\ 
=\sum_{P\in \{B,G\}}(-1)^{\dim A_P}
\sum_{\eta \in M_\mu ( {\mathbb{Q}})\backslash G({\mathbb{Q}})}
f(x^{-1}\eta ^{-1}\mu \eta x)\tau _P^{\wedge}(H_B(\eta x)-T) \\ 
=\sum_{\eta \in G_\mu ({\mathbb{Q}})\backslash G({\mathbb{Q}} )}
f(x^{-1}\eta ^{-1}\mu\eta x)
\sum_{P\in \{B,G\}}(-1)^{\dim A_P} \tau_P^{\wedge }(H_B(\eta x)-T).
\end{array}
\end{equation}
Also

\begin{equation}\label{eq:4.16}
\begin{array}{c}
J_{\mathcal{O}}^T(f)
=\int_{G({\mathbb{Q}})\backslash
\overline{G({\mathbb{A}})}}
j_{\mathcal{O}}^T(x)dx \\ 
=\int_{G({\mathbb{Q}})\backslash \overline{G({\mathbb{A}})}}
\sum_{\delta \in G_\mu ({\mathbb{Q}})\backslash G({\mathbb{Q}})}
f(x^{-1}\delta ^{-1}\mu \delta x)
\sum_{P\in\{B,G\}}(-1)^{\dim A_P}\tau _P^{\wedge }(H_B(\delta x)-T)dx \\ 
=\int_{G_\mu ({\mathbb{Q}})\backslash \overline{G({\mathbb{A}})}}
f(x^{-1}\mu x)\sum_{P\in\{B,G\}}(-1)^{\dim A_P}
\tau _P^{\wedge }(H_B(x)-T)dx \\ 
=vol(G_\mu ({\mathbb{Q}})\backslash \overline{G({\mathbb{A}})}_\mu)
\int_{\overline{G({\mathbb{A}})}_\mu\backslash \overline{G({\mathbb{A}})}}
f(x^{-1}\mu x)\sum_{P\in \{B,G\}}(-1)^{\dim A_P}
\tau_P^{\wedge }(H_B(x)-T)dx \\ 
=vol(A_{\mathcal{O}}({\mathbb{R}})^oG_\mu ({\mathbb{Q}})\backslash
\overline{G ({\mathbb{A}})}_\mu)
\int_{\overline{G ({\mathbb{A}})_\mu}\backslash \overline{G({\mathbb{A}})}}
f(x^{-1}\mu x)v(x,T)dx
\end{array}
\end{equation}
(see also \cite[(5.5)]{Sho}), where 
\begin{equation}\label{eq:4.17}
v(x,T)=\int_{A_{\mathcal{O}}({\mathbb{R}})^o}
\sum_{P\in \{B,G\}}(-1)^{\dim A_P}\tau_P^{\wedge }(H_B(ax)-T)da.
\end{equation}

If $A_{\mathcal{O}}=A$, that is to say if there is an
$a\in A({\mathbb{Q}})-\{\pm I_2\}$ in the class ${\mathcal{O}}$
then $v(x,T)$ is the ``volume'' of the convex hull of 
\[
\{T-H(x),-(T-H(wx))\}. 
\]
Thus 
\begin{equation}\label{eq:4.18}
v(x,T)=2T-\log \delta(x)-\log \delta(wx).
\end{equation}

This proves the following

\begin{theorem}\label{thm:4.10}
Let ${\mathcal{O}}\in \mathcal{C}$ and suppose that
${\mathcal{O}}\cap (A({\mathbb{Q}})-\{\pm I_2\})\not=\emptyset $.
Then 
\[
J_{\mathcal{O}}^T(f)
=vol(A_{\mathcal{O}}({\mathbb{R}})^oG_\mu ({\mathbb{Q}})\backslash G_\mu ({\mathbb{A}}))
\int_{G_\mu ({\mathbb{A}})\backslash G({\mathbb{A}})}f(x^{-1}\mu x)v(x,T)dx. 
\]
\end{theorem}

These are called \textbf{weighted orbital integrals}.
The case where $1\in {\mathcal{O}}$ will be discussed below.

\subsection{The definition of $J_M(u,f)$}\label{ssec:JMu}

Let $f\in C_c^\infty (\overline{G({\mathbb {A}}_S)})_{gen}$ and let
$\overline{\gamma}=(\gamma,1) \in \overline{M({\mathbb {A}}_S)}$
have Jordan decomposition
$\gamma =\sigma u$. Suppose that
$\overline{M({\mathbb{A}}_S)}_{\overline{\gamma}}
=\overline{G({\mathbb{A}}_S)}_{\overline{\gamma}} $.
Assume that $S$ is a ``sufficiently large'' finite set which contains
the archimedian places. In this case, define 
\begin{equation}\label{eq:4.19}
J_M(\overline{\gamma} ,f)=|D(\gamma )|^{1/2}
\int_{\overline{M({\mathbb {A}}_S)}
\backslash \overline{G({\mathbb {A}}_S)}}
\int_{{\mathcal O}_{\overline{\gamma}} (\overline{M({\mathbb {A}}_S)})}
f(x^{-1}\mu x)v_M(x)d\mu dx,
\end{equation}
where ${\mathcal O}_{\overline{\gamma}} (\overline{M({\mathbb {A}}_S)})
=\{m^{-1}\overline{\gamma} m\;m\in \overline{M({\mathbb {A}}_S)})$.
Here $v_M$ is the ``volume'' of the convex hull of the points
$\{H_B(x),H_{B^{-}}(x)\}$ in ${\mathfrak a}$. If $x=nk$ then
$H_B(x)=0$, so $v_M(nk)=|H_{B^{-}}(x)|$.
In any case, this integral can be written as 
\[
J_M(\overline{\gamma} ,f)=|D(\gamma )|^{1/2}
\int_{\overline{G({\mathbb {A}}_S)}_{\overline{\gamma}}
\backslash \overline{G({\mathbb{A}}_S)}}
f(x^{-1}\overline{\gamma} x)v_M(x)dx. 
\]
This is a function on $\overline{M({\mathbb {A}}_S)}$.

If $\gamma \in \overline{M({\mathbb {A}}_S)}$
is {\em arbitrary} and $a\in \overline{A_M({\mathbb {A}}_S)}$ is
chosen so that
$M({\mathbb {A}}_S)_{a\gamma }=G({\mathbb {A}}_S)_{a\gamma }$
then $J_M(a\gamma ,f)$ is well-defined. We may consider
$a\rightarrow 1$.

As in \cite[\S 3]{A7}, let 
\[
r_\beta (\lambda ,u,a)
=|a^\beta -a^{-\beta }|^{\rho (\beta ,u)\lambda (\beta^{\vee })},
\;\;\;\;\lambda \in {\mathfrak {a}}_{M,{\mathbb {C}}}^{*}, 
\]
and let 
\[
r_P(\lambda ,u,a)
=\prod_{v\in S}\prod_\beta
r_\beta (\frac \lambda 2,u_v,a_v), 
\]
where $\beta $ runs over the reduced roots of
$(P_{\sigma _v},A_{M_{\sigma_v}})$. By definition, 
\[
r_M(u ,a)=\lim_{\lambda \rightarrow 0}
\sum_{P\in {\mathcal P}(M)}r_P(\lambda ,u,a)\theta _P(\lambda )^{-1}, 
\]
and 
\[
v_M(x)=\lim_{\lambda \rightarrow 0}
\sum_{P\in {\mathcal P}(M)}e^{-\lambda (H_P(x))}
\theta _P(\lambda )^{-1}. 
\]
Following \cite[(8.1)]{A7},
we define $J_M(\gamma ,f)$ to be the limit at $a=1$
of 
\begin{equation}\label{eq:4.20}
\begin{array}{c}
|D(a\gamma )|^{1/2}
\int_{\overline{G({\mathbb {A}}_S)}_\sigma \backslash
\overline{G({\mathbb {A}}_S)}}
\int_{\overline{M({\mathbb {A}}_S)}_\gamma
\backslash \overline{G({\mathbb {A}}_S)}_\sigma }
f(y^{-1}\sigma x^{-1}a u xy)\times \\
\times
\left( \sum_{L\in {\mathcal{L}}(M)}
r_M^L(\gamma ,a)v_L(xy)\right) dxdy,
\end{array}
\end{equation}
where $\gamma\in \overline{M({\mathbb A}_S)}$ is as above.
\begin{example}\label{ex:4.11}
Let $G=SL(2),\;M=A,\;a=\left( 
\begin{array}{cc}
t & 0 \\ 
0 & t^{-1}
\end{array}
\right) $, where $t\neq \pm 1$. In this case, 
\[
r_A(\gamma ,a)=1,
\]
and 
\[
v_A(x)=\sum_{p\in S}\log (\sqrt{|1+n_p^2|_p}),
\]
where $x=\prod_{p\in S}a_p\cdot \left[ 
\begin{array}{cc}
1 & n_p \\ 
0 & 1
\end{array}
\right] \cdot k_p,$ for
$a_p\in \overline{A({\mathbb {Q}}_p)}$, $k_p\in \overline{K_{0,p}}$
(\cite{A6}, Intro., \cite{Sho}, chap. 6). Therefore,
$J_A(z,f)$, for $z\in Z(G({\mathbb Q}))$,
is the limit as $a\rightarrow 1$ of 
\begin{equation}\label{eq:4.21}
\ |D(a)|^{1/2}
\int_{\overline{A({\mathbb {A}}_S)}\backslash \overline{G({\mathbb {A}}_S)}}
f(y^{-1} az y)v_A(y)dy 
\end{equation}
\cite[section 6.1]{Sho}, p. 64.
\end{example}

\subsection{A more explicit form of the geometric terms}\label{ssec:explicit}

\begin{definition}\label{defn:4.12}
Let 
\[
K_{P,unip}(x,y)
=\sum_{\gamma \in {\mathcal U}_{M_P}{\mathbb {(Q)}}}
\int_{N_P({\mathbb {A})}}f(x^{-1}\gamma ny)dn,
\]
and for $T_0=0$, let 
\begin{eqnarray*}
J_{unip}(f)
&=&J_{unip}^{T_0}(f) \\
&=&\int_{G({\mathbb{Q}})\backslash \overline{G({\mathbb{A}})}}
\sum_P(-1)^{\dim(A_P/A_G)}
\sum_{\delta \in P({\mathbb {Q}})\backslash G({\mathbb{Q}})}
K_{P,unip}(\delta x,\delta x)
\tau _P^{\wedge }(H(\delta x)-T_0)dx.
\end{eqnarray*}
\end{definition}

\begin{example}\label{ex:4.13}
Let $G=SL(2),\;M=A$. Then 
\[
K_{P,unip}(x,y)
=\sum_{\gamma \in {\mathcal U}_{M_P}({\mathbb {Q}})}
\int_{N_P({\mathbb {A})}}f(x^{-1}\gamma ny)dn
\]
and 
\begin{eqnarray*}
J_{unip}(f)
=\int_{G({\mathbb {Q}})\backslash \overline{G({\mathbb{A}})}}
\sum_P(-1)^{\dim(A_P/A_G)}
\sum_{\delta \in P({\mathbb {Q}})\backslash G({\mathbb{Q}})}
K_{P,unip}(\delta x,\delta x)
\tau _P^{\wedge }(H(\delta x)-T_0)dx \\
=\int_{G({\mathbb{Q}})\backslash \overline{G({\mathbb{A}})}}
(K_{G,unip}(x,x)-
\sum_{\delta \in B({\mathbb {Q}})\backslash G({\mathbb{Q}})}
K_{B,unip}(\delta x,\delta x)
\tau _B^{\wedge }(H(\delta x)-T_0))dx
\end{eqnarray*}
\end{example}

\begin{definition}\label{defn:4.14}
Let $J_M(u,f)$ denote the weighted unipotent orbital integrals defined in
the previous subsection,
$u\in M({\mathbb A}_S)$ unipotent.
We define $a^M(S,u)$ by the property that 
\[
J_{unip}(f)
=\sum_P
\sum_{u\in {\mathcal U}_{M_P}({\mathbb {Q}})}
|W_0^M|\cdot |W_0^G|^{-1}a^M(S,u)J_M(u,f).
\]
(The verification that this definition makes sense can be found
for example in \cite[section 6.2, pp.~66--67]{Sho}.) For
$\gamma\in \overline{M({\mathbb A}_S)}$, we define $a^M(S,\gamma )$ by 
\[
a^M(S,\gamma )
=\varepsilon ^M(\sigma )|\iota ^M(\sigma)|^{-1}
\sum_{\{u\;|\;\sigma u\sim \gamma \}}a^{M_\sigma }(S,u),
\]
where 
\[
\varepsilon ^M(\sigma )=\left\{ 
\begin{array}{c}
1,\;\;{\rm if}\ \sigma \;{\rm is\ }k-{\rm elliptic\;on}\;M \\ 
0,\;\;\;{\rm otherwise}.\;\;\;\;\;\;\;\;
\end{array}
\right. 
\]
Furthermore, in the case of $G=SL(2)$, if $\gamma =\sigma u$
is the Jordan decomposition then
$a^M(S,\gamma)=\epsilon^M(\sigma)a^M(S,u)$.
\end{definition}

The relationship between the $J_{{\mathcal O}}(f)$ and the
$J_M(f)$ is given by the following:

\begin{proposition}\label{prop:4.15}
(\cite[Theorem 8.1]{A12}) If $S$ is a sufficiently large finite set of
valuations and
$f\in C_c^\infty (\overline{G({\mathbb {A}}_S)})_{gen}$ then 
\[
J_{{\mathcal O}}(f)=\sum_{M\in \{A,G\}}
|W_0^M|\cdot |W_0^G|^{-1}
\sum_{\gamma \in M({\mathbb {Q}})\cap {\mathcal{O}}}
a^M(S,\gamma )J_M(\gamma ,f) 
\]
where the constants $a^M(S,\gamma )$ defined above
satisfy, if $\gamma $ is semisimple, the property 
\[
a^M(S,\gamma )=\left\{ 
\begin{array}{c}
vol(M({\mathbb {Q}})_\gamma \backslash \overline{M({\mathbb {A}}_S)}_\gamma ^1),
\;\;\;\;{\rm if\;}\gamma \;{\rm is}\;
{\mathbb {Q}}-{\rm elliptic\;in}\;M({\mathbb {Q}}), \\ 
0,\;\;\;\;\;\;\;\;\;\;\;\;\;\;\;\;\;\;{\rm otherwise},\;\;\;\;\;
\end{array}
\right. 
\]
\end{proposition}

\begin{example}\label{ex:4.16} For $G=SL(2)$, $1\in {\mathcal O}$, we have
\[
J_{\mathcal O}(f)=\sum_{\gamma\in {\mathcal{O}}}
a^G(S,\gamma)J_G(\gamma,f)
+{1\over 2}a^A(S,1)J_A(1,f)
\]
where $J_A(1,f)$ is in Example~\ref{ex:4.13} above and $J_G(\gamma,f)$
is defined in \eqref{eq:4.19} above.
A similar result holds for the ${\mathcal{O}}$ with $-1\in {\mathcal{O}}$.
\end{example}



\section{The spectral side of the trace formula}\label{sec:spectral}

The spectral side of the trace formula comprises the collection of terms
associated with the representations of the group. This side, though
analogous to the already discussed geometric side of the trace formula in
terms of the overall structure, requires more background to state.
Essentially, we will have to leave the entire theory of Eisenstein series as
``known'' to the reader. Fortunately, there are several excellent
expositions of this theory in the case of $SL(2)$ and related groups:
\cite{L1}, \cite{K2}, \cite{KP} for example.

A general remark may be in order here. The spectral theory used in \cite{A2}
relies on the general theory of Eisenstein series on a quotient $\Gamma
\backslash H$ developed in \cite{L1}, where $H$ denotes a Lie group with finitely
many connected components and $\Gamma $ denotes a discrete subgroup
satisfying certain assumptions. By passing to the archimedian place, the
quotient $G({\mathbb{Q}})\backslash \overline{G({\mathbb{A}})^1}$ may be regarded as a
finite union of quotients $\Gamma _i\backslash H$ as in \cite{L1}
(see Lemma~\ref{lem:2.7} above). Thus we may apply the results of \cite{L1} here.

We begin with a brief overview of some of the background. Some results we
will eventually need are the decomposition \eqref{eq:5.4}, Proposition~\ref{prop:5.7}
and the definition following it, and Theorem~\ref{thm:5.8}. Other than \eqref{eq:5.4},
much of \S\ref{sssec:cuspidals} can be ignored or skimmed over as we sketch the
(slightly different) decompositions used in \cite{A1} and in \cite{A2,A3}.
We only need the decomposition (over a set ${\mathcal{X}}$ introduced below)
used in \cite{A2,A3}.

\subsection{The spectral decomposition}

A basic result in representation theory \cite[ch. 14]{Ro}, \cite[Lemma 3.2]{L2} is the
following

\begin{lemma}\label{lem:5.1}
Let $\Gamma \subset G$\ be a discrete subgroup of a locally compact
unimodular group $G$. The right regular representation of $G$
acting on $L^2(\Gamma \backslash G)$
decomposes as a (not necessarily discrete)
Hilbert sum of irreducible representations, each occuring with finite
multiplicity.
\end{lemma}

For $H<\overline{G({\mathbb{A}})}$ a locally
compact totally disconnected group of
type 1, define $\Pi_u(\overline{H})$ as in \S\ref{ssec:rank0lift} above.

Another way to look at the spectral decomposition
is adelically: The structure of the
unitary dual $\Pi_u(\overline{G({\mathbb{A}})})$ can be determined from the
structure of $\Pi_u(\overline{G({\mathbb{Q}}_p)})$.
(The proof of this is given in
\cite{Ge2} and is a consequence of a theorem of Gelbart
mentioned in Example~\ref{ex:5.2} below.) The {\bf
spectrum} of $G({\mathbb{Q}} )\backslash \overline{G({\mathbb{A}})}$ is the set of
elements of $\Pi_u(\overline{G({\mathbb{A}})})$
which occur in the representation
space $L^2(G({{\mathbb{Q}}} )\backslash \overline{G({\mathbb{A}})})$, on which
$\overline{G({\mathbb{A}})}$ acts on the right by translation.

\subsubsection{The partial decomposition in the adelic case: cuspidals}\label{sssec:cuspidals}

For each parabolic $P=MN$ of $G$ defined over ${\mathbb{Q}}$, define 
\begin{equation}\label{eq:5.1}
\begin{array}{c}
L^2(M({\mathbb{Q}})\backslash \overline{M({\mathbb{A}})^1})_{cusp} \\ 
=\{\phi \in L^2(M({{\mathbb{Q}}})\backslash \overline{M({\mathbb{A}})^1})
\ {\rm \ genuine\ }|\;{\rm for\;any\;parabolic}\;
P_1\subset P,\;P_1\neq P, \\ 
{\rm we\;have}
\int_{N_1({\mathbb{Q}})\cap M({\mathbb{Q}})\backslash
N_1({\mathbb{A}})\cap M({\mathbb{A}})}\phi (nm)dn=0, \\ 
{\rm for\;almost\;all}\;m\in \overline{M({\mathbb{A}})^1}\}.
\end{array}
\end{equation}
This is the space of genuine {\bf cusp forms} on
$\overline{M({\mathbb{A}})^1}$.

\begin{example}\label{ex:5.2}
(representations of adele groups) {\sf Case} $P=M=G=SL(2)$:
The genuine unitary representations $\pi $ of
$\overline{G({{\mathbb{A}}})}$ are all
``factorizable'', i.e.,
$\pi =\otimes _{p\leq \infty}\pi _p$ on
$\overline{G({{\mathbb{A}}})}$, where
$\pi _p$ is a genuine irreducible unitary representation of
$\overline{G({\mathbb{Q}}_p)}$ and almost all $\pi _p$
are unramified (i.e., has a non-trivial
$\overline{K_p}$-fixed vector). Conversely, given a collection 
\[
\{\pi _p\;|\;\pi _p\in \Pi _u(\overline{G({\mathbb{Q}}_p)}),
\;{\rm almost\;all}\;\pi _p\;{\rm are\;unramified}\},
\]
the restricted tensor product
$\pi =\otimes _{p\leq \infty}\pi _p$
defines an element of $\Pi _u(\overline{G({\mathbb{A}})})$. For a proof,
see for example, \cite[\S 3]{Ge2}.

\end{example}

\begin{proposition}\label{prop:5.3}
For each $\rho \in \Pi _u(\overline{M({\mathbb{A}})^1})$
there is an integer $m(\rho )\geq 0$ such that 
\[
L^2(M({{\mathbb{Q}}})\backslash \overline{M({\mathbb{A}})^1})_{cusp}\cong
\bigoplus_{\rho \in \Pi _u(\overline{M({\mathbb{A}})^1})}m(\rho )\rho .
\]
\end{proposition}

If
$\rho \in \Pi _u(\overline{M({\mathbb{A}})^1})$ and if $m(\rho )\neq 0$
then we call $\rho $ a genuine {\bf cuspidal
representation} of $\overline{M({\mathbb{A}})^1}$.

Two parabolics $P=MN$ and $P^{\prime }=M^{\prime }N^{\prime }$ are said to
be {\bf associate} parabolics if 
\[
\Omega ({\mathfrak{a}},{\mathfrak{a}}^{\prime })
=\{s\in \Omega \;|\;s{\mathfrak{a}}_{{\mathbb C}}
={\mathfrak{a}}_{{\mathbb C}}^{\prime }\}\neq \emptyset . 
\]
Here ${\mathfrak{a}}_{\mathbb{C}}$ denotes the
complexification of the Lie algebra
${\mathfrak{a}}$ of the center $A$ of $M$. In particular, a parabolic is
always associate to its opposite.
We identify the Weyl group $\Omega$ with a
subgroup of $G({\mathbb{Q}})$ (hence as a subgroup of
$\overline{G({\mathbb{A}})}$) by $s\longmapsto w_s\in G({\mathbb{Q}})$.
If $P,P^{\prime }$ are associate and if 
$\rho \in \Pi _u(\overline{M({\mathbb{A}})^1}),
\;\rho ^{\prime }\in \Pi _u(\overline{M^{\prime }({\mathbb{A}})^1})$
are cuspidal then we call $(M,\rho )$ {\bf equivalent} to
$(M^{\prime },\rho ^{\prime })$ if there exists an
$s\in \Omega ({\mathfrak{a}}, {\mathfrak{a}}^{\prime })$ such that
$s\rho $ is unitarily equivalent (as a
representation of $\overline{M^{\prime }({\mathbb{A}})^1})$
to $\rho ^{\prime }$, where 
\[
(s\rho )(m^{\prime })=\rho (w_s^{-1}m^{\prime }w_s),
\;\;\;\;m^{\prime }\in \overline{M^{\prime }({{\mathbb{A}}})^1}. 
\]
{\bf Notation}:
Let ${\mathcal{X}}={\mathcal{X}}^G$ denote the set of equivalent classes of such
pairs $(M,\rho )$.
For each $\chi \in {\mathcal X}$, we denote by
${\mathcal P}_\chi $ the corresponding class of associated parabolics.
Let ${\mathcal{X}}(G)$ denote the set of equivalent classes of
cuspidal pairs $(G,\rho )$.

\begin{example}\label{ex:5.4}
Two parabolics $P,P^{\prime }$ of $G=SL(2)$ with Levi $A$
are associate if and only they are opposite.
Therefore $(M,\rho )$ is equivalent to $(M^{\prime },\rho ^{\prime })$
if and only if $M=M^{\prime }=G$ and $\rho =\rho ^{\prime }$
(as equivalence classes of representations of
$\overline{G({\mathbb{A}})}$),
or $M=M'=A$ and $\rho ={\rho'}^w$, where 
$w=
\left(
\begin{array}{cc}
0 &-1\\
1&0
\end{array}
\right)$
denotes the non-trivial element of the Weyl group of $A$.
\end{example}

For each $P=MN$ and each $\chi \in {\mathcal X}$, we let 
\begin{equation}\label{eq:5.2}
L^2(M({{\mathbb{Q}}})\backslash \overline{M({\mathbb{A}})^1})_{cusp,\chi }
\cong
\bigoplus_{\rho \in \Pi _u(\overline{M({\mathbb{A}})^1}) ,
(M,\rho )\in \chi }
m(\rho )\rho .
\end{equation}
Note that if $P\notin {\mathcal P}_\chi $ then this sum is empty (which, by
convention, we define to be $0$).

The above theorem has the following

\begin{corollary}\label{cor:5.5}
$L^2(M({\mathbb{Q}})\backslash \overline{M({\mathbb{A}})^1})_{cusp}
=\bigoplus_{\chi \in {\mathcal X}}
L^2(M({\mathbb{Q}})\backslash \overline{M({\mathbb{A}})^1})_{cusp,\chi }$.
In particular, 
\[
L^2(G({\mathbb{Q}})\backslash \overline{G({\mathbb{A}})})_{cusp}
=\bigoplus_{\chi \in {\mathcal X}}
L^2(G({\mathbb{Q}})\backslash \overline{G({\mathbb{A}})})_{cusp,\chi }.
\]
\end{corollary}

In other words, the space of cuspidal automorphic functions decomposes into
``irreducible'' subspaces on which the action of $\overline{G({\mathbb{A}})}$ is the
more-or-less the same.

\subsubsection{The partial decomposition in the adelic case: theta functions}

We seek an analogous result for all of
$L^2(G({\mathbb{Q}})\backslash \overline{G({\mathbb{A}})})_{gen}$.

Let ${\mathfrak{Z}}_M$ denote the center of the
universal enveloping algebra ${\mathfrak{U}}_M$ of $M$. Each
$X\in {\mathfrak{Z}}_M$ defines a densely defined
unbounded operator on
$L^2(M({\mathbb{Q}})\backslash \overline{M({\mathbb{A}})^1})_{cusp}$.

Let ${\mathcal H}_P^0$ denote 
\begin{equation}\label{eq:5.3}
\begin{array}{c}
{\mathcal H}_P^0
=\{\varphi :N({\mathbb{A}})M({\mathbb{Q}})
\backslash \overline{G({\mathbb{A}})}\rightarrow {\mathbb{C}}
\ {\rm genuine}\;|\; \\ 
(i)\;\forall x\in \overline{G({\mathbb{A}})},
\;{\rm the\;function}\;m\longmapsto \varphi (mx),
\;m\in M({\mathbb{A}}), \\ 
{\rm is}\;{\mathfrak{Z}}_{M({\mathbb{R}})}{\rm -finite}, \\ 
(ii)\;{\rm the\;span\;of\;all\;functions}
\;\varphi _k:x\longmapsto \varphi (xk),\;k\in \overline{K}, \\ 
{\rm is\;finite\;}{\rm dimensional}, \\ 
(iii)\;||\varphi ||^2=
\int_{\overline{K}}
\int_{M({\mathbb{Q}})\backslash \overline{M({\mathbb{A}})^1}}
|\varphi (mk)|^2dmdk<\infty \} .
\end{array}
\end{equation}
Let ${\mathcal H}_P$ denote the closure of ${\mathcal H}_P^0$ with respect to this
norm in (iii).

For $\varphi \in {\mathcal H}_P$,
$\Lambda \in {\mathfrak{a}}_{{\mathbb C}},\;x,y\in \overline{G({\mathbb{A}})}$,
define $I_P(\Lambda ):\overline{G({\mathbb{A}})}\rightarrow End({\mathcal H}_P)$ by 
\[
(I_P(\Lambda ,y)\varphi )(x)=\varphi (xy)\exp (<\Lambda +\rho
_P,H_P(xy)>)\exp (-<\Lambda +\rho _P,H_P(x)>). 
\]
Here $\rho _P\in {\mathfrak{a}}$ is the element such that
$|\det (Ad(m))_{{\mathfrak{n}}({\mathbb A})}|
=\exp (<2\rho _P,H_P(m)>),\;m\in \overline{M({\mathbb{A}})}$.

For each $P=MN$ and each $\chi \in {\mathcal X}$, we define
$L^2(M({\mathbb{Q}})\backslash \overline{M({\mathbb{A}})^1})_\chi $
as follows. If there is no $P_1\in {\mathcal P}_\chi $
with $P_1\subset P$ then we define
$L^2(M({{\mathbb{Q}}})\backslash \overline{M({\mathbb{A}})^1})_\chi $
to be $0$. Otherwise, let $\Psi (P_1,P,\chi )$ denote the space of
all genuine complex-valued functions
$\psi :N({\mathbb{A}})M_1({\mathbb{Q}})\backslash \overline{G({\mathbb{A}})}
\rightarrow {\mathbb C}$, $P_1=N_1M_1\in {\mathcal{P}}_\chi $ 
with $P_1\subset P$, such that

\begin{enumerate}
\item  $\psi (amk)$, as a function of 
$a\in A_1({\mathbb{Q}})\backslash \overline{A_1({\mathbb{A}})}$ (fixed 
$m\in M_1({\mathbb{Q}})\backslash \overline{M_1({\mathbb{A}})^1},\;k\in K$),
is compactly supported,

\item  $\psi (amk)$, as a function of $k\in \overline{K_\infty}$,
belongs to the
vector space defined by the span of the matrix coefficients of some
irreducible representation $W$ of $\overline{K_\infty}$,

\item  $\psi (amk)$, as a function of $m\in M$, belongs to 
$L^2(M({{\mathbb{Q}}})\backslash \overline{M({\mathbb{A}})^1})_{cusp,\chi }$.
\end{enumerate}

Given $\psi \in \Psi (P_1,P,\chi )$, $P_1\subset P$, let 
\[
\psi ^{\wedge ,M}(m)=
\sum_{\delta \in P_1({\mathbb{Q}})\cap
M({\mathbb{Q}})\backslash M({\mathbb{Q}})}
\psi (\delta m),\;\;\;\;\;
m\in M({\mathbb{Q}})\backslash \overline{M({\mathbb{A}})^1}, 
\]
so that
$\psi ^{\wedge ,M}\in
L^2(M({\mathbb{Q}})\backslash \overline{M({\mathbb{A}})^1})$
\cite[Lemma 1]{L1}. Let
$L^2(M({\mathbb{Q}})\backslash \overline{M({\mathbb{A}})^1})_\chi $ denote
the $L^2$-closure of the span of all the
$\psi ^{\wedge ,M}$ with $\psi \in \Psi (P_1,P,\chi )$,
some $P_1\subset P$ as above. \cite[Lemma 2]{L1} implies 
\begin{equation}\label{eq:5.4}
L^2(M({\mathbb{Q}})\backslash \overline{M({\mathbb{A}})^1})
=\bigoplus_{\chi \in {\mathcal X}^G}
L^2(M({\mathbb{Q}})\backslash \overline{M({\mathbb{A}})^1})_\chi .
\end{equation}

\subsection{Eisenstein series in the adelic case}

For $\varphi \in {\mathcal H}_P^0,\;\lambda \in {\mathfrak{a}}_{{\mathbb C}},\;s\in
\Omega ({\mathfrak{a}},{\mathfrak{a}}_1),\;x\in \overline{G({\mathbb{A}})}$, let
$M(s,\lambda)$ denote the intertwining integral
\begin{equation}\label{eq:5.5}
\begin{array}{c}
M(s,\lambda )\phi(x) \\ 
=\exp (-<s\lambda +\rho _1,H_P(x)>)\times \;\;\;\;\;\;\;\;\;\;\;\;\; \\ 
\;\;\;\;\;\;\times
\int_{N_1({\mathbb{A}})\cap w_sN({\mathbb{A}})w_s^{-1}\backslash
N_1({\mathbb{A}})}\varphi (w_s^{-1}nx)
\exp (<\lambda +\rho_P,H_P(w_s^{-1}nx)>)dn,
\end{array}
\end{equation}
with a suitably normalized measure $dn$.
We define $M(s,\pi_\zeta)$ below (see \eqref{eq:5.9}).

Let ${\mathfrak{a}}^{+}$ denote the positive Weyl chamber ${\mathfrak{a}}^{+}=\{X\in 
{\mathfrak{a}}\;|\;\alpha (X)>0,\;\forall \alpha >0\}$.
For $\varphi \in {\mathcal{H}}_P^0,
\;\lambda \in {\mathfrak{a}}_{{\mathbb C}},
\;Re(\lambda )\in \rho _P+{\mathfrak{a}}^{+},
\;x\in \overline{G({\mathbb{A}})}$, let 
\begin{equation}\label{eq:5.6}
E(x,\varphi ,\Lambda )=
\sum_{\delta \in P({\mathbb{Q}})\backslash G({\mathbb{Q}})}
\varphi (\delta x)\exp (<\lambda +\rho _P,H_P(\delta x)>).
\end{equation}
This is the {\bf Eisenstein series} associated to $P$ and $\varphi $ and is
known to converge absolutely. 

Let $P=MN$ denote a parabolic and let
$\pi \in \Pi _u(\overline{M({\mathbb{A}})^1})$.
Following \cite{A2}, we define 
\begin{equation}\label{eq:5.7}
\begin{array}{c}
{\mathcal H}_P^0(\pi )=\{\varphi :N({\mathbb{A}})M({\mathbb{Q}})
\backslash \overline{G({\mathbb{A}})}\rightarrow {\mathbb C}\;
{\rm genuine\ }|\; \\ 
(i)\;\phi \;{\rm is\;right}\;K{\rm -finite,} \\ 
(ii)\;\forall x\in \overline{G({\mathbb{A}})},
\;{\rm the\;function}\;m\longmapsto \varphi (mx)
\;(m\in \overline{M({{\mathbb{A}}})^1}) \\ 
{\rm is\;a\;matrix\;coefficient\;of}\;\pi , \\ 
(iii)\;||\varphi ||^2
=\int_K\int_{M({\mathbb{Q}})\backslash \overline{M({\mathbb{A}})^1}}|
\varphi (mk)|^2dmdk<\infty \}.
\end{array}
\end{equation}

Let ${\mathcal H}_P(\pi )$ denote the completion of ${\mathcal H}_P^0(\pi )$ with
respect to the norm $||...||$ in (iii). For
$\phi \in {\mathcal H}_P(\pi )$, $\varsigma \in i{\mathfrak{a}}^{*},$ let
\begin{eqnarray*}
\phi _\varsigma (x)
&=&\phi (x)e^{\varsigma (H_P(x))},\;\;\;x\in \overline{G({\mathbb{A}})},
\\
\pi _\varsigma  &=&\pi \otimes \chi _\varsigma ,\;\;\;\chi _\varsigma
(x)=e^{\varsigma (H_P(x))},\;\;x\in \overline{G({\mathbb{A}})}, \\
(I_P(\pi _\varsigma ,y)\phi _\varsigma )(x) &=&\phi _\varsigma (xy)\delta
_P(xy)^{1/2}\delta _P(x)^{-1/2},\;\;\;x,y\in \overline{G({\mathbb{A}})}.
\end{eqnarray*}

Let $P=MN$ denote a parabolic, let $\pi \in \Pi _u(\overline{M({\mathbb{A}})})$,
let $f$ be a function on $\overline{M({\mathbb{A}})}$,
and let $\zeta\in {\mathfrak{a}}^*_{\mathbb{C}}$. Define $\pi_\zeta$ by
\begin{equation}\label{eq:5.8}
\begin{array}{c}
f_\zeta (m)=f(m)e^{\zeta(H_M(m))},\\
\pi_\zeta(m)=\pi(m)e^{\zeta(H_M(m))},\ \ \ m\in \overline{M({\mathbb{A}})}.
\end{array}
\end{equation}

The space of (genuine) matrix coefficients of $\pi$, which we denote by
${\mathcal A}(\pi)$, is isomorphic to that of $\pi_\zeta$, for
$\zeta\in i{\mathfrak{a}}^*$:
$$
{\mathcal A}(\pi)\cong {\mathcal A}(\pi_\zeta).
$$
(proof: We may let $\pi,\ \pi_\zeta$ have the same representation space
$V$ with inner product $(...,...)_V$. If $v_1,v_2\in V$ and
$f(g)=(\pi(g)v_1,v_2)_V$ then $f_\zeta\in {\mathcal A}(\pi_\zeta)$ and every
element of ${\mathcal A}(\pi_\zeta)$ arises in this way in a unique
fashion. $\Box$) We have a natural isomorphism
$$
\begin{array}{c}
{\mathcal H}_P(\pi)\rightarrow {\mathcal H}_P(\pi_\zeta),\\
\phi\longmapsto \phi_\zeta,
\end{array}
$$
For $\pi \in \Pi _u(\overline{M({\mathbb{A}})})$,
$\phi \in {\mathcal H}_P^0(\pi )$,
$\varsigma \in {\mathfrak{a}}_{P,{\mathbb C}}^{*}$,
$s\in \Omega ({\mathfrak{a}},{\mathfrak{a}}^{\prime})$,
$Re(\varsigma )\in \rho _P+({\mathfrak{a}}_P^{*})^{+}$, define
\begin{equation}\label{eq:5.9}
\begin{array}{c}
E(x,\phi _\varsigma )
=\sum_{\delta \in P({\mathbb{Q}})\backslash G({\mathbb{Q}})}
\phi _\varsigma (\delta x)\delta _P(\delta x)^{1/2}, \\
(M(s,\pi _\varsigma )\phi _\varsigma )(x)
=\delta_P(x)^{-1/2}
\int_{N^{\prime }({\mathbb{A}})\cap w_sN({\mathbb{A}})w_s^{-1}\backslash
N^{\prime }({\mathbb{A}})}
\phi _\varsigma (w_s^{-1}nx)\delta_P(w_s^{-1}nx)^{1/2}dn.
\end{array}
\end{equation}
Analogously, if $P\subset P_2$, $\varsigma \in {\mathfrak{a}}_{{\mathbb C}}^{*}$ is
sufficiently large then define
\[
E_{P_2}(x,\phi _\varsigma )
=\sum_{\delta \in P({\mathbb{Q}})\backslash P_2({\mathbb{Q}})}
\phi _\varsigma (\delta x)\delta _P(\delta x)^{1/2}.
\]

Let
\[
\begin{array}{c}
{\mathcal H}_P(\pi )_\chi
=\{\phi \in {\mathcal H}_P(\pi )\;|\;m\longmapsto \varphi (mx)
\;(m\in M({\mathbb{Q}})\backslash \overline{M({\mathbb{A}})}^1) \\
{\rm \;\;\;\;\;\;\;\;\;\;\;\;belongs\;to}
\;L^2(M({\mathbb{Q}})\backslash \overline{M({\mathbb{A}})^1})_\chi \}.
\end{array}
\]

\begin{proposition}\label{prop:5.6}
${\mathcal H}_P(\pi )=\bigoplus_{\chi \in {\mathcal X}}{\mathcal H}_P(\pi )_\chi $.
\end{proposition}

Let $K\subset G({\mathbb{A}}_f)\cap K_0$ be an open compact subgroup and let 
\begin{equation}\label{eq:5.10}
{\mathcal H}_P(\pi )_{\chi ,\overline{K}}\subset {\mathcal H}_P(\pi )_\chi 
\end{equation}
denote the subspace which is right invariant under
$\overline{K}$. Let $W$ denote a genuine
irreducible representation of $\overline{K_\infty} $ and let 
\begin{equation}\label{eq:5.11}
{\mathcal H}_P(\pi )_{\chi ,\overline{K},W}
\subset {\mathcal H}_P(\pi )_{\chi ,\overline{K}}
\end{equation}
denote the subspace of functions $\phi $ such that
$k\longmapsto \phi (xk)\;\;(k\in \overline{K_\infty} )$
is a matrix coefficient of $W$.

\begin{proposition}\label{prop:5.7}
$\dim {\mathcal H}_P(\pi )_{\chi ,\overline{K},W}<\infty .$
\end{proposition}

We apply this to fixing a choice of a basis:
For each connected component $\Pi^0\subset
\Pi_u(\overline{M({\mathbb{A}})}^1)$, fix an element $\pi^0\in \Pi^0$.
Every other $\pi\in\Pi^0$ may be uniquely written as
$\pi=\pi^0_\zeta$, for some $\zeta\in i{\mathfrak{a}}^*_P$.
Let ${\mathcal B}_P(\pi^0 )_\chi $
denote a basis of ${\mathcal H}_P(\pi^0 )_\chi $ satisfying the following
properties:

\begin{enumerate}
\item  each $\phi \in {\mathcal B}_P(\pi^0 )_\chi $ belongs to
${\mathcal H}_P(\pi^0)_{\chi ,\overline{K},W}$
for some $K,W$ (possibly depending on $\phi $),

\item  the elements of ${\mathcal B}_P(\pi^0 )_\chi $
form an orthonormal set with
respect to the inner product
\[
(\phi _1,\phi _2)=
\int_K\int_{M({\mathbb{Q}})\backslash \overline{M({\mathbb{A}})}^1}
\phi _1(mk)\overline{\phi _2(mk)}dmdk,
\]

\item  if $\varsigma \in i{\mathfrak{a}}^{*}$ then 
\[
{\mathcal B}_P(\pi^0_\varsigma )_\chi =\{\phi _\varsigma \;|\;
\phi \in {\mathcal B}_P(\pi^0 )_\chi \}.
\]
\end{enumerate}

${\mathcal B}_P(\pi)_\chi$ is a basis for
${\mathcal H}_P(\pi)_\chi$, if $\pi=\pi^0$, and is a subset of
${\mathcal H}_P(\pi)$ in general.

Let ${\mathcal P}$ denote a class of associated parabolics. Let
$L_{{\mathcal P}}^{\wedge }$ denote the
set of collections $F=\{F_P\;|\;P\in {\mathcal P}\}$ of
measurable functions
$F_P:\Pi _u(\overline{M({\mathbb{A}})})\rightarrow {\mathcal H}_P$ such
that

(i) if $s\in \Omega ({\mathfrak{a}},{\mathfrak{a}}_1)$ then 
\[
F_{P_1}(s\pi )=M(s,\pi )F_P(\pi ),
\]

(ii) if $N(A)$ denotes the number of chambers in ${\mathfrak{a}}$
(which is equal to 1 for $SL(2)$) then 
\[
||F||^2=\sum_{P\in {\mathcal{P}}}{1\over 2\pi }^{dim(A)}
\int_{\Pi _u(\overline{M({\mathbb{A}})})}
||F_P(\pi )||^2d\pi <\infty .
\]

The following result describes some basic properties of Eisenstein series.

\begin{theorem}\label{thm:5.8} For
$s\in \Omega ({\mathfrak{a}},{\mathfrak{a}}_1)$, we have
{\bf \ }(Langlands \cite[p. 927]{A2})

(a)\ If $\varphi \in {\mathcal H}_P^0(\pi )$\ then
$E(x,\varphi _\varsigma )$ and
$M(s,\pi _\varsigma )\varphi _\varsigma $
can be analytically continued
to a meromorphic function on ${\mathfrak{a}}_{{\mathbb C}}^{*}$. On 
$i{\mathfrak{a}}$ the
Eisenstein series $E(x,\varphi_\varsigma )$
is regular and the intertwining operator $M(s,\pi_\varsigma )$
is unitary. For $\overline{K_0}$-finite
$f\in C_c^\infty (\overline{G({\mathbb{A}})})_{gen}$,
$t\in \Omega ({\mathfrak{a}}_1,{\mathfrak{a}}_2)$, we have 
\[
\begin{array}{c}
(i)\;\;\;\;E(x,I_P(\pi ,f)\varphi )
=\int_{\overline{G({\mathbb{A}})}}f(y)E(xy,\varphi )dy,
\\ 
(ii)\;\;\;\;\;E(x,M(s,\pi )\varphi )=E(x,\varphi ), \\ 
(iii)\;\;\;\;M(ts,\pi )=M(t,s\pi )M(s,\pi ).
\end{array}
\]
(b) There is a dense subspace of $L_{{\mathcal P}}^{\wedge }$\ on which the map 
\[
F\longmapsto \sum_{P\in {\mathcal P}}n(A)^{-1}
\left( \frac 1{2\pi i}\right)^{\dim (A)}
\int_{\Pi _u(\overline{M({\mathbb{A}})})}E(x,F_P(\pi ))d\pi 
\]
is defined. This extends to all of $L_{{\mathcal P}}^{\wedge }$\ and defines a
unitary map of $L_{{\mathcal P}}^{\wedge }$\ onto a closed subspace 
\[
L_{{\mathcal P}}^2(G({{\mathbb{Q}}})\backslash \overline{G({\mathbb{A}})})
\subset L^2(G({\mathbb{Q}})\backslash \overline{G({\mathbb{A}})})_{gen}.
\]
Moreover, 
\[
L^2(G({\mathbb{Q}})\backslash \overline{G({\mathbb{A}})})_{gen}
=\bigoplus_{\mathcal{P}}
L_{\mathcal{P}}^2(G({\mathbb{Q}})\backslash \overline{G({\mathbb{A}})}).
\]
\end{theorem}



\subsection{The integrability of $k_\chi ^T(x,f)$}

For $K$-finite compactly supported $f$ as above, let 
\begin{equation}\label{eq:5.12}
K_{P,\chi }(x,y)=\sum_{P_1\subset P}n_P(A_1)^{-1}
\int_{\Pi _u(\overline{M(\aaa )})}
\sum_{\phi \in {\mathcal{B}}_{P_1}(\pi )_\chi }
E_P(x,I_{P_1}(\pi ,f)\phi )\overline{E_P(y,\phi )}d\pi
\end{equation}
and 
\begin{equation}\label{eq:5.13}
k_\chi ^T(x,f)=\sum_P(-1)^{\dim (A_P/Z)}
\sum_{\delta \in P(\qqq )\backslash G(\qqq )}
K_{P,\chi }(x,x)\tau _P^{\wedge }(H(\delta x)-T).
\end{equation}

\begin{lemma}\label{lem:5.9}
$K_{P,\chi }(x,y)$ is the kernel function for the regular 
representation $R_P(f)$ of $\overline{G(\aaa )}$ on 
\[
L^2(G(\qqq )\backslash \overline{G(\aaa )})_{\mathcal{P},\chi }
=L^2(G(\qqq )\backslash \overline{G(\aaa )})_{\mathcal{P}}
\cap L^2(G(\qqq )\backslash \overline{G(\aaa )})_{\chi},
\]
where $\mathcal{P}$ denotes the associate class of $P$. In particular,
$K_{G,\chi }(x,y)=K_\chi (x,y)$ is the kernel function for the regular
representation $R_P(f)$ of $\overline{G(\aaa )}$ on
$L^2(G(\qqq )\backslash \overline{G(\aaa )})_\chi $.
\end{lemma}

\begin{remark}\label{rem:5.10} In \cite{A15}, different notation is used: $R_P(f)$ denotes
the convolution operator associated to right translation on
\[
L^2(N_P(\aaa )M_P(\qqq )\backslash \overline{G(\aaa )}))
=\oplus_\chi
L^2(N_P(\aaa )M_P(\qqq )\backslash \overline{G(\aaa )}))_\chi .
\]
\end{remark}

\begin{theorem}\label{thm:5.11}
$\sum_{\chi \in \mathcal{X}}K_{P,\chi }(x,y)
=\sum_{{\mathcal O}\in \mathcal{C}}
K_{P,{\mathcal O}}(x,y)$.
\end{theorem}

proof: Both sides represent the kernel function for the regular
representation $R_P(f)$ of $\overline{G(\aaa )}$
on $L^2(G(\qqq )\backslash \overline{G(\aaa )})$. $\Box $

The trace formula is obtained by integrating both sides of 
\[
\sum_{\chi \in \mathcal{X}}\Lambda ^TK_{P,\chi }(x,x)
=\sum_{{\mathcal O}\in \mathcal{C}}
\Lambda ^TK_{P,{\mathcal O}}(x,x) 
\]
over $G(\qqq )\backslash \overline{G(\aaa )}$,
where $\Lambda ^T$ is the
truncation operator defined below. We know from earlier sections that the
right-hand side is a polynomial in $T$, so the left-hand side is too. The 
\textbf{global trace formula} is the constant term of this expression. 

To prove the integrability of $k_\chi ^T(x,f)$ a truncation operator is
introduced.
If $\phi \in C(G(\qqq )\backslash \overline{G(\aaa )})_{gen}$
then define,
as in \cite[p 89]{A3}, the \textbf{truncation operator} by 
\begin{equation}\label{eq:5.14}
\Lambda ^T\phi (x)
=\sum_P(-1)^{\dim (A/Z)}
\sum_{\delta \in P(\qqq )\backslash G(\qqq )}
\int_{N(\qqq )\backslash N(\aaa )}\phi (n\delta
x)dn\cdot \tau _P^{\wedge }(H(\delta x)-T).
\end{equation}
More generally, if $P_1\subset G$ is a parabolic and if
$\phi \in C(P_1(\qqq )\backslash \overline{G(\aaa )})_{gen}$ then
define, as in \cite[p 97]{A3}, the \textbf{truncation operator} by 
\begin{equation}\label{eq:5.15}
\Lambda ^{T,P_1}\phi (x)
=\sum_{R , P_0\subset R\subset P_1}
(-1)^{\dim (A_R/A_1)}
\sum_{\delta \in R(\qqq )\backslash P_1(\qqq )}
\int_{N_R(\qqq )\backslash N_R(\aaa )}
\phi (n\delta x)dn\cdot \tau_R^{\wedge ,1}(H(\delta x)-T).
\end{equation}
In the classical case of $SL(2)$, the truncation operator acting
on the Eisenstein series has the effect of subtracting
off the ``$y^s$ term'' in its Fourier expansion. This allows one 
to compute the inner product of two truncated Eisenstein series
(see Theorem~\ref{thm:5.27} below).

\begin{example}\label{ex:5.12}
For $G=SL(2),P_1=G$ this formula becomes 
\begin{equation}\label{eq:5.16}
\Lambda ^T\phi (x)=\phi (x)-\sum_{\delta \in B(\qqq )\backslash G(\qqq %
)}\int_{N(\qqq )\backslash N(\aaa )}\phi (n\delta x)dn\cdot \tau
_B^{\wedge }(H(\delta x)-T).
\end{equation}
\end{example}

\begin{lemma}\label{lem:5.13}
The truncation operator is a self-adjoint idempotent on
$C(G(\qqq )\backslash \overline{G(\aaa )})_{gen}
\cap L^2(G(\qqq )\backslash \overline{G(\aaa )})$ .
\end{lemma}
proof: This is an immediate consequence of the non-metaplectic
version \cite{A3} since the metaplectic cover splits over $N$. $\Box$

\begin{lemma}\label{lem:5.14}
(\cite[Lemma 1.5]{A3}) For
$\phi \in C(P(\qqq )\backslash \overline{G(\aaa )})_{gen}$, we
have 
\[
\sum_{ P_1 , P_0\subset P_1\subset P }
\sum_{\delta \in P_1(\qqq )\backslash P(\qqq )}
\Lambda ^{T,P_1}\phi (\delta x)\tau _1^P(H_0(\delta x)-T)
=\int_{N(\qqq )\backslash N(\aaa )}\phi (nx)dn.
\]
\end{lemma}

proof in the case $G=SL(2)$: Suppose $P=G$. First, from the fact that 
\begin{equation}\label{eq:5.17}
\sum_{ P_1 , R\subset P_1\subset Q}
 (-1)^{\dim (A_R/A_1)}\tau_R^{\wedge ,1}(X)\tau _1^Q(X)
=\left\{ 
\begin{array}{c}
1,\;\;\;\;Q=R, \\ 
0,\;\;\;\;Q\neq R,
\end{array}
\right.
\end{equation}
(this is Langlands' combinatorial lemma; see \cite{A3})
% Claude (2026): the original cited a nonexistent "lemma 6.9".
, we have $\tau
_B^{\wedge ,B}(X)\tau _B^G(X)=\tau _B^{\wedge ,G}(X)\tau _G^G(X)=\tau
_B^{\wedge }(X)$. The LHS of the equation in the above lemma is 
\begin{eqnarray*}
&&\sum_{ P_1, B\subset P_1\subset G } 
\sum_{\delta \in P_1(\qqq )\backslash G(\qqq )}
\Lambda ^{T,P_1}\phi (\delta x)\tau _1^G(H_0(\delta x)-T) \\
&=&\Lambda ^{T,G}\phi (x)\tau _G^G(H_0(x)-T)+ \\
&&+\sum_{\delta \in B(\qqq )\backslash G(\qqq )}\Lambda ^{T,B}\phi
(\delta x)\tau _B^G(H_0(\delta x)-T) \\
&=&\Lambda ^{T,G}\phi (x)+\sum_{\delta \in B(\qqq )\backslash G(\qqq %
)}\int_{N(\qqq )\backslash N(\aaa )}\phi (n\delta x)dn\cdot \tau
_B^{\wedge ,B}(H_0(\delta x)-T)\cdot \tau _B^G(H_0(\delta x)-T) \\
&=&\Lambda ^{T,G}\phi (x)+\sum_{\delta \in B(\qqq )\backslash G(\qqq %
)}\int_{N(\qqq )\backslash N(\aaa )}\phi (n\delta x)dn\cdot \tau
_B^{\wedge }(H_0(\delta x)-T) \\
&=&\phi (x),
\end{eqnarray*}
by the previous example.

Suppose $P=B$. The LHS of the equation in the above lemma is 
\[
\Lambda ^{T,B}\phi (x)\tau _B^B(H_0(x)-T)
=\int_{N(\qqq )\backslash N(\mathbb{A})}\phi (nx)dn, 
\]
since $\tau _B^B(X)=1$.

This verifies the lemma in both cases. $\Box $

\begin{lemma}\label{lem:5.15}
(\cite[Lemma 2.2]{A3})

\begin{eqnarray*}
k_\chi ^T(x,f)=\sum_{ P_1,P_2 {\rm s.t.} P_0\subset P_1\subset P_2}
\sum_{\delta \in P_1(\qqq )\backslash G(\qqq )}
\sigma _1^2(H_0(\delta x)-T)\times  \\
\times \left\{ \sum_{ P, P_1\subset P\subset P_2 }
(-1)^{\dim (A/Z)}\Lambda _2^{T,P}K_{P,\chi }(\delta x,\delta x)\right\} .
\end{eqnarray*}
(The subscript $2$ in $\Lambda _2^{T,P}$ means that it is to be applied to
the second variable.)
\end{lemma}

\begin{example}\label{ex:5.16}
Let $G=SL(2)$. The above lemma becomes 
\begin{eqnarray*}
k_\chi ^T(x,f)=
\sum_{ P_1,P_2 {\rm s.t.} (P_1,P_2)\in  \{(B,B),(B,G),(G,G)\}} 
\sum_{\delta \in P_1(\qqq )\backslash G(\qqq )}
\sigma_1^2(H_0(\delta x)-T)\times  \\
\ \times \left\{ \sum_{ P , P_1\subset P\subset P_2 }
(-1)^{\dim (A/Z)}\Lambda _2^{T,P_1}K_{P,\chi }(\delta x,\delta x)\right\}  \\
=-\sum_{\delta \in B(\qqq )\backslash G(\qqq )}\sigma _B^B(H_0(\delta
x)-T)\Lambda _2^{T,B}K_{B,\chi }(\delta x,\delta x)\;\;(\mathrm{the}\;(B,B)\;%
\mathrm{term}) \\
+\sigma _G^G(H_0(x)-T)\Lambda _2^{T,G}K_{G,\chi }(x,x)\;\;(\mathrm{the}%
\;(G,G)\;\mathrm{term}) \\
+\sum_{\delta \in B(\qqq )\backslash G(\qqq )}\sigma _B^G(H_0(\delta
x)-T)\times  \\
\;\;\times \left\{ -\Lambda _2^{T,B}K_{B,\chi }(\delta x,\delta x)+\Lambda
_2^{T,B}K_{G,\chi }(\delta x,\delta x)\right\} \;\;\;\;\;\;\;(\mathrm{the}%
\;(B,G)\;\mathrm{term}) \\
=\Lambda _2^{T,G}K_{G,\chi }(x,x)+ \\
+\sum_{\delta \in B(\qqq )\backslash G(\qqq )}\tau _B(H_0(\delta
x)-T)\left\{ -\Lambda _2^{T,B}K_{B,\chi }(\delta x,\delta x)+\Lambda
_2^{T,B}K_{G,\chi }(\delta x,\delta x)\right\}  \\
=\Lambda _2^{T,G}K_{G,\chi }(x,x)+\sum_{\delta \in B(\qqq )\backslash G(%
\qqq )}\tau _B(H_0(\delta x)-T)\Lambda _2^{T,B}K_{G,\chi }(\delta x,\delta
x).
\end{eqnarray*}
by \eqref{eq:2.5}.
% Claude (2026): the original cited "lemma 4.7" here.
The definition is 
\begin{equation}\label{eq:5.18}
k_\chi ^T(x,f)=K_{G,\chi }(x,x)-\sum_{\delta \in B(\qqq )\backslash G(\mathbb{%
Q})}K_{B,\chi }(\delta x,\delta x)\tau _B^{\wedge }(H_0(\delta x)-T).
\end{equation}

We claim now that 
\begin{equation}\label{eq:5.19}
\Lambda _2^{T,G}K_{G,\chi }(x,x)=K_{G,\chi }(x,x)-\sum_{\delta \in B(\qqq %
)\backslash G(\qqq )}\tau _B(H(\delta x)-T)\Lambda _2^{T,B}K_{G,\chi
}(\delta x,\delta x).
\end{equation}
To see this, note 
\[
\Lambda _2^{T,G}K_{G,\chi }(x,x)=K_{G,\chi }(x,x)-\sum_{\delta \in B(\qqq %
)\backslash G(\qqq )}\int_{N(\qqq )\backslash N(\aaa )}\Lambda
_2^{T,B}K_{G,\chi }(x,n\delta x)dn\cdot \tau _B^{\wedge }(H(\delta x)-T).
\]
On the other hand, 
\[
\Lambda _2^{T,B}K_{G,\chi }(x,x)=\int_{N(\qqq )\backslash N(\aaa %
)}K_{G,\chi }(x,nx)dn\cdot \tau _B^{\wedge ,B}(H(\delta x)-T),
\]
so 
\begin{eqnarray}
\sum_{\delta \in B(\qqq )\backslash G(\qqq )}\int_{N(\qqq )\backslash
N(\aaa )}\Lambda _2^{T,B}K_{G,\chi }(x,n\delta x)dn\cdot \tau _B^{\wedge
}(H(\delta x)-T)  \nonumber \\
=\sum_{\delta \in B(\qqq )\backslash G(\qqq )}\Lambda _2^{T,B}K_{G,\chi
}(x,\delta x)\cdot \tau _B^{\wedge }(H(\delta x)-T)  \nonumber \\
=\sum_{\delta \in B(\qqq )\backslash G(\qqq )}\Lambda _2^{T,B}K_{G,\chi
}(\delta x,\delta x)\cdot \tau _B(H(\delta x)-T),
\label{eq:5.20}\end{eqnarray}
since the volume of $N(\qqq )\backslash N(\aaa )$ and $\tau _B(X)=\tau
_B^{\wedge }(X)$ in this case. The claimed result follows.
\end{example}

\begin{example}\label{ex:5.17}
Let $G=SL(2),\;P=G$. Then 
\[
E_G(x,\phi )=E(x,\phi )=\sum_{\delta \in B(\qqq )\backslash G(\qqq %
)}\phi (\delta x)\delta _B(\delta x)^{1/2}.
\]

If $P=B$ then $E_B(x,\phi )=\phi (x)\delta _B(x)^{1/2}$.
\end{example}

\begin{example}\label{ex:5.18}
Let $G=SL(2),\;P=B$. For $\pi \in \Pi _u(\overline{A(\aaa )})$,
$\chi \in \mathcal{X}$,
the set ${\mathcal{B}}_B(\pi )_\chi $ is an orthonormal basis for
${\mathcal{H}}_B(\pi )_\chi $. We have 
\[
{\mathcal{H}}_B(\pi )_\chi =\left\{
\begin{array}{c}
\phi:N(\aaa )A(\qqq )\setminus \overline{G(\aaa )}\rightarrow
{\mathbb C}\ {\rm genuine\ }|\ \\
(1)\ \  \phi {\rm\ is\  right}\  \overline{K_0}\ {\rm -finite},\\
(2) \ \ \phi(ag)=\pi(a)\phi(g),
\ \ a\in \overline{A(\aaa )},\ g\in \overline{G(\aaa )},\\
(3)\ \  ||\phi ||<\infty
\end{array}
\right\}.
\]
\end{example}

\begin{example}\label{ex:5.19}
Let $G=SL(2),\;P=G$. For $\pi \in \Pi _u(\overline{G(\aaa )})$
cuspidal,$\;\chi \in 
\mathcal{X}$, the set ${\mathcal{B}}_G(\pi )_\chi $ is an orthonormal basis
for ${\mathcal{H}}_G(\pi )_\chi $. If $\chi = (G,\pi )$
(in $\mathcal{X}$) then this space is non-zero and consists of
(measurable equivalence classes of) square-integrable
genuine functions on $G(\qqq )\backslash \overline{G(\aaa )}$
which are cuspidal automorphic forms of type $\pi_\infty $.
Since almost every $\pi _p$ is unramified, the action of almost
all $\overline{K_{0,p}}$ is trivial on the
functions in this space. Since the functions in
this space are all $\overline{K_0}$-finite,
we may decompose it by the action of $\overline{K_0}$ on
the right: 
\[
{\mathcal{H}}_G(\pi )_\chi =
\bigoplus_{\sigma \in \Pi(\overline{K_0})}
{\mathcal{H}}_G(\pi )_{\chi ,\sigma }
\]
where each irreducible representation
$\sigma \in \Pi(\overline{K_0})$ occurs at
most with finite multiplicity. For the basis above, one must pick a element
in each one of these ``$\overline{K_0}$-types''.
\end{example}

\begin{example}\label{ex:5.20}
Let $G=SL(2),\;P=B$. Recalling the determination of
${\mathcal{B}}_B(\pi)_\chi $ above, we have 
\[
K_{B,\chi }(x,y)
=\int_{\Pi _u(\overline{A(\aaa )})}
\sum_{\phi \in {\mathcal{B}}_B(\pi)_\chi }
I_B(\pi ,f)\phi (x)d\pi \cdot \delta _B(x)\overline{\phi (y)}
\]
\end{example}

\begin{example}\label{ex:5.21}
Let $G=SL(2),\;P=G$. Recalling the determination
of ${\mathcal{B}}_B(\pi)_\chi $ above, we have 
\begin{eqnarray*}
K_{G,\chi }(x,y)
=\int_{\Pi _u(\overline{G(\aaa )})}
\sum_{\phi \in {\mathcal{B}}_G(\pi)_\chi }
E(x,I_G(\pi ,f)\phi )d\pi \cdot \overline{E(y,\phi )}+ \\
\int_{\Pi _u(\overline{A(\aaa )})}
\sum_{\phi \in {\mathcal{B}}_B(\pi )_\chi }
I_B(\pi,f)\phi (x)d\pi \cdot \delta _B(x)\overline{\phi (y)}.
\end{eqnarray*}
\end{example}

\begin{example}\label{ex:5.22}
Let $G=SL(2)$. We have 
\begin{eqnarray*}
k_\chi ^T(x,f)
=K_{G,\chi }(x,x)-\sum_{\delta \in B(\qqq )\backslash G(\mathbb{Q})}
K_{B,\chi }(\delta x,\delta x)
\tau _B^{\wedge }(H_0(\delta x)-T) \\
=\int_{\Pi _u(\overline{G(\aaa )})}
\sum_{\phi \in {\mathcal{B}}_G(\pi )_\chi}
E_G(x,I_G(\pi ,f)\phi )d\pi \cdot \overline{E_G(x,\phi )}+ \\
+\int_{\Pi _u(\overline{A(\aaa )})}
\sum_{\phi \in {\mathcal{B}}_B(\pi )_\chi }
I_B(\pi,f)\phi (\delta x)d\pi \cdot \delta _B(\delta x)
\overline{\phi (\delta x)}(1-\tau _B^{\wedge }(H(x)-T)) \\
-\sum_{ \delta \in B(\qqq )\backslash G(\qqq ), \delta \neq 1 }
\int_{\Pi _u(\overline{A(\aaa )})}
\sum_{\phi \in {\mathcal{B}}_B(\pi )_\chi }I_B(\pi,f)
\phi (\delta x)d\pi \cdot \delta _B(\delta x)
\overline{\phi (\delta x)}\tau _B^{\wedge }(H(x)-T)
\end{eqnarray*}
\end{example}

Let 
\begin{equation}\label{eq:5.21}
K_\chi ^T(f)=\int_{G(\qqq )\backslash \overline{G(\aaa )}}
\Lambda _2^TK_\chi (x,x)dx.
\end{equation}
Define 
\begin{equation}\label{eq:5.22}
J_\chi ^T(x,y)=\sum_Pn(A)^{-1}
\int_{\Pi _u(\overline{M(\aaa )})}
\sum_{\phi \in {\mathcal{B}}_P(\pi )_\chi }
E(x,I_P(\pi ,f)\phi )\overline{\Lambda ^TE(y,\phi )}d\pi
\end{equation}
and $j_\chi ^T(x,f)=J_\chi ^T(x,x)$. Let 
\[
J_\chi ^T(f)
=\int_{G(\qqq )\backslash \overline{G(\aaa )}}
J_\chi ^T(x,x)dx 
\]
and let
\begin{equation}\label{eq:5.23}
J_\chi (f)=J_\chi ^T(f)|_{T=0}.
\end{equation}

\begin{lemma}\label{lem:5.23}
The integrals $K_\chi ^T(f),J_\chi ^T(f)$ are absolutely convergent and $%
K_\chi ^T(f)=J_\chi ^T(f)$ if $T$ is sufficiently regular.
\end{lemma}

\begin{remark}\label{rem:5.24}
The equality above is the metaplectic analog of \cite[Lemma 2.4]{A3}
for $SL(2)$.
\end{remark}

proof in the case $G=SL(2)$: The integral 
\[
K_\chi ^T(f)
=\int_{G(\qqq )\backslash \overline{G(\aaa )}}
\Lambda _2^TK_\chi (x,x)dx. 
\]
is absolutely convergent since the integrand is compactly
supported on $G(\qqq )\backslash \overline{G(\aaa )}$.

Let $F(P_1,P_2)$ denote the subset of elements of
$P_1(\qqq )\backslash P_2(\qqq )$ which do not belong
to for some parabolic $P$ with $P_1\subset
P\subset P_2,\;P\neq P_2$. Then 
\[
F(P_1,P_2)=\left\{ 
\begin{array}{c}
\{1\},\;\;\;\;P_1=P_2, \\ 
B(\qqq )\backslash G(\qqq )-\{1\},\;\;P_1=B,\;P_2=G,
\end{array}
\right. 
\]
and, from \cite[p. 105]{A3}, 
\begin{eqnarray*}
&&\int_{G(\qqq )\backslash \overline{G(\aaa )}}
k_\chi ^T(x,f)dx \\
&&=\sum_{P_1\subset P_2}(-1)^{\dim (A_2/Z)}
\int_{P_1(\qqq )\backslash \overline{G(\aaa )}}
\sum_{\gamma \in F(P_1,P_2)}\sigma _1^2(H(x)-T)
\Lambda_2^{T,P_1}K_{P_1,\chi }(\gamma x,x)dx \\
&&=\int_{G(\qqq )\backslash \overline{G(\aaa )}}
\sigma _G^G(H(x)-T)\Lambda_2^{T,G}K_{G,\chi }(x,x)dx\\
&&\;\;\;\;\;(the\ P_1=P_2=G\ term) \\
&&+\int_{B(\qqq )\backslash \overline{G(\aaa )}}
\sigma _B^G(H(x)-T)
\sum_{\gamma\in B(\qqq )\backslash G(\qqq )-\{1\}}
\Lambda _2^{T,B}K_{B,\chi }(\gamma x,x)dx\\
&&\;\;\;\;\;(the\ P_1 =B,\;P_2=G\ term) \\
&&-\int_{B(\qqq )\backslash \overline{G(\aaa )}}
\sigma _B^B(H(x)-T)\Lambda_2^{T,B}K_{B,\chi }(x,x)dx
\;\;\;\;\;(P_1 =P_2=B)
\end{eqnarray*}
Each of these integrals converge absolutely since the integral is continuous
and compactly supported. $\Box $

\begin{lemma}\label{lem:5.25}
$\sum_{ \chi =(M,\rho )\in \mathcal{X}, M\neq G }
\int_{G(\qqq )\backslash \overline{G(\aaa )}}
|j_\chi ^T(x,f)|dx<\infty $.
\end{lemma}

proof in the case $G=SL(2)$: Since
$K_P(x,y)=\sum_{\chi \in {\mathcal{X}}}K_{P,\chi }(x,y)$
is absolutely and uniformly convergent on compacta in
$x,y $, $K_P(x,y)$ is continuous in $x,y$. The same
reasoning as in the
previous lemma therefore applies here as well. $\Box $

\begin{lemma}\label{lem:5.26}
For $\chi \in \mathcal{X}$, we have 
\[
\int_{G(\qqq )\backslash \overline{G(\aaa )}}
j_\chi^T(x,f)dx=\sum_Pn(A)^{-1}
\int_{\Pi _u(\overline{M(\aaa )})}
\sum_{\phi \in {\mathcal{B}}_P(\pi )_\chi }
\Lambda ^TE(x,I_P(\pi ,f)\phi )\overline{\Lambda ^TE(y,\phi )}d\pi .
\]
\end{lemma}

proof: This follows from Lemma~\ref{lem:5.23} above on the
equality of $J_\chi ^T(f)$ and $K_\chi ^T(f)$, along with the
metaplectic analog of the fact that
fact $\Lambda^T\circ \Lambda ^T=\Lambda ^T$
\cite[Corollary 1.2]{A3} and that $\Lambda ^T$ is
self-adjoint \cite[Lemma 1.3]{A3} (see Lemma~\ref{lem:5.13} above). $\Box $

Let $L_P\subset {{\mathfrak{a}}}_P^G$ denote the integral
lattice generated by the
co-roots $\{\alpha ^{\vee }\;|\;\alpha \in \Delta _P\}$. For $\phi \in 
{\mathcal{H}}_{P_1}^0(\pi ),
\;s\in \Omega ({{\mathfrak{a}}}_1,{{\mathfrak{a}}}_2),\varsigma
\in {{\mathfrak{a}}}_{1,\mathbb{C}}^{*}$, let 
\[
M(s,\pi ,\varsigma )\phi =(M(s,\pi _\varsigma )\phi _\varsigma
)_{-s\varsigma }. 
\]
\begin{theorem}\label{thm:5.27}
If $\chi \in \mathcal{X}$ is an unramified class then (\cite[section 4]{A3}) 
\[
J_\chi ^T(f)=\sum_{P\in \mathcal{P}_\chi }
n(A)^{-1}\int_{\Pi _u(\overline{M(\aaa )})}
tr(M_P^T(\pi )_\chi I_P(\pi ,f)_\chi )d\pi ,
\]
where 
\[
M_P^T(\pi )_\chi =vol({{\mathfrak{a}}}_P^G/L_P)\lim_{\varsigma \rightarrow
0}\sum_{P_2\in \mathcal{P}_\chi }
\sum_{s\in \Omega ({{\mathfrak{a}}},{{\mathfrak{a}}}_2)}
\frac{e^{(s\varsigma )(T)}M(s,\pi )^{-1}M(s,\pi ,\varsigma )}{\prod_{\alpha
\in \Delta _2}(s\varsigma )(\alpha ^{\vee })}.
\]
\end{theorem}

For a similar but more general result (in the non-metaplectic case),
see \cite[Theorem 8.2]{A14}, which does not assume that $\chi$ is unramified.

\begin{example}\label{ex:5.28} For $G=SL(2)$, each $(G,\pi)\in\chi\in{\mathcal{X}}$ is
unramified. Moreover, $(A,\mu)\in\chi\in{\mathcal{X}}$ is unramified
if and only if $\mu^w\not=\mu$.
($P=G,\chi =(G,\rho )$) $M_G^T(\pi )_\chi =1$

($P=B,\chi =(G,\rho )$) $M_B^T(\pi )_\chi =1$

($P=B,\chi =(A,\rho )$ unramified) 
\[
\begin{array}{c}
M_B^T(\pi )_\chi \\
=vol({\mathbb{R}}/{\mathbb{Z}})
\lim_{\varsigma \rightarrow 0}
\left[
\frac{e^{\varsigma (T)}
M(1,\pi )^{-1}M(1,\pi ,\varsigma )}{\varsigma (\alpha^{\vee })}
+\frac{e^{-\varsigma (T)}
M(s,\pi )^{-1}M(s,\pi ,\varsigma )}{-\varsigma (\alpha^{\vee })}
\right]\\
=M(s,\pi)^{-1}
\lim_{\varsigma \rightarrow 0}
{M(s,\pi)-e^{-2\varsigma (T)}M(s,\pi,\varsigma)
\over \varsigma(\alpha^\vee)} ,
\end{array}
\]
where $\alpha $ denotes the unique positive root.
We identify ${\mathfrak{a}}^*$ with ${\mathbb{R}}$, ${\mathfrak{a}}$ with
${\mathbb{R}}$, and
${\mathfrak{a}}^*\times {\mathfrak{a}}\rightarrow {\mathbb{R}}$ with
the map $(x,y)\longmapsto 2xy$
of real numbers. Thus, at a function $\phi$
on $K$, the above limit becomes 
\[
\begin{array}{c}
M(s,\pi)^{-1}
\lim_{\varsigma \rightarrow 0}
{M(s,\pi)\phi(k)-e^{-2\varsigma (T)}e^{\varsigma (H(k))}
M(s,\pi_\varsigma)\phi_\varsigma(k)
\over \varsigma(\alpha^\vee)} \\
=M(s,\pi)^{-1}
\lim_{\varsigma \rightarrow 0}
{M(s,\pi)\phi(k)-e^{-2\varsigma (T)}M(s,\pi_\varsigma)\phi_\varsigma
\over \varsigma(\alpha^\vee)} \\
=-M(s,\pi)^{-1}M'(s,\pi)\phi(k)+4 \phi(k) T,
\end{array}
\]
where the prime denotes differentiation with respect to
$\varsigma$. Thus
\[
M_B^T(\pi )_\chi =
-M(s,\pi)^{-1}M'(s,\pi)+4\cdot I\cdot T,
\]
where $I$ denotes the identity. (There appears to be a sign
mistake here - the minus sign perhaps differs with some
calculations of Arthur. However, the sign it gives for the
corresponding expression in the trace formula
seems to agree with the corresponding sign in \cite{LL}.)

For
examples and properties of more general limits, see \cite[section 6]{A6}.
\end{example}


\subsection{The trace formula}

The geometric side of the trace formula is the sum of the following terms:

(a) the sum over classes ${\mathcal{O}}\subset G(\qqq )$
such that ${\mathcal{O}}\cap Z(G(\qqq ))\not= \emptyset$ of
\begin{equation}\label{eq:5.24}
J_{\mathcal{O}}(f)=
a^G(S,z)f(z)+{1\over 2}a^A(S,z)J_A(z,f),
\end{equation}
where $J_A$ and $a^M$ are given in \S\S\ref{ssec:JMu}--\ref{ssec:explicit} (these include the
``nilpotent terms''),

(b) the sum over elliptic regular classes ${\mathcal{O}}\subset G(\qqq )$
of
\begin{equation}\label{eq:5.25}
J_{\mathcal{O}}(f)=
a^G(S,\gamma)J_G(\gamma,f),
\end{equation}
where $\gamma\in {\mathcal{O}}$ and
$J_G$ (an orbital integral) and $a^G$ (a volume)
are given in \S\S\ref{ssec:JMu}--\ref{ssec:explicit},

(c) the sum over ``hyperbolic'' regular classes
${\mathcal{O}}\subset G(\qqq )$
of
\begin{equation}\label{eq:5.26}
J_{\mathcal{O}}(f)=
a^A(S,\gamma)J_A(\gamma,f),
\end{equation}
where $J_A$ (a weighted orbital integral) and $a^A$ (a volume)
are given in \S\S\ref{ssec:JMu}--\ref{ssec:explicit}.

The spectral side of the trace formula is the sum of the following terms:

(a) the sum over the cuspidal classes $\chi\in {\mathcal{X}}(G)$ of
\begin{equation}\label{eq:5.27}
J_\chi(f)=
{\rm tr}(\pi(f)),\ \ \ (\pi,G)\in\chi,
\end{equation}

(b) the sum over classes $\chi\in {\mathcal{X}}-{\mathcal{X}}(G)$ of
\begin{equation}\label{eq:5.28}
J_\chi(f)=
{1\over 2}\int_{\Pi_u(\overline{A(\aaa _S)})}
{\rm tr}(M_A(\pi)_\chi I_B(\pi,f)_\chi)d\pi,
\end{equation}
where $M_A(\pi)_\chi =M_A^T(\pi)_\chi|_{T=0}$ is
(essentially a logarithmic derivative of an intertwining operator)
given by Example~\ref{ex:5.28}.




\section{The global trace formula in invariant form}\label{sec:invariant}

The global trace formula is an equality of the form 
\begin{equation}\label{eq:6.1}
\sum_{\chi \in {\mathcal X}}J_\chi (f) =\sum_{{\mathcal{O}}\in {\mathcal C}}
J_{{\mathcal{O}}}(f),
\end{equation}
where $J_\chi (f)=J_\chi ^T(f)|_{T=0}$,
$J_{{\mathcal{O}}}(f)=J_{{\mathcal{O}}}^T(f)|_{T=0},$ and where
$J_\chi ^T(f)$ is given in Section~\ref{sec:spectral} and
$J_{{\mathcal{O}}}^T(f)$ is given in Section~\ref{sec:geometric}.

We can very crudely index the terms $J_{\mathcal{O}}$ by the centralizer
of the semi-simple part $s$ of a $\gamma =su\in {\mathcal{O}}$. As a
general principle, the ``smaller'' the centralizer the easier it is
to deal (analytically) with the term. 

Let ${\mathcal C}(G)$ denote the subset
of elliptic regular equivalence classes in ${\mathcal C}$ and let
${\mathcal{X}}(G)$
denote the subset of $\chi \in {\mathcal{X}}$ of the form $\chi=(G,\rho )$ for
some cuspidal
$\rho \in \Pi _u(\overline{G({\mathbb {A}})})$.

Not all the expressions in the above trace formula \eqref{eq:6.1} are invariant in
the sense that $J(f)=J(f^y)$ for all
$f\in {\mathcal{H}}(\overline{G({\mathbb {A}})}) $ and all
$y\in \overline{G({\mathbb {A}})}$. The terms $J_{{\mathcal{O}}}(f)$
with ${\mathcal{O}}\in {\mathcal{C}}(G)$ are invariant and the terms
$J_\chi (f)$ with $\chi \in {\mathcal{X}}(G)$ are invariant.
We follow \cite{A6}, \cite{A9}, \cite{A10}, in putting
the remaining terms in invariant form.

Fix a finite set of places $S$ of ${\mathbb {Q}}$ which
unless stated otherwise
will contain the archimedian place and the primes $p|2n$.
Let ${\mathbb {A}}_S=\Pi _{p\in S}{\mathbb {Q}}_p$ and embed
${\mathcal{H}}(\overline{G({\mathbb{A}}_S)})\hookrightarrow {\mathcal{H}}(G({\mathbb {A}))}$
in the natural way. We identify $f\in {\mathcal{H}}(\overline{G({\mathbb{A}}_S)})$
with the $f_S\in {\mathcal{H}}(G({\mathbb {A}}))$ which agrees
with $f$ on $\overline{G({\mathbb{A}}_S)}$ and with the unit
element of ${\mathcal{H}}(\overline{G({\mathbb{Q}}_p)},K_p)$
on $\overline{G({\mathbb{Q}}_p)} $ for $p\notin S$.

Define 
\begin{equation}\label{eq:6.2}
I_\chi (f)=I_\chi ^S(f)=J_\chi (f),\;\;\;\;\;\;\;\;\chi \in {\mathcal{X}}(G),
\end{equation}
and 
\begin{equation}\label{eq:6.3}
I_{{\mathcal{O}}}(f)=I_{{\mathcal{O}}}^S(f)=J_{{\mathcal{O}}}(f),
\;\;\;\;\;\;\;\;{\mathcal{O}}\in {\mathcal C}(G),
\end{equation}
for $f\in {\mathcal{H}}(\overline{G({\mathbb{A}}_S)})$. One of the objects to this subsection
is to define $I_\chi (f)=I_\chi ^S(f)$,
$\chi \in {\mathcal{X}}-{\mathcal{X}}(G), $
and $I_{{\mathcal{O}}}(f)=I_{{\mathcal{O}}}^S(f)$,
${\mathcal{O}}\in {\mathcal C}-{\mathcal C}(G),\;$%
following \cite{A6}. Once this is done, we will use the formalism of \cite{A6} to
deduce the global invariant trace formula in the form 
\begin{equation}\label{eq:6.4}
\sum_{\chi \in {\mathcal{X}}}I_\chi (f)=\sum_{{\mathcal{O}}\in {\mathcal C}}
I_{{\mathcal{O}}}(f).
\end{equation}

To do this requires, in particular, verifying for each $p\in S$ the local
assumptions of \cite{A6}, of which there are three types:

(1) \cite[Assumption 5.1, p. 30]{A6},

(2) \cite[Assumption 5.2, p. 32]{A6},

(3) the statements \cite[(7.1)--(7.6)]{A6} regarding normalized intertwining operators.
(In our case, only the statements \cite[(7.1)--(7.4)]{A6} need proving since
$SL(2)/k$ is rank one.)

\subsection{Local assumptions}\label{ssec:localass}

We shall now make precise three types of statements alluded to
above which must be verified.

Let $Q=M_Q N_Q$ be a parabolic of $G$ and let $y\in G({\mathbb{A}}_S)$.
For $m\in M({\mathbb{A}}_S)$, let
\begin{equation}\label{eq:6.5}
f_{Q,y}(m)=\delta_Q(m)^{1/2}\int_K\int_{N_Q({\mathbb{A}}_S)}
f(k^{-1}mnk)u'_Q(k,y)dndk,
\end{equation}
where
\[
u'_Q(k,y)=\int_{{\mathfrak{a}}_Q}
\sum_{ P ,Q\subset P}
(-1)^{dim(A_Q/A_G)}
\tau_Q^P(X)\hat{\tau}_P(X+H_Q(ky))dX ,
\]
as in \cite[{3.3}]{A6} (modulo a typo).
(Recall \eqref{eq:2.4}, \eqref{eq:2.5} for $\tau_Q^P,...$.)
\begin{example}\label{ex:6.1} Let $G=SL(2)$. Then
\[
\begin{array}{c}
u'_G(k,y)=\tau_G^G(X)\hat{\tau}_G(X+H_G(ky))=1,\\
u'_Q(k,y)=\int_{{\mathbb{R}}}
-\tau_B^B(X)\hat{\tau}_B(X+H_B(ky))
+\tau_B^G(X)\hat{\tau}_G(X+H_B(ky)) dX \\
=\int_{{\mathbb{R}}}
-\hat{\tau}_B(X+H_B(ky))+\tau_B(X)dX\\
=|H_B(ky)|.
\end{array}
\]
\end{example}

\subsubsection{Assumption 5.1}\label{sssec:ass51}

First, a

\begin{definition}\label{defn:6.2}
Let $\phi :U\rightarrow V$ denote a continuous linear transformation between
two complete topological vector spaces. Let $U^{\prime }$ denote the space
of continuous linear functionals on $U$ (similarly for $V^{\prime }$). The 
{\bf transpose} $\phi ^{\prime }:V^{\prime }\rightarrow U^{\prime }$ is the
map defined by 
\[
\phi ^{\prime }(\lambda ):f\longmapsto \lambda (\phi (f)),\;\;\;\lambda \in
V^{\prime },\;f\in U. 
\]
\end{definition}

Let $C_N^\infty({\mathfrak{a}}_{M,S})$ denote the topological vector space of
smooth functions on ${\mathfrak{a}}_{M,S}$ which are supported on
\[
\{X\in {\mathfrak{a}}_{M,S}\ |\ ||X||\leq N\},
\]
where
\begin{equation}\label{eq:6.6}
{\mathfrak{a}}_{M,S}
=\{H_M(m)\ |\ m\in \overline{M({\mathbb{A}}_S)}\}.
\end{equation}

We follow \cite[section 11]{A11} (as opposed to \cite[sections 4,5]{A6} where
the assumptions are stated), restricting ourselves to
$G=SL(2)$ and simplifying whenever possible.

 Let $N,\Gamma$ be as in \S\ref{sssec:defs} and let
$I_N(\overline{G({\mathbb{A}}_S)})_\Gamma$ denote the vector
space of complex valued functions $\phi$
on the tempered dual $\Pi_t(\overline{G({\mathbb{A}}_S)})$ such
that

(a) if $\pi|_K$ contains no representation in $\Gamma$ then
$\phi (\pi)=0$,

(b) if $M\subset G$ is an admissible Levi (e.g., $M=A$) and
$\sigma\in \Pi_t(\overline{M({\mathbb{A}}_S)})$ then
\begin{equation}\label{eq:6.7}
\phi(\sigma,X)
=\int_{{\mathfrak{a}}_{M,S}}
\phi(\sigma,X)e^{-\Lambda (X)}d\Lambda ,\ \ \ \ X\in {\mathfrak{a}}_{M,S},
\end{equation}
converges to a function of $X$ which belongs to
$C_N^\infty({\mathfrak{a}}_{M,S})$.

If we fix a family of semi-norms $||...||_{M,\sigma}$
on $C_N^\infty({\mathfrak{a}}_{M,S})$ then there are seminorms
\begin{equation}\label{eq:6.8}
\phi\longmapsto ||\phi(\sigma,*)||_{M,\sigma},\ \ \ \
\sigma\in \Pi_t(\overline{M({\mathbb{A}}_S)}),
\end{equation}
which make $I_N(\overline{G({\mathbb{A}}_S)})_\Gamma$ into a topological vector space.
By \cite[ch 1]{Ru}, we may complete
$I_N(\overline{G({\mathbb{A}}_S)})_\Gamma$ using this family of seminorms to
obtain a Frech\'et space.

Define
\begin{equation}\label{eq:6.9}
I(\overline{G({\mathbb{A}}_S)})
=\lim_{\Gamma}
I(\overline{G({\mathbb{A}}_S)})_\Gamma
\end{equation}
where
\[
I(\overline{G({\mathbb{A}}_S)})_\Gamma
=\lim_N
I_N(\overline{G({\mathbb{A}}_S)})_\Gamma .
\]

In \cite{A11}, the spaces $I_{ac}(\overline{G({\mathbb{A}}_S)})$ are also
introduced, as well as the subspace of moderate functions,
which we denote by $I_{mod}(\overline{G({\mathbb{A}}_S)})$.
Actually only the non-metaplectic versions of these spaces
were introduced in \cite{A11}, their metaplectic analogs are left
to the reader.
Note that for $G=SL(2)$ or $G=A$, the diagonal maximal split torus
of $SL(2)$, we have
\[
I_{ac}(\overline{G({\mathbb{A}}_S)})=I(\overline{G({\mathbb{A}}_S)})
=I_{mod}(\overline{G({\mathbb{A}}_S)}),
\ \ \ \ 
I_{ac}(\overline{A({\mathbb{A}}_S)})
=I(\overline{A({\mathbb{A}}_S)})
=I_{mod}(\overline{A({\mathbb{A}}_S)}).
\]

Let $M\in \{A,G\}$, let ${\mathcal{H}}_{ac}$ be as in
\cite{A11} (see also \S\ref{sssec:defs} above),
and let 
\[
U_M={\mathcal{H}}_{mod}(\overline{M({\mathbb{Q}}_p)})
={\mathcal{H}}(\overline{M({\mathbb{Q}}_p)}) .
\]
These are smooth, compactly supported (in our case), $K$-finite
functions. Let 
\[
V_M=I_{mod}(\overline{M({\mathbb{Q}}_p)}) =I(\overline{M({\mathbb{Q}}_p)})
\]
denote the space of functions on $\Pi _u(\overline{M({\mathbb{Q}}_p)})$ as
in \cite{A11} (and as above).
Note that these function spaces are not the same as those in \cite{A6}.
I will assume that \cite[Assumption 5.1]{A6} is intended
for the spaces in \cite[section 11]{A11}.
Then the assumptions are

\begin{conjecture}\label{conj:6.3}
The map $\phi=\phi_M^M :U_M\rightarrow V_M$ given by
$\phi (f):\pi \longmapsto tr(\pi (f)) $ is surjective.
\end{conjecture}

The non-metaplectic analog of this
was proven in \cite[Lemma 6.1]{A10} (see also \cite{BDK}).
For a discussion of the metaplectic conjecture, see \cite{J1}.

\begin{conjecture}\label{conj:6.4}
The image of $\phi ^{\prime }:V_M^{\prime }\rightarrow U_M^{\prime }$ is the
space of all invariant distributions on
$\overline{M({\mathbb {Q}}_p)}$.
\end{conjecture}

In \cite[Lemma 8.3, Theorem 12.1]{A6}, it is shown that these imply

\begin{theorem}\label{thm:6.5} Assume the conjectures above.
There is a map $\phi^G_M :U_G\rightarrow V_M$ satisfying

(a) if $G=M$ then $\phi^G_M$ is as above,

(b) in general, we have
\[
\phi_M^G(f^y)=
\sum_{Q ,M\subset M_Q}
\phi_M^{M_Q}(f_{Q,y}),\ \ \ \ f\in U_G.
\]
In particular, $\phi_G^G(f^y)=\phi_G^G(f)$.

(c) for $f\in U_G$, $\pi\in \Pi_t(\overline{M({\mathbb{A}}_S)})$,
we have $\phi^G_M(f,\pi)={\rm tr}(I_{P_0}(\pi,f)R_M^G(\pi,P_0))$,
where $R_M^G(\pi,P_0))$ is defined in \cite[section 7]{A6}
via the $(G,M)$-family
\[
R_P(\lambda, \pi,P_0))=
R_{P|P_0}(\pi)^{-1}R_{P|P_0}(\pi_\lambda),
\]
where the $R_{P|P_0}(\pi_\lambda)$ are as in \S\ref{sssec:normint} below.
\end{theorem}

\begin{example}\label{ex:6.6} Let $G=SL(2)$. In this case,
\[
\phi^G_A(f,\pi)={\rm tr}(I_{B}(\pi,f)R_A^G(\pi,B))
\]
is a global ``weighted character'' similar to the integrand of
the expression in Theorem~\ref{thm:5.27} above. The term $R_A^G(\pi,B)$
calculated in a similar fashion to $M_B^T(\pi)_\chi$
in Example~\ref{ex:5.28}.
\end{example}

Let $\sigma\in \Pi(\overline{A({\mathbb{A}}_S)})$ and let
\[
R_{\overline{B}|B}(\sigma)
=m(\sigma)^{-1}M_{\overline{B}|B}(\sigma)
\]
Then we may write 
\[
R_{\overline{B}|B}(\sigma_\lambda)^{-1}
R'_{\overline{B}|B}(\sigma_\lambda)|_{\lambda =0}
=-m(\sigma_\lambda)^{-1}m'(\sigma_\lambda)|_{\lambda =0}
+M_{\overline{B}|B}(\sigma_\lambda)^{-1}
M'_{\overline{B}|B}(\sigma_\lambda)|_{\lambda =0}
\]
Eventually we will want to plug this into the term \eqref{eq:5.28}
(using the last equation in Example~\ref{ex:5.28}). 

\subsubsection{Assumption 5.2}\label{sssec:ass52}

Let $M\in \{A,G\}$, $S=\{p\}$.

A complex valued function $\phi$ on
$\Pi(\overline{M({\mathbb{Q}}_p)})$ extends to the
complex Grothendieck group
$R(\Pi(\overline{M({\mathbb{Q}}_p)}))\otimes {\mathbb{C}}$ by linearity.
Thus, if $L$ is any Levi of $M$ then $\phi$ is defined on the image of
the induction map
$I_L^M: \Pi(L({\mathbb{Q}}_p))\rightarrow R(\Pi(\overline{M({\mathbb{Q}}_p)}))$.
In particular, for $\sigma\in \Pi(\overline{M({\mathbb{Q}}_p)})$ and
$\zeta\in \oplus_{v\in S}i{\mathfrak{a}}^*_v$, $\phi(I_L^M(\sigma_\zeta))$ is
defined. Let
\begin{equation}\label{eq:6.10}
D_\sigma \phi(\sigma^M)=
\left\{
\begin{array}{c}
{d\over d\zeta}\phi(I_L^M(\sigma_\zeta))|_{\zeta =0},
\ \ \ \ {\rm if\ it\ exists},\\
\infty,\ \ \ \ {\rm otherwise}.
\end{array}
\right.
\end{equation}
If $p<\infty$ is a prime and
$\sigma\in \Pi(\overline{M({\mathbb{Q}}_p)})$ then define
$||\sigma||=0$. If $p=\infty$,
$\sigma\in \Pi(\overline{M({\mathbb{Q}}_p)})$,
and $\Delta$
is a Casimer operator on $\overline{M({\mathbb{Q}}_p)}$
(which is the same as the Casimer operator on
$M({\mathbb{Q}}_p)$)
then define $||\sigma||$ to be the
smallest (positive) eigenvalue of $\sigma(\Delta)$. For
$\sigma\in \Pi(\overline{M({\mathbb{Q}}_p)})$, define
\[
||\sigma||=\sup_{p\in S}||\sigma_p||.
\]

If $U_M={\mathcal{C}}(\overline{M({\mathbb{Q}}_p)})$, let
$V_M$ be defined as follows.
Let
\[
\begin{array}{c}
{\mathcal{I}}(\overline{M({\mathbb{Q}}_p)})_{\overline{K}}
=\{\phi: \Pi(\overline{M({\mathbb{Q}}_p)})\rightarrow {\mathbb{C}} \ | \\
(a) \ \phi(\pi)=0\ {\rm \ unless\ }\pi{\rm \ has\ a\ }
\overline{K\cap M({\mathbb{Q}}_p)}-{\rm \ fixed\ vector},\\
(b)\ {\rm for \ each\ Levi\ }L{\rm\ of\ }M,\\
||\phi||_{D,n}=\sup_{\sigma\in \Pi_t(\overline{L({\mathbb{Q}}_p)})}
(1+||\sigma||)^n|D_\sigma \phi(\sigma^M)|\\
{\rm is\ finite}\},
\end{array}
\]
topologized by the semi-norm topology. Let
\begin{equation}\label{eq:6.11}
V_M=
{\mathcal{I}}(\overline{M({\mathbb{Q}}_p)})
=\lim_{K}{\mathcal{I}}(\overline{M({\mathbb{Q}}_p)})_{\overline{K}},
\end{equation}
topologized using the direct limit topology.
See also \cite[section 3]{A16}.

Then the assumptions are

\begin{conjecture}\label{conj:6.7}
The map $\phi =\phi_M^M :U_M\rightarrow V_M$ given by
$\phi (f):\pi \longmapsto tr(\pi (f)) $ is surjective.
\end{conjecture}

This is proven in \cite[appendix]{A16}.

\begin{conjecture}\label{conj:6.8}
The image of $\phi ^{\prime }:V_M^{\prime }\rightarrow U_M^{\prime }$ is the
space of all tempered invariant distributions on
$\overline{M({\mathbb{Q}}_p)}$.
\end{conjecture}

When $M=G=SL(2)$ these conjectures are discussed in \cite{J1}, where
they are shown to be a consequence of other conjectures (!).

In \cite{A6}, it was shown that these imply
\begin{theorem}\label{thm:6.9} Assume the two conjectures above.
There is a map
$\phi^G_M :U_G\rightarrow V_M$ satisfying

(a) if $M=G$ then $\phi_M^G$ is as above.

(b) In general, we have
\[
\phi_M^G(f^y)=
\sum_{ Q , M\subset M_Q}
\phi_M^{M_Q}(f_{Q,y}),\ \ \ \ f\in U_G.
\]
In particular, $\phi_G^G(f^y)=\phi_G^G(f)$.

(c) $\phi_M^G$ is given as in Theorem~\ref{thm:6.5}(c) above.
\end{theorem}

\subsubsection{Normalized intertwining operators}\label{sssec:normint}

Let $P,P^{\prime },P^{\prime \prime }$ be parabolics of
$G$ defined over ${\mathbb{Q}}$
with Levi $A$. Let $\chi \in \Pi _u({\overline{A({\mathbb{Q}}_p)}})$ have
representation space $V_\chi $ - a finite dimensional complex vector space
with an ${A({\mathbb {Q}}_p)}$-invariant inner product
$(...,...)_\chi $ - and let ${\mathcal{K}}_P(\chi )$ denote the space of all
genuine $V_\chi $-valued functions $\phi $
whose norms are square-integrable over
${K_p}$: 
\[
\int_{K_p}(\phi ,\phi )_\chi dk<\infty . 
\]

\begin{conjecture}\label{conj:6.10}
There are canonically normalized intertwining operators 
\[
R_{P^{\prime }|P}(\chi ):{\mathcal{K}}_P(\chi )\rightarrow
{\mathcal{K}}_{P^{\prime}}(\chi ) 
\]
such that

(1) for all $x\in \overline{G({\mathbb{Q}}_p)}$, we have 
\[
R_{P^{\prime }|P}(\chi )I_P(\chi ,x)
=I_P(\chi ,x)R_{P^{\prime }|P}(\chi ), 
\]

(2) $R_{P^{\prime }|P}(\chi )$ is unitary and 
\[
R_{P^{\prime \prime }|P}(\chi )
=R_{P^{\prime \prime }|P^{\prime }}(\chi)R_{P^{\prime }|P}(\chi ), 
\]

(3) if $\psi \in {\mathcal K}_P(\chi )$ is such that
$\psi (k)=\psi (1)$
for all $k\in \overline{K_p}$ then 
\[
R_{P^{\prime }|P}(\chi )\psi =\psi , 
\]

(4) (for $p=\infty $) if $W\in \Pi (\overline{K_p})$ and 
\[
R_{P^{\prime }|P}(\chi )_W=R_{P^{\prime }|P}(\chi )|_{{\mathcal{K}}_P(\chi )_W}
\]
then for all $D\in Diff(i{\mathfrak {a}}_P^{*})$ there is a $C>0,\;N>0$ such
that 
\[
||D_\chi R_{P^{\prime }|P}(\chi )_W||\leq C\cdot (1+||\chi ||)^N\cdot
(1+||W||)^N,
\]
for all $\chi \in \Pi _u(\overline{A({\mathbb{Q}}_p)})$.
Here $||\chi||,\;||W||$ are as in \cite[section 5]{A6}
(see also \S\ref{sssec:ass52} above).
\end{conjecture}

If $p=\infty $, $G=SL(2)$ then these follow either from \cite{Sa}.
For $p<\infty$, they are discussed in \cite{J1} (see \cite[section 2]{A16}
for the case $n=1$).

\subsection{The terms in the invariant trace formula}

Recall ${\mathcal C}$ denotes the set of equivalence classes in $G({\mathbb {Q}})$
(two elements are equivalent if and only if their semi-simple
Jordan components are $G({\mathbb{Q}})$-conjugate). Let
$\sigma \in G({\mathbb{Q}})$ be semi-simple. The
$(G,S)${\bf -equivalence class of }$\sigma $ is
the set 
\[
G({\mathbb{Q}})\cap (\sigma {\mathcal{U}})^{G({\mathbb{Q}})}
=\{\delta ^{-1}\sigma u\delta
\;|\;\delta \in G({\mathbb{Q}}),\;u\in {\mathcal{U}}\cap G({\mathbb{Q}})\}, 
\]
where ${\mathcal{U}}$ is a unipotent conjugacy class in
$\overline{G({\mathbb {A}}_S)}_\sigma$.

\begin{lemma}\label{lem:6.11}
Each ${\mathcal{O}}\in {\mathcal{C}}$ breaks up into a finite set of
$(G,S)$-classes ${\mathcal{O}}_{(G,S)}$.
\end{lemma}

More generally, we may define this
finer notion of equivalence for any Levi $M$. Let
$\gamma \in M({\mathbb{Q}})$ have semi-simple Jordan
component $\sigma$. Let $\gamma ^{\prime }\in M({\mathbb{Q}})$.
We say that $\gamma ^{\prime }$
is $(M,S)${\bf -equivalent} to $\gamma $ if there is a
$\delta \in M({\mathbb{Q}})$ such that

(1) $\sigma $ is the semi-simple Jordan component of
$\delta ^{-1}\gamma^{\prime }\delta $,

(2) $\sigma ^{-1}\gamma $ and $\sigma ^{-1}\delta ^{-1}
\gamma ^{\prime}\delta $ are
$\overline{M({\mathbb{A}}_{S})}_\sigma$-conjugate.

\begin{example}\label{ex:6.12}
Let $G=SL(2)$, $\gamma,\gamma^{\prime}\in G({\mathbb{Q}})$
and let $S$ be non-empty.

Let $M=G$ and $\gamma =1$. $\gamma$ is $(G,S)$-equivalent to
$\gamma^{\prime}$ if and only if they are equal.

Let $M=A$ and $\gamma =1$. $\gamma$ is $(A,S)$-equivalent to
$\gamma^{\prime}$ if and only if they are equal.

Let $M=G$ and let $\gamma $ be regular semi-simple.
$\gamma$ is $(G,S)$-equivalent to
$\gamma^{\prime}$ if and only if they are $G({\mathbb{Q}})$-conjugate.

Let $M=A$ and $\gamma \in A({\mathbb{Q}})$ be regular.
$\gamma$ is $(A,S)$-equivalent to
$\gamma^{\prime}$ if and only if they are equal.
\end{example}

\begin{lemma}\label{lem:6.13}
If $G=SL(2)$ and $S$ is non-empty then
$\gamma^{\prime}$ is $(G,S)$-equivalent to $\gamma\in G({\mathbb{Q}})$ only if 
$\gamma^{\prime}$ is equivalent to $\gamma$.
\end{lemma}

\begin{example}\label{ex:6.14}
Let $G=SL(2)$ and let ${\mathcal{O}}\in {\mathcal C}-{\mathcal C}(G)$, so
${\mathcal{O}}$ is not an elliptic class. If $a$ denotes the Jordan component of
some $\gamma \in {\mathcal{O}}$ and $a$ is regular then we have
${\mathcal{O}}\cap A({\mathbb{Q}})=\{a,a^w\}$.
\end{example}

Let $G=SL(2)$. Assume that
${\mathcal{O}}\in {\mathcal{C}}-{\mathcal{C}}(G)$
contains a regular element.
We next define a term $J_{{\mathcal{O}}}^{\prime }(f)$
such that 
\begin{equation}\label{eq:6.12}
I_{{\mathcal{O}}}(f)=J_{{\mathcal{O}}}(f)-J_{{\mathcal{O}}}^{\prime }(f)
\end{equation}
is an invariant distribution on ${\mathcal{H}}(\overline{G({\mathbb{A}}_S)})$.
To define $J_{\mathcal{O}}^{\prime }(f)$, let
$\chi \in \Pi _u(\overline{A({\mathbb{A}}_S)})$, let
$M(\chi )=M(w,\chi)$ denote the intertwining operator
\[
M(w,\chi):I_A(\chi )\rightarrow I_A(\chi ^w),
\]
and let $M(\chi ,\eta )$ denote the operator
$M(w,\eta )=M(w,I_A(\chi ),\eta)$ in \cite[section 4, p. 113]{A2},
where $\eta \in {\mathfrak {a}}_{{\mathbb {C}}}^{*},\;P_1=B,\;P_2=B^w$.
Let $m(\chi ,\eta )$ denote the scalar factor such that 
\begin{equation}\label{eq:6.13}
R(\chi ,\eta )=m(\chi ,\eta )^{-1}M(\chi ,\eta ):
{\mathcal K}_B(\chi )\rightarrow {\mathcal K}_B(\chi ^w)
\end{equation}
satisfies the normalized intertwining operator satisfying the conditions in
\S\ref{sssec:normint} above. The relation between the logarithmic
derivatives of the two is given by
\begin{equation}\label{eq:6.14}
M(\chi ,\eta )^{-1}M'(\chi ,\eta )=
m(\chi ,\eta )^{-1}m'(\chi ,\eta )+
R(\chi ,\eta )^{-1}R'(\chi ,\eta ).
\end{equation}

Let 
\begin{equation}\label{eq:6.15}
\varphi _A(f,\chi )
={\rm tr}(R(\chi ,0)^{-1}R(\chi ,0)^{\prime }I_A(\chi ,f)),
\end{equation}
where the derivative (the prime on $R$) is with respect to $\eta $ and then
evaluated at $0$. This makes sense since the range of $R(\chi ,\eta )$ is
independent of $\eta $.
Roughly speaking, the correction term $J_{\mathcal{O}}^{\prime }(f)$ can be
regarded as an average of the inverse Fourier transform of $\varphi _A(f)$
over ${\mathcal{O}}\cap A({\mathbb{Q}})$
(by analogy with the case of $SL(2)$ or $GL(2)$
- see the introduction to \cite{A6} and \cite{LL}). This is
essentially the way
it works in the formalism here but more work will be required to set things
up properly.

We now turn to \cite[section 4]{A6} and recall from \S\ref{ssec:rank0corr} above  the
definition of the invariant distributions
$I_{{\mathcal{O}}}^A,\;I_\chi ^A$ on $\overline{A({\mathbb{A}}_S)}$.
These will be used to define $I_{{\mathcal{O}}}^G,\;I_\chi ^G$ on 
$\overline{G({\mathbb{A}}_S)}$ using the metaplectic analog of
\cite[Proposition 4.1 and its proof]{A6}. In fact, we
have 
\begin{equation}\label{eq:6.16}
J_{{\mathcal{O}}}^{\prime }(f)
=|W_G|^{-1}(I_{{\mathcal{O}}}^A)^{\wedge }(\varphi_A(f)),
\end{equation}
where $(I_{{\mathcal{O}}}^A)^{\wedge }$ denotes the distribution on the
Schwartz space ${\mathcal{S}}(\Pi _u(A({\mathbb {A}}_S)))$ such that 
\[
(\varphi _G^A)^{\prime }((I_{\mathcal{O}}^A)^{\wedge })=I_{\mathcal{O}}^G, 
\]
where 
\[
\varphi_G^A:{\mathcal{H}}(\overline{G({\mathbb {A}}_S)})\rightarrow
I(\Pi _u(\overline{A({\mathbb {A}_S)})) }
\]
is the Fourier transform map.

\begin{proposition}\label{prop:6.15}
(\cite[Introduction and Proposition 4.1]{A6}) The distribution
\eqref{eq:6.12} on ${\mathcal{H}}(\overline{G({\mathbb{A}}_S)})$ is invariant,
where $J_{\mathcal{O}}^{\prime }(f)$ is given
by \eqref{eq:6.16}.
\end{proposition}

Assume that ${\chi }\in {\mathcal{X}}-{\mathcal{X}}(G)$.
Next we define a term $J_{{ \chi }}^{\prime }(f)$ such that 
\begin{equation}\label{eq:6.17}
I_{{\chi }}(f)=J_{{\chi }}(f)-J_{{\chi }}^{\prime }(f)
\end{equation}
is an invariant distribution on
${\mathcal{H}}(\overline{G({\mathbb{A}}_S)})$.
As in \eqref{eq:6.16}, let 
\begin{equation}\label{eq:6.18}
J_{{\chi }}^{\prime }(f)=|W_G|^{-1}(I_{{\chi }}^A)^{\wedge
}(\varphi _A(f)),
\end{equation}
where $(I_{{\chi }}^A)^{\wedge }$ denotes the distribution on the
space $I(\Pi _u(\overline{A({\mathbb {A}}_S)}))$ such that 
\[
(\varphi _G^A)^{\prime }((I_{{\chi }}^A)^{\wedge })=I_{{\chi }}^G. 
\]

\begin{proposition}\label{prop:6.16}
(\cite[Introduction and Proposition 4.1]{A6}) The distribution \eqref{eq:6.17}
on ${\mathcal{H}}(\overline{G({\mathbb{A}}_S)})$ is invariant,
where $J_{{\chi }}^{\prime }(f)$ is
given by \eqref{eq:6.18}.
\end{proposition}

The invariant formula is the equality given by \eqref{eq:6.4}.

\subsection{The terms $I_M(\gamma, f)$ and $I_M(\pi,f)$}

The term $I_M(\gamma,f)$ is defined inductively by
\begin{equation}\label{eq:6.19}
I_M(\gamma,f)
=J_M(\gamma,f)-
\sum_{ L, M\subset L\\ L\not= G}
\hat{I}_M^L(\gamma,\phi_L(f)),
\end{equation}
where $\phi_L$ is as in \S\ref{ssec:localass}.
For a general explicit expression, see \cite{A16}.

There is an ``invariant version'' of Proposition~\ref{prop:4.15}, where $J$ is
replaced by $I$, (substitute \eqref{eq:6.12},
\eqref{eq:6.16}, and \eqref{eq:6.19} into Proposition~\ref{prop:4.15}):

Let $(M({\mathbb{Q}})\cap {\mathcal{O}})_{M,S}$ denote the set of
$(M,S)$-equivalence classes in $M({\mathbb {Q}})\cap {\mathcal{O}}$.

\begin{proposition}\label{prop:6.17}
(\cite[Theorem 3.3]{A10}) 
\[
I_{{\mathcal{O}}}(f)
=\sum_{M\in {\mathcal{L}}(G)}|W_0^M|\cdot |W_0^G|^{-1}
\sum_{\gamma \in (M({\mathbb{Q}})\cap {\mathcal{O}})_{M,S}}
a^M(S,\gamma )I_M(\gamma ,f),
\]
where $S$ is a sufficiently large set of valuations and the
$a^M(S,\gamma)$ are as in \S\ref{ssec:explicit}.
\end{proposition}

In \S\ref{sssec:normint} we introduced normalized intertwining operators
$R_{P^{\prime}|P}(\pi)$, for $P^{\prime},P\in {\mathcal P}(M)$.
We will use these to make the distributions
$I_M(\gamma)$ and $I_M(\pi)$ more explicit.
For a fixed $P_0\in {\mathcal P}(M)$,
$\nu\in i{\mathfrak{a}}_M^*$,
$\pi\in \Pi_u(\overline{M({\mathbb{A}}_S)})$,
$\lambda\in {\mathfrak{a}}_{M,{\mathbb{C}}}^*$, let
\begin{equation}\label{eq:6.20}
R_P(\mu,\pi_\lambda,P_0)
=R_{P|P_0}(\pi_\lambda)^{-1}R_{P|P_0}(\pi_{\lambda+\mu}) .
\end{equation}
This more than likely (the metaplectic analog of) a
$(G,M)$-family \cite[section 6]{A11}. In any case, the quantity
\[
R_M(\pi_\lambda,P_0)
=\lim_{\mu\rightarrow 0}
\sum_{P\in {\mathcal P}(M)}
R_P(\mu,\pi_\lambda,P_0)\theta_P(\mu)^{-1}
\]
exists. Here $R_G(\pi_\lambda,P_0)=1$ and
\[
R_A(\pi_\lambda,P_0)=
R_{\overline{P_0}|P_0}(\pi_\lambda)^{-1}
{d\over d\lambda}
R_{\overline{P_0}|P_0}(\pi_\lambda)
\]
by the computation in Example~\ref{ex:5.28} (see also \cite{A11} and \cite[(6.5)]{A6}).

\begin{example}\label{ex:6.18}
If $G=SL(2)$ and $M=A$ then
\[
I_A(\gamma,f)
=J_A(\gamma,f)-
\hat{I}_A^A(\gamma,\phi_A(f)),
\]
for $\gamma\in A({\mathbb{Q}})$.
The term $\hat{I}_A^A(\gamma,\phi_A(f))$
is the inverse Fourier transform of
$\phi_A(f)(\chi)=tr(I_A(\chi)R_A(\chi,B))(f)$ evaluated at
$\gamma$. 
\end{example}

Let
\begin{equation}\label{eq:6.21}
J_M(\pi_\lambda,f)=
{\rm tr}(R_M(\pi_\lambda,P_0)
I_{P_0}(\pi_\lambda,f)),
\end{equation}
where $f\in {\mathcal{H}}(\overline{G({\mathbb{A}}_S)})$ and where
$I_{P_0}(\pi)$ denotes the usual induced representation.
Define
\begin{equation}\label{eq:6.22}
J_M(\pi,X,f)
=\int_{i{\mathfrak{a}}^*_{M,S}}
J_M(\pi_\lambda,f)e^{\lambda(X)}d\lambda,
\end{equation}
where $f$ is in the Hecke algebra as above, $X\in {\mathfrak{a}}_{M,S}$
is as in \eqref{eq:6.6}, and where

\begin{eqnarray*}
&&{\mathfrak{a}}^*_{M,S}
={\mathfrak{a}}^*_{M}/{\mathfrak{a}}^{\vee}_{M,S}\\
&&{\mathfrak{a}}^{\vee}_{M,S}={\rm Hom}({\mathfrak{a}}_{M,S},2\pi {\mathbb Z}).
\end{eqnarray*}
For
$\pi\in \Pi_u(\overline{M({\mathbb{A}}_S)})$,
inductively define $I_M(\pi,X,f)=I_M^G(\pi,X,f)$ by
\begin{equation}\label{eq:6.23}
I_M(\pi,X,f)
=J_M(\pi,X,f)-
\sum_{ L, M\subset L\\ L\not= G}
\hat{I}_M^L(\pi,X,\phi_L(f)),
\end{equation}
where $\phi_L$ is as in \S\ref{ssec:localass} and \cite[section 3]{A9}.
In general, define $I_A(\pi,f)=I_A(\pi,0,f)$.

\begin{example}\label{ex:6.19}
If $G=SL(2)$ and $M=A$ then
\[
I_A(\pi,X,f)
=J_A(\pi,X,f)-
\hat{I}_A^A(\pi,X,\phi_A(f))
=J_A(\pi,X,f)-(\phi_A(f))(\pi),
\]
for $\pi\in \Pi_u(\overline{A({\mathbb{A}}_S)})$.
\end{example}

% Claude (2026): the original read "The analog of Proposition 4.16 is";
% the intended target is unclear (the 1997 book has "Proposition 320").
We have

\begin{proposition}\label{prop:6.20}
(\cite[Lemma 3.1]{A9}) If $\pi\in\Pi_t(\overline{M({\mathbb{A}}_S)})$ then
\[
I_M(\pi,X,f)=
\left\{
\begin{array}{c}
 f_G(\pi,X),\ \ M=G,\\
 0,\ \ M\not= G,
\end{array}
\right.
\]
where
\[
f_G(\pi,X)=
\int_{i{\mathfrak{a}}^*_{G,S}}
{\rm tr}(\pi_\lambda(f))e^{\lambda(X)}d\lambda,
\]
and $X\in {\mathfrak{a}}_{G,S}$.
\end{proposition}

(This vanishing result for $I_M,\ M\not=G$ is false in
general if $\pi$ is not tempered.

\begin{example}\label{ex:6.21} If $G=SL(2)$ and
$\pi\in\Pi_t(\overline{M({\mathbb{A}}_S)})$ then
\[
I_M(\pi,X,f)=
\left\{
\begin{array}{c}
 tr(\pi(f)),\ \ M=G,\\
 0,\ \ M= A.
\end{array}
\right.
\]
\end{example}

We assume (as usual) that $S$ contains the archimedian place.
If $\pi$ is an automorphic representation of $\overline{M({\mathbb{A}}_S)}^1$,
the infinitesimal character of the archimedian
component of $I_M(\pi)$, is a $W$-orbit in the set of
${\mathfrak{a}}_{M,{\mathbb{C}}}^*$-orbits in ${\mathfrak{h}}_{\mathbb{C}}^*$.
In this double orbit there is a unique
$W$-orbit in ${\mathfrak{h}}_{\mathbb{C}}^*$ of smallest norm, which we denote by
$\nu_\pi$. If $\chi\in {\mathcal{X}}$ is the class of the orbit of
$(M,\pi)$, $\pi$ an
irreducible cuspidal representation of $M$ then let
$\nu_\chi =\nu_\pi$.

\begin{theorem}\label{thm:6.22} (\cite[Corollary 4.5]{A10}) 
\[
I_\chi
(f)
=\sum_{M\in L(G)}|W_0^M|\cdot |W_0^G|^{-1}
\int_{\Pi (M,\chi)}a^M(\pi
)I_M(\pi ,f)d\pi ,
\]
where 
\[
\Pi (M,\chi )=
\{\pi \in \Pi (\overline{M({\mathbb{A}}_S)^1})\;|\;v_\pi =\nu_\chi \},
\]
$a^M(\pi )$ satisfies \cite[(4.5)]{A10} and $I_\chi$ is
defined by \eqref{eq:6.17}.
\end{theorem}

\subsection{The invariant trace formula}

The geometric side of the trace formula is the sum of the following terms:

(a) the sum over classes ${\mathcal{O}}\subset G({\mathbb{Q}})$
such that ${\mathcal{O}}\cap Z(G({\mathbb{Q}}))\not= \emptyset$ of
\begin{equation}\label{eq:6.24}
I_{\mathcal{O}}(f)=\sum_{\gamma\in {\mathcal{O}}_{(G,S)}}
a^G(S,\gamma)I_G(\gamma,f)+{1\over 2}a^A(S,z)I_A(z,f),
\end{equation}
where $I_M$ and $a^M$ are given in \eqref{eq:6.19} and \S\ref{ssec:explicit},
resp., and the $(G,S)$-classes wre described in Example~\ref{ex:6.12}
above (these include the ``nilpotent terms''),

(b) the sum over elliptic regular classes 
${\mathcal{O}}\subset G({\mathbb{Q}})$
of
\begin{equation}\label{eq:6.25}
I_{\mathcal{O}}(f)=
a^G(S,\gamma)J_G(\gamma,f),
\end{equation}
where $\gamma\in {\mathcal{O}}$ and
$J_G$ (an orbital integral) and $a^G$ (a volume)
are given in \S\ref{ssec:explicit},

(c) the sum over ``hyperbolic'' regular classes
${\mathcal{O}}\subset G({\mathbb{Q}})$
of
\begin{equation}\label{eq:6.26}
I_{\mathcal{O}}(f)=
a^A(S,\gamma)I_A(\gamma,f),
\end{equation}
where $I_A$ and $a^A$ (a volume)
are given in above and in \S\ref{ssec:explicit}.

The spectral side of the trace formula is the sum of the following terms:

(a) the sum over classes $\chi\in {\mathcal{X}}(G)$ of
\begin{equation}\label{eq:6.27}
I_\chi(f)=
{\rm tr}(\pi(f)),\ \ \ (\pi,G)\in\chi,
\end{equation}

(b) the sum over classes $\chi\in {\mathcal{X}}-{\mathcal{X}}(G)$ of
\begin{equation}\label{eq:6.28}
I_\chi(f)=
{1\over 2}\int_{\Pi_u(A,\chi)}
a^A(\pi)I_A(\pi,f)d\pi,
\end{equation}
where $I_A$ is given above and $\Pi_u(A,\chi)$ is as in Theorem~\ref{thm:6.22}
above. (Recall that $I_A(\pi,f)=0$ for $\pi$ {\em tempered}.)





\section{The global trace formula in stable(r) form}\label{sec:stable}

Our primary references for this section are \cite{LL}, \cite{L3}.
\subsection{Definitions and a general discussion of the results}

We define $G_1=GL(2)$. 

Suppose that $\gamma ,\gamma ^{\prime }\in G(\qqq )$. Following the
well-known definition, if there is a $g\in G(\overline{\qqq })$ such that $%
\gamma ^{\prime }=g^{-1}\gamma g$ then we call $\gamma ,\gamma ^{\prime }$ 
\textbf{stably conjugate}. Here $\overline{\qqq }$ denotes an algebraic
closure of $\qqq $. As is well-known, two elements $\gamma ,\gamma
^{\prime }\in G(\qqq )$ are conjugate under
$G(\overline{\qqq })$ if and
only if they are conjugate under $G_1(\qqq )$. See \cite{L3}, \cite{LL}.

Here is reformulation of this definition: If there is a
$g\in G_1(\qqq )\hookrightarrow G_1(\aaa _S)$ such
that $\gamma ^{\prime}=g^{-1}\gamma g$ then
$\gamma ,\gamma ^{\prime }$ are \textbf{stably conjugate}.
This condition has the following local analog which we make our
local definition:
Suppose that $\gamma ,\gamma ^{\prime }\in \overline{G(F)}$. If there
is a $g\in G_1(F\mathbb{)}$ such that 
$\gamma ^{\prime }=g^{-1}\gamma g$ then $\gamma ,\gamma ^{\prime }$ are 
\textbf{stably conjugate}.

\begin{remark}\label{rem:7.1}
We shall mention briefly that the paper of Springer-Steinberg
contains a comment (actually an exercise)
which may be worth mentioning in connection with
these definitions: \cite[ch IV, (2.15)]{SpSt} says that if $E$ is
algebraically closed and $G\subset GL(n)$ is either a symplectic
group or an orthogonal group then then $g_1,g_2\in G(E)$ are
conjugate in $G(E)$ if and only if they are conjugate in $GL(n,E)$.
(I believe it is assumed throughout that $E$ is not of characteristic $2$.)
In particular, if $F$ is a p-adic field, $\overline{F}$ is its
algebraic closure and if $g_1,g_2\in G(F)$ are conjugate in
$GL(n,F)$ then they are conjugate in $G(\overline{F})$.
Recall the fact that if $g_1,g_2\in GL(n,F)$
are conjugate in $GL(n,\overline{F})$
then they are conjugate in $GL(n,{F})$ (``each conjugacy class
in $GL(n)$ is stable'', i.e., ``$GL(n)$ has no L-indistinguishability'').
These in conjunction imply that stable conjugacy in the
symplectic or orthogonal groups may be ``realized'' as
$GL(n,F)$-conjugacy.
\end{remark}

The global invariant trace formula is an equality of the form \eqref{eq:6.4}.
We will say that a distribution $T$ on ${\mathcal{H}}(\overline{G(\aaa _S)})$ is 
\textbf{stably invariant} in case $T(f)=T(f^y)$ for all
$f\in {\mathcal{H}}(\overline{G(\aaa _S)})$ and all
$y\in G_1(\aaa _S)$. Not all the expressions
in the invariant trace formula \eqref{eq:6.4} are stably invariant.
The object
of this section is to write the invariant trace formula
as the sum of stable
terms which we write in the form (as in \cite{L3})

\begin{equation}\label{eq:7.1}
\begin{array}{c}
{\rm tr}R_{cusp}(f)=
\sum_{\chi \in {\mathcal{X}}(G)}I'_\chi (f)\\
=-\sum_{\chi \in {\mathcal{X}}-{\mathcal{X}}(G)}I_\chi (f)
+S(f)
+\sum_{\mathcal{O}\in \mathcal{C}}SI_{\mathcal{O}}(f)
+{\rm\ ``endoscopic\ terms''}
\end{array}
\end{equation}
where $S=S_1+S_2+S_3$ is defined in \eqref{eq:7.22}--\eqref{eq:7.24},
$I'_\chi$ is equal to $I_\chi$ plus a term described in
\S\ref{ssec:stabletf} (g) below ($\chi\in
{\mathcal{X}}-{\mathcal{X}}(G)$), and
the endoscopic terms are a sum of terms of the form

\begin{equation}\label{eq:7.2}
\begin{array}{c}
{\rm\ ``endoscopic\ terms''}
=
{1\over 4}
\sum_{ H\not= G, {\rm cuspidal}}
\left[
\sum_{{\mathcal{O}}\in {\mathcal{C}}^H}SI^H_{\mathcal{O}}(f^H)
-\sum_{\theta\in E_H^G}<\theta,f^H>
\right]\\
={1\over 4}
\sum_{ H\not= G, {\rm cuspidal}}
\sum_{
\theta\in \Pi_u(H(\qqq )\backslash H(\aaa _S)),
\theta\not= 1}
<\theta,f^H>
\end{array}
\end{equation}
which are stable distributions on endoscopic groups $H$ of $G$
(not $\overline{G(\aaa _S)}$)
arising from the stable trace formula of $H$.
Here ${\mathcal{C}}^H=\{\{\gamma\}\ |\ \gamma\in H(\qqq )\}$,
$E_H^G$ is the subset of
$\Pi_u(H(\qqq )\backslash {H(\aaa _S)})$,
and $f^H$ is a function on
$H(\aaa _S)$ associated to
$f\in {\mathcal{H}}(\overline{G(\aaa _S)})_{gen}$.
In the version presented below we will not
stabilize the $I_\chi$ in this paper.
We shall explain the terms ``$H$ a cuspidal endoscopic group of
$G$'' in \eqref{eq:7.2} in the next subsection.

\subsection{Some basics on endoscopic groups}

We shall define what an endoscopic group is in the simple case
of $G=SL(2)$ and discuss the sets $D(T/F)$, $D(T/{\mathbb Q})$,
and $D(T/{\mathbb A})$. The important but more technical distinction
between an endoscopic group and the sharper notion of
endoscopic data (\cite{L3}, \cite{She}) will be glossed over.
Our main reference is \cite{LL}.

In the case of $SL(2)$, a group $H$ is an endoscopic group
if $H$ is a Cartan subgroup of $G$. Moreover, up to ``equivalence''
these are the only ones. In general an endoscopic group $H$
is not a subgroup of $G$ nor are they in general abelian. A
cuspidal endoscopic group of $SL(2)$ over $F$
is an elliptic Cartan (i.e., not conjugate to $A(\qqq )$ by an
element of $G(\qqq )$) of $SL(2)$. 

If $T$ is an elliptic Cartan of $SL(2)$ defined over
a local or global field $F$ then let
\begin{equation}\label{eq:7.3}
D(T/F)=T_1(F)\backslash G_1(F)/G(F),
\end{equation}
where $G_1=GL(2)$ and $T_1=Cent(T,G_1)$.
We identify without further mention this double coset space with the
space
$\overline{T_1(F)}\backslash \overline{G_1(F)}/ \overline{G(F)}$,
where $\overline{G_1(F)}$ denotes the n-fold cover defined using the
cocyle $\beta$ in \S\ref{ssec:local}.
If $T$ is an elliptic Cartan of $SL(2)$ defined over
${\mathbb Q}$ then let
\begin{equation}\label{eq:7.4}
D(T/\aaa )
=\oplus_p T_1(\qqq _p)\backslash G_1(\qqq _p)/G(\qqq _p).
\end{equation}
We embed $D(T/\qqq )$ into $D(T/\aaa )$
or $D(T/\aaa _S)$ diagonally.

For general definitions, see \cite{L3}.
The expressions $SI_{{\mathcal{O}}_H}$ and $f^H$ are discussed there but
will also be mentioned (indirectly) below.

\subsection{Stabilizing the elliptic regular terms}

We shall call an element
$\gamma\in G(\qqq )$ (resp.,
$(\gamma,\zeta) \in \overline{G(F)}$) \textbf{semisimple} if $\gamma $ is
diagonalizable and {\bf regular} if its eigenvalues
$\gamma_1,\gamma_2$ are distinct.
We shall call a semisimple element
$\gamma \in G(\qqq )$
(resp., $(\gamma,\zeta) \in \overline{G(F)}$) \textbf{elliptic}
(or $G$-\textbf{elliptic}) if $\gamma $ belongs to
an elliptic Cartan.

We first shall recall the idea of stabilizing the elliptic regular terms of
\eqref{eq:7.1} following \cite{LL}, \cite{L3}, \cite{Ko1}.

For $\gamma \in G(\qqq )$, let
${\mathcal{O}}_\gamma \in {\mathcal{C}}$ denote
the equivalence class containing $\gamma $. 

We now turn to the trace formula in the rank $0$ case, as discussed in
Section~\ref{sec:rank0} above.

\begin{proposition}\label{prop:7.2} (Poisson summation)
Let $T\subset G$ be an elliptic Cartan defined over $\qqq $. For any
$\phi \in C_c^\infty (\overline{T(\aaa _S)})_{gen}$, we have 
\begin{eqnarray*}
&&{\mathrm{vol}}(T(\qqq )\backslash \overline{T(\aaa _S)})
\sum_{\gamma \in T(\qqq )}\phi (\gamma ) \\
&=&\sum_{\theta \in \Pi _u(T(\qqq )\backslash \overline{T(\aaa _S)})}
<\theta,\phi>.
\end{eqnarray*}
If $\theta =\otimes _p\theta _p$ and $\phi=\otimes _p \phi_p$ then 
\[
<\theta ,\phi>=\prod_p{\rm tr}\theta _p(\phi_p).
\]
\end{proposition}

\begin{remark}\label{rem:7.3}
({\bf A preview of coming attractions})
We shall eventually apply this with $\phi (\gamma )$ equal to

\begin{equation}\label{eq:7.5}
I_\gamma ^\kappa (f)=
[D(T/\qqq ): D(T/\aaa _S)]^{-1}
\sum_{\delta\in
D(T/\qqq )\backslash D(T/\aaa _S)}
\kappa (\det \delta)I_{\mathcal{O}}(f^\delta),
\end{equation}
where $\kappa \in \Pi _u(D(T/\qqq )\backslash D(T/\aaa _S))$,
$\gamma\in T(\qqq )$ elliptic, and
$f\in {\mathcal H}(\overline{G(\aaa _S)})$.
These terms \eqref{eq:7.5} do not occur in the invariant trace formula.
However, after grouping terms as in \eqref{eq:7.6} below and substituting
the expression in Lemma~\ref{lem:7.5} below, all the terms \eqref{eq:7.5},
$\gamma\in T(\qqq )$, will occur in the trace formula
except for $\gamma =\pm 1$. The terms corresponding to
$\gamma =\pm 1$ will be discussed in \S\ref{ssec:singular} below.
The terms \eqref{eq:7.5} with $\kappa = 1$, $\gamma\not= \pm 1$,
will comprise part of the sum over
${\mathcal{O}}\in {\mathcal{C}}$ in \eqref{eq:7.1}. The
terms \eqref{eq:7.5} with $\kappa \not= \pm 1$,
will comprise part of the sum in \eqref{eq:7.2}.  
\end{remark}

Next we record a corollary of the above proposition 
which we will substitute into the trace formula later.

\begin{corollary}\label{cor:7.4}
For each elliptic torus T, we have 
\[
vol[T(\qqq )\setminus \overline{T(\aaa _S)}]\cdot
\sum_{\gamma \in T({\mathbb Q})}
I_\gamma ^\kappa (f)=\sum_\theta <\theta ,I_{*}^\kappa (f)>
\]
where $I_{*}^\kappa $ is the unstable orbital integral defined in \eqref{eq:7.5}
above (for its properties, as a function on $T$, in particular on the
singular set of $T$, see \cite{LL} or \S\ref{ssec:singular} below).
\end{corollary}

This expression will form part of what is meant by \eqref{eq:7.2}.

As we will see in Lemma~\ref{lem:7.21} below,
we may use \cite{LL} to to relate $<\theta,I_{*}^\kappa (f)>$ to
a linear combination of characters of automorphic representations of
$\Pi_u(\overline{G(\aaa _S)})$. If $\mathcal{O}=\mathcal{O}_\gamma $
and $\gamma \in \overline{T(\aaa _S)}$ is elliptic regular, let 
\begin{equation}\label{eq:7.6}
\begin{array}{c}
I^{st}_{\mathcal{O}}(f)=\sum_{\delta\in D(T/\qqq )}
I_{\mathcal{O}}(f^\delta),\\
SI_{\mathcal{O}}(f)=I^1_{\gamma}(f)\; .
\end{array}
\end{equation}
These distributions $SI_{\mathcal{O}}$ are stable.

\begin{lemma}\label{lem:7.5}
(a)\ If $\mathcal{O}=\mathcal{O}_\gamma $ and
$\gamma \in \overline{T(\aaa _S)}$ is
elliptic regular, we have 
\[
SI_{\mathcal{O}}(f)
=\sum_{\kappa \in \Pi _u(D(T/{\mathbb Q})\backslash D(T/\aaa ))}
I_\gamma ^\kappa (f)
\]
(b) If $\mathcal{O}=\mathcal{O}_\gamma $ and
$\gamma \in \overline{T(\aaa _S)}$ is elliptic regular, we have 
\[
I_{h^{-1}\gamma h}^\kappa (f)=\kappa (\det h)I_\gamma ^\kappa (f)=I_\gamma
^\kappa (f^h),\ \ \ \ \ \ h\in \overline{G_1(\aaa _S)}.
\]
\end{lemma}

Any function $g$ satisfying $g(h^{-1}\gamma h)=\kappa (\det (h))g(\gamma )$,
for all $h\in \overline{G_1(\aaa _S)}$, will be called
\textbf{$\kappa $-stable}.

proof: It is elementary that if $A$ and $B$ are
discrete countable abelian groups, with $A\subset B$ of finite
index, and if $f$ is any compactly supported function on $B$ then 
\begin{equation}\label{eq:7.7}
\sum_{a\in A}f(a)=[B:A]^{-1}\sum_{\kappa \in (B/A)^{\wedge }}\sum_{b\in
B}\kappa (b)f(b),
\end{equation}
whenever the sums converge absolutely. We apply this to our sum over $\delta 
$: 
\[
\sum_{\delta \in D(T/\qqq )}
I_{\mathcal{O}}(f^\delta)
=
\sum_{\kappa \in (D(T/\aaa )/D(T/\qqq ))^{\wedge}}
[D(T/\aaa ):D(T/\qqq )]^{-1}
\sum_{\delta \in D(T/\aaa )}
\kappa (\delta )I_{\mathcal{O}}(f^\delta),
\]
which we may equate with
\begin{equation}\label{eq:7.8}
\sum_{\kappa \in \Pi_u(D(T/\aaa )/D(T/\qqq ))}
I^\kappa_\gamma (f).
\end{equation}
$\Box$

\subsection{Stabilizing the $A$-elliptic regular terms}

We shall call a semisimple element
$\gamma \in G(\qqq )$
(resp., $\gamma \in \overline{G(F)}$) $A$-\textbf{elliptic} regular if
$\gamma $ is conjugate to an element $a\in A(F)$
which is elliptic regular in $A(F)$ (i.e., $a\not= \pm 1$).

If $\mathcal{O}=\mathcal{O}_\gamma $ and
$\gamma \in G(\qqq )$ is $A$-elliptic regular, let 
\begin{equation}\label{eq:7.9}
\begin{array}{c}
I^{st}_{\mathcal{O}}(f)=
\sum_{\delta\in D(A/\qqq )}
I_{\mathcal{O}} (f^\delta),\\
I^\kappa_\gamma(f)
=[D(A/\qqq ):D(A/\aaa _S)]^{-1}
\sum_{\delta\in D(A/\aaa _S)}
\kappa(\delta) I_{\mathcal{O}} (f^\delta),
\end{array}
\end{equation}
where
\[
\begin{array}{c}
D(A/\qqq )
=A(\qqq )Z(G_1(\qqq ))\backslash G_1(\qqq )/G(\qqq ) \\
D(A/\aaa _S)
=A(\aaa _S)Z(\overline{G_1(\aaa _S)})
\backslash \overline{G_1(\aaa _S)}/G(\aaa _S) \\
\kappa\in \Pi_u(D(A/\qqq )\backslash D(A/\aaa _S)).
\end{array}
\]
Let
\[
SI_{\mathcal{O}}(f)=I^1_{\gamma} (f),
\]
From \eqref{eq:7.7}, we have
\begin{equation}\label{eq:7.10}
I^{st}_{\mathcal{O}}(f)=
\sum_{\kappa\in \Pi_u(D(A/\qqq )\backslash D(A/\aaa _S))}
I^\kappa_\gamma(f)
\end{equation}
\begin{lemma}\label{lem:7.6}
If $\mathcal{O}=\mathcal{O}_\gamma $ and
$\gamma \in G(\qqq )$ is $A$-elliptic regular, then
$SI_{\mathcal{O}}(f)$ is stable.
\end{lemma}

Recall that this term $SI_{\mathcal{O}}$ is defined using
\begin{equation}\label{eq:7.11}
I_{\mathcal{O}}(f)
=\sum_{\gamma\in {\mathcal{O}}}a^G(S,\gamma)I_G(\gamma,f)
+{1\over 2}\sum_{\gamma\in {\mathcal{O}}\cap A(\qqq )}
a^A(S,\gamma)I_A(\gamma,f).
\end{equation}
The $I_M(\gamma,f)$ were defined by
\[
I_G(\gamma,f)=J_G(\gamma,f),
\]
which is an ``ordinary'' orbital integral over
$\overline{G(\aaa _S)}_\gamma\backslash \overline{G(\aaa _S)}$, and
\[
I_A(\gamma,f)=J_A(\gamma,f)-\hat{I}_A^A(\gamma,\phi_A(f)).
\]
(It turns out that
$\overline{G(\aaa _S)}_\gamma\subset \overline{A(\aaa _S)}$,
with equality if $\gamma$ is an $N$-th power, and
$J_G(\gamma,f)=0$ if $\gamma$ is not an $N$-th power - provided
$f$ is genuine as we have assumed.)
To see that the sum
$\sum_{\gamma\in {\mathcal{O}}}a^G(S,\gamma)I_G(\gamma,f)$
is stably invariant, we claim that $I_G(\gamma,f)$ is stably invariant.
This is proven in \cite{J3}.

The term $J_A(\gamma,f)$ is not invariant. We examine it
more closely in Lemma~\ref{lem:7.9} and \eqref{eq:7.14} below.
(See also \cite[Corollary 4.2]{A16} for a general formula for it
in the case $n=1$.)

The term $\hat{I}_A^A(\gamma,\phi_A(f))$ is not invariant. It is the
Fourier transform of the map
\[
\chi\rightarrow {\rm tr}(I_A(\chi,f)R_A(\chi,B)),
\]
see \eqref{eq:7.18}. The instability of $f\longmapsto \phi_A(f)$ is
examined in Lemma~\ref{lem:7.23} below.

Let $F=\qqq _p$, $p<\infty$, $K=G({\mathbb{Z}}_p)$ which we
identify with a {\em subset} $(K,1)$ of $\overline{G(F)}$,
$\sigma\in \Pi(\overline{A(F)})$. Let $K_{I_A(\sigma)}(k,k')$
denote the kernel function of the induced representation
$I_A(\sigma)$ on the space of functions on $K$, so
\begin{equation}\label{eq:7.12}
{\rm tr}(I_A(\sigma,f))
=\int_K K_{I_A(\sigma)}(k,k)dk .
\end{equation}

\begin{lemma}\label{lem:7.7}
If $F(\sigma)={\rm tr}(I_A(\sigma,f))$ then
\[
F^\vee(\gamma)
=|a_1-a_2|J_G(\gamma,f),
\]
for $\gamma =(a,\zeta)\in \overline{A(F)}$, $a\not= \pm 1$.
\end{lemma}

proof: This will be proven in \cite{J3}.

As a corollary, it follows from this lemma without to much work
that $f\longmapsto I_G(\gamma,f)$ is stable, for
$\gamma =(a,\zeta)\in \overline{A(F)},\ a\not= \pm 1$.

Let $\delta(g)=|a_1/a_2|$, if $g=(nak,\zeta)\in \overline{G_1(F)}$. Let
$N(g)$ denote the operator on the space of functions on $K$ given by
\[
f(k)\longmapsto f(k)\log\delta(kg).
\]
This operator also occurs in connection with the terms in the
invariant trace formula associated to the normalized intertwining
operators.

\begin{lemma}\label{lem:7.8}
For $g\in G_1(F)$, we have
\[
{\rm tr}(I_A(\sigma,f)N(g))
=\int_K K_{I_A(\sigma)}(k,k)\log\delta(kg)dk.
\]
\end{lemma}

\begin{lemma}\label{lem:7.9}
For $g\in G_1(F)$, if
$F(\sigma)={\rm tr}(I_A(\sigma,f)N(g))$ then
\[
J_A(a,f^{g^{-1}})-J_A(a,f)
=F^\vee(a)+F^\vee(a^w)
\]
\end{lemma}
proof: This will be proven in \cite{J3}.

If $v>2$ then we may represent $D(A/F)$ by
\[
\{
\left(
\begin{array}{cc}
a&0\\
0&1
\end{array}
\right)\ |\ a=1,\epsilon, \pi, \epsilon\pi \},
\]
where $\pi$ is a uniformizer and $\epsilon$ is a unit
which is not a quadratic residue. If $v=2$ the situation is similar...
Therefore, if $\kappa$ is non-trivial
we may apply Lemma~\ref{lem:7.9} above to compute
\[
\sum_{\delta\in D(A/F)}
\kappa(\delta)J_A(\gamma,f^\delta).
\]

If $x=a
\left(
\begin{array}{cc}
1&n\\
0&1
\end{array}
\right) k$, $n\in \aaa _S$, then
\begin{equation}\label{eq:7.13}
v_A(x)
=v(x,0)=\sum_v v(0,x_v)=
\sum_v \log\sqrt{|1+n_v^2|_v}
\end{equation}
and
\begin{equation}\label{eq:7.14}
\begin{array}{c}
J_A(\gamma,f)
=|D(\gamma)|^{1/2}
\int_{\overline{G(\aaa _S)}_\gamma
\backslash \overline{G(\aaa _S)}}
f(x^{-1}\gamma x)v_A(x)dx\\
=
\sum_w\prod_v
|D(\gamma_v)|_v^{1/2}
\int_{\overline{G(\qqq _p)}_\gamma
\backslash \overline{G(\qqq _p)}}
f_v(x_v^{-1}\gamma_v x_v)v_A(x_w)dx_v\\
=\sum_v |D(\gamma_v)|_v^{1/2} J_A(\gamma_v,f_v)
\prod_{w\not= v}|D(\gamma_w)|_w^{1/2} J_G(\gamma_w,f_w).
\end{array}
\end{equation}

\begin{remark}\label{rem:7.10} This comprises metaplectic analog of
the $a\not= \pm 1$
part of the distribution denoted $S_2$ in \cite{LL}. It
will be combined with other distribution
described in the next section to yield a stably
invariant distribution.
\end{remark}

\subsection{Stabilizing the singular terms}\label{ssec:singular}

A term $I_{\mathcal{O}}(f)$ is called {\bf singular} if ${\mathcal{O}}$
does not contain any $M$-elliptic elements, for any $M\in \{A,G\}$.
An element $\gamma\in G(\qqq )$ is called {\bf singular}
if its eigenvalues $\gamma_1,\gamma_2$ are equal (i.e.,
$\gamma$ is not regular).

In this subsection, our goal is to show that part of
the singular terms of the geometric side of the invariant
trace formula \eqref{eq:6.4} ``combine'' with the sum of the terms
$I_\gamma ^\kappa (f)$,
with $\kappa \neq 1$ and with
$\gamma \in T(\qqq )$ elliptic \emph{singular}
(i.e., $\gamma =\pm 1$) to contribute to \eqref{eq:7.2}. 
If ${\mathcal O}=\{z\}$, $z=\pm 1$, then define
\begin{equation}\label{eq:7.15}
SI_{\mathcal O}(f)=c\cdot f(z),
\end{equation}
where $c=vol(G({\mathbb Q})\backslash \overline{G(\aaa _S)})$ is a volume.

\begin{lemma}\label{lem:7.11} If $z\in {\mathcal{O}}$, with $z=\pm 1$, then
\[
I_{\mathcal O}(f)-SI_{\mathcal O}(f)=
\sum_{ \gamma\in {\mathcal{O}}, \gamma\not= z}
a^G(S,\gamma)I_G(\gamma,f)+
{1\over 2}a^A(S,z)I_A(z,f).
\]
\end{lemma}
proof: This equation follows from Proposition~\ref{prop:6.17}
(also Proposition~\ref{prop:4.15}, Example~\ref{ex:4.16}). $\Box$

\begin{definition}\label{defn:7.12}
Let $\kappa$ be a character of
$D(T/F)$ of order $2$.
Define $\Phi^{T/\kappa}$ on the regular set of $\gamma\in T(F)$ by
\[
\begin{array}{c}
\Phi^{T/\kappa}(\gamma,f)=
\Delta(\gamma,T,\kappa)
\int_{\overline{G_1(F)}_\gamma\backslash \overline{G_1(F)}}
\kappa(\det{g})f(g^{-1}\gamma g)dg\\
=
\lambda(E/F,\psi)
\kappa({\gamma_{1,v}-\gamma_{2,v}\over \gamma_{1,v}^0-\gamma_{2,v}^0})
\int_{\mathcal{O}}\epsilon(x)f(x)|\omega_\gamma|,
\end{array}
\]
where $\epsilon(g^{-1}\gamma g)=\kappa(\det(g))$,
$\gamma^0\in \overline{T(F)}$
is fixed, and $\mathcal{O}$ denotes the conjugacy class of
$\gamma$ in $\overline{G(F)}$.
\end{definition}

We leave the relation between the global expressions
$\Phi^{T/\kappa}(\gamma,f)$ and $I^\kappa_\gamma(f)$ to
the reader (use the Hilbert product formula).

Let ${\mathcal{O}}(\gamma)$ denote the orbit of
$\gamma$ in $\overline{G(F)}$ and,
for $a\in F^\times$, which we identify with the center of $G_1(F)$, let
\[
n(a)=a
\left(
\begin{array}{cc}
1&1\\ 0&1
\end{array}
\right).
\]

\begin{definition}\label{defn:7.13} If $z\in Z(\overline{G(F)})$, then we define
\[
\Phi^{T/\kappa}(z,f)
=
\left\{
\begin{array}{c}
\int_{{\mathcal{O}}(n(z))}
f(x)dx,\ \ \ \ T=A,\ \kappa =1\\
\lambda(E/F,\psi)
\kappa({\pi_E-\overline{\pi_E}\over \gamma_{1,v}^0-\gamma_{2,v}^0})
\int_{{\mathcal{O}}(n(z))}\epsilon(x)f(x)dx,\ \ \ \ T_1(F)=E^\times,
\kappa\in \Pi_u(F^\times),
\end{array}
\right.
\]
where $\pi_E$ is a uniformizer for $E/F$, $T_1(F)=E^{\times 1}$,
and
\[
\epsilon(g^{-1}n(z)g)=
\left\{
\begin{array}{c}
\kappa(-b),\ \ \ \ b\not= 0,\\
\kappa(c),\ \ \ \ b= 0,\\
\end{array}
\right.
\]
if
\[
g^{-1}n(z)g=
\left(
\begin{array}{cc}
a&b\\c&d
\end{array}
\right).
\]
\end{definition}

\begin{lemma}\label{lem:7.14} If $\kappa\in \Pi_u(F^\times)$ is the character
of order 2
associated to $E/F$ be local class field theory ($\kappa =1$
if $E=F\times F$) and $T_1(F)=E^\times$, then
$\gamma\longmapsto \Phi^{T/\kappa}(\gamma,f)$ is a smooth function on
$T(F)$.
\end{lemma}

proof: This follows from the results on Shalika germs. It is proven
in \cite{J3}.

\begin{lemma}\label{lem:7.15} For $z\in Z(\overline{G(F)})$, we have
\[
\begin{array}{c}
\lim_{\gamma\rightarrow z}
\int_{\overline{G_1(F)}_\gamma\backslash \overline{G_1(F)}}
\kappa(\det{g})f(g^{-1}\gamma g)dg\\
c{1\over L(1,\kappa)}
\int_{\overline{G_1(F)}_{n(1)}\backslash \overline{G_1(F)}}
\kappa(\det{g})f(g^{-1}n(z) g)dg
\end{array}
\]
where $c=|\gamma_{1,v}-\gamma_{2,v}|$.
\end{lemma} 

proof: This will be proven in \cite{J3}.

\begin{lemma}\label{lem:7.16} Let $\kappa$ be a character of
$\aaa _S^\times $ of order 2 which we extend to
$\overline{G_1(\aaa _S)}$ via the $\det$ map,
and $T$ be a Cartan of $G$ associated to $\kappa$ by class field
theory then
\begin{equation}\label{eq:7.16}
\sum_{\delta\in D(A/\aaa _S)}
\kappa(\delta)I_G(n(z),f^\delta)
=c
\lim_{ \gamma\rightarrow z , \gamma\in T(\qqq )}
\sum_{\delta\in D(T/\aaa _S)}
\kappa(\delta)I_G (\gamma, f^\delta),
\end{equation}
for $z =\pm 1$, $\kappa \in \Pi _u(G_1(\aaa _S))$ of order 2,
where the constant $c=c(T,\kappa,S)$ is given by
\[
c=\prod_v |\gamma_{1,v}^0-\gamma_{2,v}^0|_v^{-1}L(1,\kappa_v) .
\]
\end{lemma}

\begin{remark}\label{rem:7.17} The terms on the right hand side of the above
lemma go into \eqref{eq:7.20} below. To see this, just relate them to
\eqref{eq:7.5} by using \eqref{eq:7.11} and the fact that
${\mathcal{O}}\cap A(\qqq ) =\emptyset$.

The terms on the left hand side will examined following the proof
below.
\end{remark}

proof: This is is a consequence of the previous lemma and the
definitions. $\Box$

If $z\in {\mathcal{O}}$ is an equivalence class of $G(\qqq )$
as above then part of $I_{\mathcal{O}}(f)$ may be written as
\begin{equation}\label{eq:7.17}
\begin{array}{c}
\sum_{\gamma\in {\mathcal{O}}} J_G(\gamma,f)
=J_G(z,f)+\sum_{a\in \qqq ^{\times}/\qqq ^{\times 2}}
\int_{\overline{G(\aaa _S)}_{n(1)}
\backslash \overline{G(\aaa _S)}} f(g^{-1} z\cdot
\left(
\begin{array}{cc}
1 & a\\
0 & 1
\end{array}
\right)
g)dg\\
=J_G(z,f)+
\sum_{\kappa \in \Pi_u(\aaa _S^{\times}/\aaa _S^{\times 2})}
\sum_{a \in \aaa _S^{\times}/\aaa _S^{\times 2}}
\kappa (a)
\int_{\overline{G(\aaa _S)}_{n(1)}
\backslash \overline{G(\aaa _S)}} f(g^{-1} z\cdot
\left(
\begin{array}{cc}
1 & a\\
0 & 1
\end{array}
\right)
g)dg\\
=J_G(z,f)+
\sum_{\kappa \in \Pi_u(\aaa _S^{\times}/\aaa _S^{\times 2})}
\int_{\overline{G_1(\aaa _S)}_{n(1)}
\backslash \overline{G_1(\aaa _S)}} \kappa(\det{g})
f(g^{-1} n(z) g)dg,
\end{array}
\end{equation}
which we write as $J_G(z,f)$ plus the sum over
$\kappa \in \Pi_u(\aaa _S^{\times}/\aaa _S^{\times 2})$
of the left hand side of the equation in Lemma~\ref{lem:7.16} above.

Regarding the expression
\[
{a^A(S,z)\over 2}
\sum_\delta \kappa(\delta)
(J_A(z,f^\delta)-\hat{I}_A^A(z,\phi_A(f^\delta)))),
\]
we must also compute $a^A(S,z)$ and discuss the terms
$\sum_\delta \kappa(\delta)J_A(z,\phi_A(f^\delta))$ and
$\sum_\delta \kappa(\delta)\hat{I}_A^A(z,\phi_A(f^\delta))$.

Computation of $a^A(S,z)$: By Proposition~\ref{prop:4.15},
\[
a^A(S,z)=vol(A(\qqq )\backslash \overline{A(\aaa _S)^1}).
\]

Computation of
$\sum_\delta \kappa(\delta)J_A(z,\phi_A(f^\delta))$:
See Example~\ref{ex:4.13} above and \eqref{eq:7.14}. These terms
comprise the $a=\pm 1$ part of the distribution denoted by
$S_2(f)$ ``missing'' from \eqref{eq:7.14}.

Computation of $\sum_\delta \kappa(\delta)\hat{I}_A^A(z,\phi_A(f^\delta))$:
First, we note that the map
$\phi_A:{\mathcal{H}}(\overline{G(\aaa _S)}) \rightarrow I(A(\aaa _S))$
of \S\ref{sssec:ass51} is a ``weighted
character'' of an induced representation:
\begin{equation}\label{eq:7.18}
\begin{array}{c}
\phi_A(\sigma,f)=
{\rm tr}(R_{\overline{B}|B}(\sigma_\lambda)^{-1}
R'_{\overline{B}|B}(\sigma_\lambda)I_A(\sigma_\lambda,f))|_{\lambda =0}\\
=\sum_v
{\rm tr}(R_{\overline{B}|B}(\sigma_v)^{-1}
R'_{\overline{B}|B}(\sigma_v)I_A(\sigma_v,f_v))
\prod_{w\not= v}{\rm tr}(I_A(\sigma_w,f_w)).
\end{array}
\end{equation}
for $\sigma\in \Pi_u(\overline{A(\aaa _S)})$.
The inverse Fourier transform of the function
$\sigma\longmapsto \phi_A(\sigma, f)$ evaluated at $a=1$ is
essentially the distribution denoted by $S_1(f)$ in \cite{LL}.

\subsection{The geometric side of the stable trace formula}

The stable elliptic regular terms are 
\begin{equation}\label{eq:7.19}
\sum_{ {\mathcal{O}}\in {\mathcal{C}}_G ,
\mathcal{O}\;\mathrm{regular\ elliptic} }
SI_{\mathcal{O}}(f),
\end{equation}
where the $SI$'s are defined by \eqref{eq:7.6}, \eqref{eq:7.9}, 

\begin{definition}\label{defn:7.18} If $\pi$ is an irreducible admissible
genuine representation of
$\overline{G({\mathbb Q}_p)}$ then let
$L(\pi)$ denote the set of inequivalent representations
$\pi^g$ where $g\in \overline{G_1(\qqq _p)}$.
\end{definition}

\begin{remark}\label{rem:7.19} For $\pi$ as above and $n>1$ then
\[
\# L(\pi)=
\prod_{v|p}\#( k_v^\times/k_v^{\times 2}).
\]
\end{remark}

\begin{lemma}\label{lem:7.20} If $\theta_1$ is an irreducible genuine representation
of ${T_1({\mathbb Q}_p)}$ and $\pi_1(\theta_1)$ is the
irreducible admissible representation of
${G_1({\mathbb Q}_p)}$ associated to it in \cite{J2}
(at least if $n$ is odd), let
$\theta =\theta_1|_{{T(\qqq _p)}}$ and let
$\pi(\theta)$ denote one of the irreducible
constituents of
$\pi(\theta) =\pi_1(\theta_1)|_{G(\qqq _p)}$.
\end{lemma}

This is proven in \cite{J3}.

The following lemma will not be used here. Note that it related
the distribution
\[
\begin{array}{c}
C_c^\infty(\overline{G(\aaa _S)})_{gen}
\rightarrow
C_c^\infty(\overline{T(\aaa _S)})_{gen}\\
f\longmapsto I_*^\kappa(f)
\end{array}
\]
to a map
\[
\begin{array}{c}
C_c^\infty(T(\aaa _S))'
\rightarrow
C_c^\infty(\overline{G(\aaa _S)})'_{gen}\\
\theta \longmapsto
\prod_{p\in S}\sum_{\pi _p\in \Pi (\theta_p)}
\kappa _{*}(\pi _p)\chi _{\pi _p}(f_p),
\end{array}
\]
after relating $C_c^\infty(\overline{T(\aaa _S)})_{gen}$
to $C_c^\infty(T(\aaa _S))$.
\begin{lemma}\label{lem:7.21}
For each elliptic torus T and each factorizable
$f\in {\mathcal{H}}(\overline{G(\aaa _S)})_{gen}$, we have 
\[
\begin{array}{c}
<\theta ,I_{*}^\kappa (f )> \\ 
=[D(T/\qqq ):D(T/\aaa _S)]^{-1}
\prod_{p\in S}
\sum_{\delta\in D(T/{\mathbb Q}_p)}
\sum_{\pi _p\in \Pi (\theta_p^\delta)}
\kappa _{*}(\pi _p)\chi _{\pi _p}(f_p),
\end{array}
\]
where $\kappa _{*}(\pi _p):=\kappa (\delta)$ if
$\pi _p\in \Pi (\theta_p^\delta)$,
${\theta }_p^h$ denotes the conjugate of ${\theta }_p$ by $h$,
${\theta }_p$ is a unitary genuine irreducible
representation of $\overline{T(\aaa _S)}$, and
$f_p(g)=f_p(g,1)$ (with apologies for the abuse of
notation).
\end{lemma}

This is proven in \cite{J3}.
\begin{remark}\label{rem:7.22}
This result is essentially the metaplectic analog of
\cite[(2.4-5)]{LL}, \cite[p. 762]{LL}. This is, after a suitable 
choice of global measures, related to a special case of
\cite[(3.42)]{J2} which was only proven under in the case $n$ odd.
\end{remark}

By this lemma and the trace formula for rank 0 (in
Proposition~\ref{prop:7.2})
the terms 
\begin{equation}\label{eq:7.20}
\sum_{ T \; , G-\mathrm{elliptic} }
{\rm vol}(T(\qqq )\backslash \overline{T(\aaa _S)})
\sum_{\gamma \in T(\qqq )}
\sum_{ \kappa \in \Pi _u(D(T/{\mathbb Q})\backslash D(T/\aaa _S)),
 \kappa \neq 1}
I_\gamma ^\kappa (f)
\end{equation}
may be expressed as a linear combination of characters of representations 
\begin{equation}\label{eq:7.21}
\sum_T
\sum_\theta c_\theta tr(\pi (\theta )(f)),
\end{equation}
for some constants $c_\theta$, where $\theta$ runs over
$\Pi_u(T(\qqq )\backslash \overline{T(\aaa _S)})$.
Note that we may rewrite the terms in \eqref{eq:7.20} as
\[
\sum_{H\ cuspidal}
\sum_{ {\mathcal O}_H\not= \pm 1,
a\ conj,
class\ in\ H}
SI_{{\mathcal O}_H}^H(f^H)
\]
in the notation of \eqref{eq:7.2}.

In \eqref{eq:7.20} the terms $\gamma =\pm 1$ occuring
in the middle sum on the right
come from the singular elliptic terms in Lemma~\ref{lem:7.16} above.

\subsection{Stabilizing the remaining terms}

Let $F=\qqq _p$.

\begin{lemma}\label{lem:7.23} For all $g\in \overline{G_1(F)}$,
\[
\begin{array}{c}
{\rm tr}(R_{\overline{B}|B}(\sigma)^{-1}
R'_{\overline{B}|B}(\sigma)I_A(\sigma,f^{g^{-1}}))
-{\rm tr}(R_{\overline{B}|B}(\sigma)^{-1}
R'_{\overline{B}|B}(\sigma)I_A(\sigma,f))\\
={\rm tr}(I_A(\sigma,f)N(g))+{\rm tr}(I_A(\sigma^w,f)N(g))
\end{array}
\]
\end{lemma}

proof: Proven in \cite{J3}.

From Example~\ref{ex:5.28}, \eqref{eq:5.28}, and Theorem~\ref{thm:6.22}, let
\begin{equation}\label{eq:7.22}
S_1(f)=\sum_{\chi\in{\mathcal{X}}-{\mathcal{X}}(G)}
\sum_v {1\over 2}
\int_{\Pi(A,\chi)}
{\rm tr}(R_{\overline{B}|B}(\sigma_v)^{-1}
R'_{\overline{B}|B}(\sigma_v)I_A(\sigma_v,f_v))
\prod_{w\not= v}
{\rm tr}(I_A(\sigma_w,f_w))|d\pi |.
\end{equation}

Let
\begin{equation}\label{eq:7.23}
S_2(f)=
-{1\over 2}
\sum_v \sum_{\{a,a^w\}\in A(\qqq )}
{|D(a)|^{1/2}_v
J_A(a,f_v)\over L(1,\qqq _v)}
\prod_{w\not= v}
{|D(a)|^{1/2}_w J_G(a,f_w)\over L(1,\qqq _w)}.
\end{equation}

Let
\begin{equation}\label{eq:7.24}
\begin{array}{c}
S_3(f)=\\
=Res_{s=0}{1\over s}
\sum_{z\in Z(G(\qqq ))}
\sum_{
x\in \qqq ^\times /\qqq ^{\times 2}
x\not= 1 }
vol(Z(G(\qqq ))\backslash Z(\overline{G(\aaa _S)}))\times \\
\times
\int_{\overline{G(\aaa _S)}_{n(1)}
\backslash \overline{G(\aaa _S)}}
f(g^{-1}z
\left(
\begin{array}{cc}
1&x\\ 0&1
\end{array}
\right)
 g)\lambda(g)^{-s}dg ,
\end{array}
\end{equation}
where $\lambda(g)=\delta(g)\delta(wg)$.

The following will not be needed. It is the analog of a result in
\cite{LL}, the proof of which used the analogs of Lemmas~\ref{lem:7.9}
and~\ref{lem:7.23} above.
\begin{conjecture}\label{conj:7.24}
Let $S=S_1+S_2+S_3$. Then $S$ is stable
$S(f^g)=S(f)$ for all $g\in \overline{G_1(\aaa _S)}$.
\end{conjecture}

Perhaps the proof of this (assuming it i correct)
is similar to the proof of its analog in \cite{LL}.

\subsection{The stable trace formula}\label{ssec:stabletf}

This is the equality 
\begin{equation}\label{eq:7.25}
\sum_{\mathcal{O}\in \mathcal{C}}SI_{\mathcal{O}}(f)
=\sum_{\chi \in {\mathcal{X}}}I_\chi (f)
-\sum_{H\ cuspidal}\sum_{{\mathcal O}_H}SI_{{\mathcal O}_H}^H(f^H),
\end{equation}
where the $SI_{\mathcal{O}}$ are as in sections 7.2-7.3.

The geometric side of the stable(r)
trace formula is the sum of the following terms:

(a) the sum over $(G,S)$-classes ${\mathcal{O}}\subset G(\qqq )$
such that ${\mathcal{O}}\cap Z(G(\qqq ))\not= \emptyset$ of
\begin{equation}\label{eq:7.26}
SI_{\mathcal{O}}(f)=
a^G(S,z)f(z),
\end{equation}
where $a^G$ is (a volume) is given in \S\ref{ssec:explicit} 
(these do not include the ``nilpotent terms'', which
got combined with some other terms ...),

(b) the sum over elliptic regular classes 
${\mathcal{O}}\subset G(\qqq )$
of
\begin{equation}\label{eq:7.27}
SI_{\mathcal{O}}(f)=
\sum_{\delta\in D(T/\aaa _S)}
a^G(S,\gamma)J_G(\gamma,f^\delta),
\end{equation}
where $\gamma\in {\mathcal{O}}$, $T=G_\gamma$, and
$J_G$ (an orbital integral) and $a^G$ (a volume)
are given in \S\ref{ssec:explicit},

(c)  the sum over $A$-elliptic regular classes
${\mathcal{O}}\subset G(\qqq )$
of
\begin{equation}\label{eq:7.28}
SI_{\mathcal{O}}(f)=
\sum_{\delta\in D(A/\aaa _S)}
I_{\mathcal{O}}(f^\delta),
\end{equation}

(d) plus terms associated to the geometric side of the 
trace formula for each cuspidal endoscopic of $G$,

(e) the ``geometric part'' of the endoscopic terms \eqref{eq:7.2}.

The spectral side of the stable(r) metaplectic
trace formula is the sum of the following terms:

(f) the sum over classes $\chi\in {\mathcal{X}}(G)$ of
\begin{equation}\label{eq:7.29}
I_\chi(f)=
{\rm tr}(\pi(f)),\ \ \ (\pi,G)\in\chi.
\end{equation}

(g) The sum over classes $\chi\in {\mathcal{X}}-{\mathcal{X}}(G)$
of
\begin{equation}\label{eq:7.30}
I_\chi(f)+\sum_{ \sigma,
(\sigma,A)\in\chi}
{m'(\sigma)\over m(\sigma)}
{\rm tr}(I_A(\sigma,f)),
\end{equation}

(h) the sum of the distributions $S_1$, $S_2$, and $S_3$

(i) the ``spectral part'' of the endoscopic terms \eqref{eq:7.2}
(which we have omitted ...).

\section{Some food for thought}

It would be nice if we had

(1) provided a natural interpretation of the
stabilization of the terms in (h),

(2) uncovered endoscopic analogs of Arthur's
inductive formulas \eqref{eq:6.19}, \eqref{eq:6.23} of the invariant
distributions $I_{\mathcal{O}}$ and $I_\chi$ (if one exists),

(3) a stabilization of the terms in (f).




\begin{thebibliography}{99}

\bibitem[A1]{A1} J. Arthur,
``The Selberg trace formula for $F$-rank one groups'', Ann.
Math. \underline{100}(1974)326-385

\bibitem[A2]{A2} ------,
``A trace formula for reductive groups, I : terms associated to
classes in $G(\mathbb{Q})$'', Duke Math. J. \underline{45} (1978)911-952

\bibitem[A3]{A3}  ------,
``A trace formula for reductive groups, II : applications of a
truncation operator'', Comp. Math. \underline{40} (1980)87-121

\bibitem[A4]{A4} ------,``Eisenstein series and the trace formula'', in
{\bf Automorphic forms, Representations, and
L-functions, vol. I}, 1979 Corvallis Conference, ed. A. Borel, W. Casselman,
Proc. Symp. Pure Math. vol 33, AMS, 1979 

\bibitem[A5]{A5} ------,
``The trace formula for non-compact quotient'', Proc. Inter.
Cong. Math., Warsaw, Poland, 1983

\bibitem[A6]{A6} ------, ``The trace formula in invariant form'', Ann. Math.
\underline{114}(1981)1-74

\bibitem[A7]{A7} ------, ``The local behaviour of weighted orbital integrals'',
Duke Math. J. \underline{56}(1988)223-293

\bibitem[A8]{A8} ------, ``A measure on the unipotent variety'', Can. J. Math.
\underline{37}(1985)1237-1274

\bibitem[A9]{A9} ------, ``The invariant trace formula, I. Local theory'',
J.A.M.S.\underline{1}(1988)323-383

\bibitem[A10]{A10} ------, ``The invariant trace formula, II. Global theory'',
J.A.M.S.\underline{1}(1988)501-554

\bibitem[A11]{A11} ------, ``Intertwining operators and residues, I: weighted
characters'', J. Func. Anal.\underline{84}(1989)19-84

\bibitem[A12]{A12} ------, ``A family of distributions obtained from
orbits'', Can. J.
Math. \underline{38}(1986)179-214

\bibitem[A13]{A13} ------, ``On a family of distributions obtained from
Eisenstein series, I: application of the Paley--Wiener theorem'', Amer. J.
Math. \underline{104} (1982) 1243--1288

\bibitem[A14]{A14} ------, ``On a family of distributions obtained from
Eisenstein series, II: explicit formulas'', Amer. J. Math.
\underline{104} (1982) 1289--1336

\bibitem[A15]{A15} ------,``The trace formula for reductive groups,'' in
{\bf Journ\'ees Automorphes}, Publ. Math. Univ. Paris VII,
U.E.R. Math, 1983

\bibitem[A16]{A16} ------,``On the Fourier transform of weighted
orbital integrals,'' J. reine angew. Math. \underline{452} (1994) 163--217

\bibitem[BDK]{BDK} J. Bernstein, P. Deligne, D. Kazhdan, ``Trace
Paley-Wiener theorem for
reductive p-adic groups'', J. d'Anal. Math. \underline{47}(1986)180-192

\bibitem[Fla]{F} D. Flath, ``Decomposition of representations into tensor
products'', in
\underline{Automorphic forms, }
\underline{Representations, } \underline{and
L-functions}, 1979 Corvallis Conference, ed. A. Borel, W. Casselman,
Proc. Symp. Pure Math. vol 33, AMS, 1979

\bibitem[Fl]{Fl} 
Y. Flicker, ``Automorphic forms on covering groups of $GL(2)$'', Invent.
Math. \underline{57} (1980) 119--182

\bibitem[Ge1]{Ge1} S. Gelbart, {\bf Automorphic forms on adele groups}, 
Annals of
Math. Studies, Princeton Univ. Press, Princeton, 1975

\bibitem[Ge2]{Ge2} S. Gelbart, {\bf Weil's representation and the 
spectrum of the metaplectic group},
Springer Lecture Notes in Math \underline{530}, 1976

\bibitem[GeKn]{GeKn} S. Gelbart and A. W. Knapp, ``L-indistinguishability and R 
groups for the
special linear group'', Adv. Math. 43(1982)101-121

\bibitem[GPS]{GPS} S. Gelbart and I. Piatetski-Shapiro, ``Some remarks 
on metaplectic cusp
forms and the correspondences of Shimura and Waldspurger'', Israel 
J. Math.
\underline{44}(1983)97-126

\bibitem[JL]{JL} H. Jacquet and R. Langlands,
{\bf Automorphic forms on $GL(2)$}, Springer Lecture Notes in Math.,
vol 114, 1970

\bibitem[J1]{J1} D. Joyner, 
``A hitch-hiker's guide to genuine invariant distributions
on metaplectic covers of $SL(2)$ over a p-adic field'', preprint

\bibitem[J2]{J2} ------, 
``On the metaplectic analog of Kazhdan's `endoscopic' lifting'', Israel
J. Math. \underline{61}(1988)113-154

\bibitem[J3]{J3} ------, 
``Weighted orbital integrals and weighted characters
on metaplectic covers of $SL(2)$ over a p-adic field'', preprint

\bibitem[KP]{KP} D. Kazhdan, S. Patterson, ``Metaplectic forms'', Publ. Math.
IHES \underbar{59}(1984)35-142

\bibitem[Ko1]{Ko1} R. Kottwitz, ``Stable trace formula: elliptic singular
terms'', Math. Ann. \underline{275} (1986) 365--399

\bibitem[K1]{K1} 
T. Kubota, ``Topological coverings of $SL(2)$ over a local field'', J.
Math. Soc. Japan \underline{19}(1967)114-121

\bibitem[K2]{K2} ------, {\bf Automorphic functions and the reciprocity
law in a number field}, Kyoto U., 1969

\bibitem[LL]{LL} 
J.-P. Labesse, R. P. Langlands, ``L-indistinguishability for $SL(2)$'',
Can. J. Math. \underline{31}(1979)726-785

\bibitem[L1]{L1} R. Langlands, ``Eisenstein series'', Proc Symp Pure Math,
vol. 9, AMS,
Providence, 1966

\bibitem[L2]{L2} ------, {\bf On the functional equations satisfied
by Eisenstein series}, Springer Lecture Notes, \underline{544} (1976)

\bibitem[L3]{L3} ------, {\bf Les d\'ebuts d'une formule des traces 
stable}, Publ. Math. de L'Univ. Paris, VII, 1983

\bibitem[L4]{L4} ------, {\bf Base change for $GL(2)$},
Annals of Math. Studies \underline{96}, Princeton Univ. Press, 1980

\bibitem[MW1]{MW} C. Moeglin, J.-L. Waldspurger,
``Le spectre r\'esiduel de $GL(n)$'', Ann. Sci. \'Ec. Norm. Sup.
\underline{22} (1989) 605--674

\bibitem[MW2]{MW2} ------,
{\bf Spectral decomposition and Eisenstein series},
Cambridge Univ Press, 1995

\bibitem[M]{M} C. Moore, ``Group extensions of $p$-adic and adelic linear
groups'', Publ. Math. IHES \underline{35} (1968) 157--222

\bibitem[Mo1]{Mo1} C. Moen, ``Intertwining operators for covering groups
of $SL(2)$ over
a local field'', preprint

\bibitem[Mo2]{Mo2} ------, ``Irreducibiliy of unitary principal series
representations
for covering groups of $SL(2)$'', Pac. J. Math. 135(1988)89-110

\bibitem[O]{O} T. Ono, ``On Tamagawa numbers'', in Proc. Symp. Pure Math.,
vol. 9, AMS,
Providence, 1966

\bibitem[PR]{PR} G. Prasad, M. S. Raghunathan,
``Topological central extensions of
semi-simple groups over local fields'', Annals of Math.
\underline{119}(1984)143-268

\bibitem[Ro]{Ro} A. Robert, {\bf Introduction to the
representation theory of compact and locally compact groups},
London Math. Series, Cambridge Univ. Press, Cambridge, 1983

\bibitem[Rog]{R} J. Rogawski,
{\bf Automorphic representations of unitary groups
in three variables}
Annals of Math. Studies, \underline{123}, Princeton Univ. Press, 1990

\bibitem[Ru]{Ru} W. Rudin, {\bf Functional analysis}, 
McGraw-Hill, New
York, 1973

\bibitem[Sa]{Sa} 
P. J. Sally Jr, ``Analytic continuation of the irreducible unitary
representations of the universal covering group
of $SL(2,{\mathbb{R}})$'', Memoirs AMS \underbar{69}(1967)

\bibitem[Sh]{Sh} 
F. Shahidi, ``Local coefficients as Artin factors
for real groups'',
Duke Math J. \underbar{52}(1985)973-1007

%\bibitem[Shal]{Shal} J. Shalika, ``Representations of the two by two
%unimodular group over local fields'', IAS

\bibitem[Sha]{Sha} S. Shatz, {\bf Profinite groups, arithmetic, and
geometry},
Princeton Univ. Press, Princeton, 1972

\bibitem[She]{She} D. Shelstad, ``L-indistinguishability for real groups'',
Math
Ann \underbar{259}(1982)385-430

\bibitem[Sho]{Sho} S. Shokranian, {\bf The Selberg-Arthur trace formula},
Springer Lecture Notes in Math., vol. \underline{1503} (1992)

%\bibitem[S]{S} A. Silberger, {\bf Introduction to harmonic analysis
%on reductive p-adic groups}, Princeton Univ. Press, Princeton, 1979

\bibitem[Se]{Se} J.-P. Serre, {\bf Linear representations of finite groups},
Springer-Verlag, New York, 1977

\bibitem[SpSt]{SpSt} T. Springer and R. Steinberg, ``Conjugacy classes'',
in {\bf Seminar on Algebraic groups ...} (by A. Borel, et al),
Springer Lecture Notes \underbar{131}, 2nd printing 1986

\bibitem[Ta]{Ta} M. Taibleson, {\bf Fourier analysis on local fields}, 
Princeton Univ.
Press, Princeton, 1975


\bibitem[We1]{W1} A. Weil, {\bf Basic Number Theory},
Springer-Verlag, 1975

\bibitem[We2]{W2} -----, {\bf Adeles and algebraic groups}, 
Birkhauser, 1982


\end{thebibliography}





\end{document}


